\newcommand{\DoNotLoadEpstopdf}{}
\documentclass[11pt,letterpaper]{article}
\usepackage[nopatch=footnote]{microtype}
\usepackage[margin=0.88in,headheight=15pt,headsep=18pt,footskip=27pt]{geometry}
\usepackage{amsmath,amssymb,amsthm,mathtools}
\usepackage{booktabs,longtable,array,tabularx,graphicx,float,caption}
\usepackage{enumitem,fancyhdr,xcolor,needspace,placeins,xurl,aliascnt}
\usepackage{titlesec,setspace,etoolbox,tocloft,tikz}
\usetikzlibrary{arrows.meta,decorations.pathreplacing,patterns,calc}
\usepackage[hidelinks]{hyperref}
\usepackage{bookmark}
\usepackage[nameinlink,noabbrev]{cleveref}
\allowdisplaybreaks
\setlength{\emergencystretch}{3em}
\setlist[enumerate]{topsep=4pt,itemsep=2pt}
\setlist[itemize]{topsep=4pt,itemsep=2pt}
\setlength{\parskip}{3pt}
\setlength{\parindent}{1.3em}
\setcounter{tocdepth}{2}
\makeatletter
\renewcommand*\l@section{\@dottedtocline{1}{0em}{4.3em}}
\renewcommand*\l@subsection{\@dottedtocline{2}{1.5em}{6.0em}}
\makeatother
\hypersetup{pdftitle={Projector-Resolved Seed Calculus VII: Final Traditional Edition},pdfauthor={Komorebi P.},pdfsubject={Complete corrected finite core, strengthened quantum packet development, and mathematical proof completions}}
\pagestyle{fancy}\fancyhf{}
\fancyhead[L]{\small\textsc{Projector-Resolved Seed Calculus VII}}
\fancyhead[R]{\small Komorebi P.}
\fancyfoot[C]{\thepage}
\newtheorem{theorem}{Theorem}[section]
\newaliascnt{proposition}{theorem}\newtheorem{proposition}[proposition]{Proposition}\aliascntresetthe{proposition}
\newaliascnt{lemma}{theorem}\newtheorem{lemma}[lemma]{Lemma}\aliascntresetthe{lemma}
\newaliascnt{corollary}{theorem}\newtheorem{corollary}[corollary]{Corollary}\aliascntresetthe{corollary}
\theoremstyle{definition}
\newaliascnt{definition}{theorem}\newtheorem{definition}[definition]{Definition}\aliascntresetthe{definition}
\newaliascnt{example}{theorem}\newtheorem{example}[example]{Example}\aliascntresetthe{example}
\newaliascnt{construction}{theorem}\newtheorem{construction}[construction]{Construction}\aliascntresetthe{construction}
\newaliascnt{conjecture}{theorem}\newtheorem{conjecture}[conjecture]{Conjecture}\aliascntresetthe{conjecture}
\newaliascnt{standardfact}{theorem}\newtheorem{standardfact}[standardfact]{Standard Fact}\aliascntresetthe{standardfact}
\newtheorem*{statement}{Supplementary statement}
\theoremstyle{remark}
\newaliascnt{remark}{theorem}\newtheorem{remark}[remark]{Remark}\aliascntresetthe{remark}
\newaliascnt{boundary}{theorem}\newtheorem{boundary}[boundary]{Boundary}\aliascntresetthe{boundary}
\newaliascnt{certificate}{theorem}\newtheorem{certificate}[certificate]{Recorded Numerical Check}\aliascntresetthe{certificate}
\newaliascnt{classical}{theorem}\newtheorem{classical}[classical]{Classical Input}\aliascntresetthe{classical}
\crefname{classical}{classical input}{classical inputs}
\crefname{standardfact}{standard fact}{standard facts}
\newcommand{\cA}{\mathcal A}\newcommand{\cB}{\mathcal B}\newcommand{\cC}{\mathcal C}
\newcommand{\cE}{\mathcal E}\newcommand{\cF}{\mathcal F}\newcommand{\cG}{\mathcal G}
\newcommand{\cH}{\mathcal H}\newcommand{\cK}{\mathcal K}\newcommand{\cL}{\mathcal L}
\newcommand{\cO}{\mathcal O}\newcommand{\cD}{\mathcal D}\newcommand{\cP}{\mathcal P}
\newcommand{\cR}{\mathcal R}\newcommand{\cS}{\mathcal S}\newcommand{\cT}{\mathcal T}
\newcommand{\bbC}{\mathbb C}\newcommand{\bbR}{\mathbb R}\newcommand{\bbZ}{\mathbb Z}\newcommand{\bbN}{\mathbb N}
\newcommand{\CC}{\mathbb C}\newcommand{\RR}{\mathbb R}\newcommand{\ZZ}{\mathbb Z}\newcommand{\NN}{\mathbb N}
\newcommand{\C}{\mathbb C}\newcommand{\R}{\mathbb R}\newcommand{\N}{\mathbb N}
\newcommand{\Hcal}{\mathcal H}\newcommand{\Acal}{\mathcal A}\newcommand{\Bcal}{\mathcal B}
\newcommand{\Ccal}{\mathcal C}\newcommand{\Dcal}{\mathcal D}\newcommand{\Ecal}{\mathcal E}
\newcommand{\Fcal}{\mathcal F}\newcommand{\Gcal}{\mathcal G}\newcommand{\Lcal}{\mathcal L}
\newcommand{\Mcal}{\mathcal M}\newcommand{\Ocal}{\mathcal O}\newcommand{\Pcal}{\mathcal P}
\newcommand{\Qcal}{\mathcal Q}\newcommand{\Rcal}{\mathcal R}\newcommand{\Scal}{\mathcal S}
\newcommand{\Tcal}{\mathcal T}\newcommand{\Ucal}{\mathcal U}\newcommand{\Vcal}{\mathcal V}
\newcommand{\Id}{\operatorname{Id}}\newcommand{\id}{\mathrm{id}}
\newcommand{\Tr}{\operatorname{Tr}}\newcommand{\tr}{\operatorname{tr}}
\newcommand{\rank}{\operatorname{rank}}\newcommand{\Ran}{\operatorname{Ran}}\newcommand{\Ker}{\operatorname{Ker}}
\newcommand{\Spec}{\operatorname{Spec}}\newcommand{\diag}{\operatorname{diag}}
\newcommand{\End}{\operatorname{End}}\newcommand{\Hom}{\operatorname{Hom}}
\newcommand{\dist}{\operatorname{dist}}\newcommand{\ad}{\operatorname{ad}}\newcommand{\Ad}{\operatorname{Ad}}
\newcommand{\Ob}{\operatorname{Ob}}\newcommand{\ob}{\mathfrak o}
\newcommand{\HS}{\mathrm{HS}}\newcommand{\phys}{\mathrm{phys}}
\newcommand{\e}{\mathrm e}\newcommand{\ii}{\mathrm i}\newcommand{\dd}{\,\mathrm d}\newcommand{\one}{\bm 1}
\newcommand{\norm}[1]{\left\lVert#1\right\rVert}\newcommand{\abs}[1]{\left\lvert#1\right\rvert}
\newcommand{\ket}[1]{\left|#1\right\rangle}\newcommand{\bra}[1]{\left\langle#1\right|}
\newcommand{\ip}[2]{\left\langle#1,#2\right\rangle}
\newcommand{\comm}[2]{[#1,#2]}\newcommand{\deph}{\mathcal D_P}
\newcommand{\spanof}{\operatorname{span}}\newcommand{\srank}{\operatorname{srank}}
\newcommand{\code}[1]{{\small\nolinkurl{#1}}}\newcommand{\lean}[1]{{\small\nolinkurl{#1}}}
\newcolumntype{L}[1]{>{\raggedright\arraybackslash}p{#1}}
\newcolumntype{Y}{>{\raggedright\arraybackslash}X}
\newcommand{\startreadingpart}[1]{\clearpage\FloatBarrier\part{#1}\setcounter{section}{0}\renewcommand{\thesection}{\Roman{part}.\arabic{section}}\renewcommand{\theHsection}{\arabic{part}.\arabic{section}}}

% Authored modern monograph profile. The nine mathematical parts are unchanged.
\usepackage{komorebi-monograph}
\renewcommand{\KomorebiTitle}{Quantum Seeds and Effective Dynamics}
\renewcommand{\KomorebiSubtitle}{Symmetry, geometric transport, quantum channels, and bounded refinement towers}
\renewcommand{\KomorebiSeries}{Core series VII}
\renewcommand{\KomorebiKind}{Modern Monograph}
\renewcommand{\KomorebiEdition}{9 October 2026}
\renewcommand{\KomorebiShortTitle}{Paper VII}
\renewcommand{\KomorebiTitleSize}{38}
\renewcommand{\KomorebiTitleLeading}{43}
\renewcommand{\KomorebiCoverTitle}{Quantum Seeds\\and Effective Dynamics}
\renewcommand{\KomorebiCoverBlurb}{A quantum seed generates a space through an operator grammar. Packets resolve how symmetry, evolution, geometry, and hidden degrees of freedom change that space.}
\KomorebiSetup
\setlength{\cftsecnumwidth}{3.7em}
\setlength{\cftsubsecindent}{3.7em}
\setlength{\cftsubsecnumwidth}{4.8em}
\pdfstringdefDisableCommands{\def\quad{ }}
% Grouping adapters render legacy unbraced math alphabets without editing bodies.
\providecommand{\bm}[1]{\symbfit{#1}}
\let\KomorebiPrimitiveMathcal\mathcal
\renewcommand{\mathcal}[1]{{\KomorebiPrimitiveMathcal{#1}}}
\let\KomorebiPrimitiveMathbb\mathbb
\renewcommand{\mathbb}[1]{{\KomorebiPrimitiveMathbb{#1}}}
\let\KomorebiPrimitiveMathfrak\mathfrak
\renewcommand{\mathfrak}[1]{{\KomorebiPrimitiveMathfrak{#1}}}
\let\KomorebiPrimitiveMathsf\mathsf
\renewcommand{\mathsf}[1]{{\KomorebiPrimitiveMathsf{#1}}}
\let\KomorebiPrimitiveMathbf\mathbf
\renewcommand{\mathbf}[1]{{\KomorebiPrimitiveMathbf{#1}}}
\let\KomorebiPrimitiveMathrm\mathrm
\renewcommand{\mathrm}[1]{{\KomorebiPrimitiveMathrm{#1}}}
\let\KomorebiPrimitiveMathtt\mathtt
\renewcommand{\mathtt}[1]{{\KomorebiPrimitiveMathtt{#1}}}
\crefname{proposition}{proposition}{propositions}
\crefname{lemma}{lemma}{lemmas}
\crefname{corollary}{corollary}{corollaries}
\crefname{definition}{definition}{definitions}
\crefname{example}{example}{examples}
\crefname{construction}{construction}{constructions}
\crefname{conjecture}{conjecture}{conjectures}
\crefname{remark}{remark}{remarks}

% Keep each unboxed statement anchor with its visible heading.
\BeforeBeginEnvironment{proposition}{\par\Needspace{5\baselineskip}}
\BeforeBeginEnvironment{lemma}{\par\Needspace{5\baselineskip}}
\BeforeBeginEnvironment{corollary}{\par\Needspace{5\baselineskip}}
\BeforeBeginEnvironment{example}{\par\Needspace{5\baselineskip}}
\BeforeBeginEnvironment{remark}{\par\Needspace{5\baselineskip}}
\BeforeBeginEnvironment{certificate}{\par\Needspace{5\baselineskip}}
\BeforeBeginEnvironment{conjecture}{\par\Needspace{5\baselineskip}}
\BeforeBeginEnvironment{standardfact}{\par\Needspace{5\baselineskip}}
\BeforeBeginEnvironment{classical}{\par\Needspace{5\baselineskip}}
\BeforeBeginEnvironment{construction}{\par\Needspace{5\baselineskip}}
\BeforeBeginEnvironment{statement}{\par\Needspace{5\baselineskip}\phantomsection}
\BeforeBeginEnvironment{theorem}{\par\Needspace{7\baselineskip}}
\BeforeBeginEnvironment{definition}{\par\Needspace{7\baselineskip}}
\begin{document}
\KomorebiCover
\KomorebiOpeningNote{\textbf{Part I.} The corrected finite quantum seed calculus, with its exact released formal contract.

\medskip
\textbf{Part II.} The complete strengthened mathematical development, including all seventy-nine theorem, proposition, lemma, and corollary proofs.

\medskip
\textbf{Part III.} The substantive mathematics of all six formal proof supplements and three retained alternative derivations.

\medskip
\textbf{Part IV.} Forty worked visual pages connect the formulas to finite models and identify the hypotheses needed at each continuum boundary.}

\section*{Reading this edition}\addcontentsline{toc}{section}{Reading this edition}
A seed is useful only together with the operators allowed to act on it. Paper VII follows that principle through symmetry packets, quantum circuits, moving projectors, channels, effective dynamics, and refinement. The finite calculus in Part I supplies the common language. The strengthened development in Part II gives broader bounded constructions and limit criteria. Part III retains the detailed proof completions rather than reducing them to a status summary.

The original finite item numbers keep the prefix I. The strengthened development has the prefix II, and the supplementary constructions have the prefix III. All nine retrieved mathematical source parts are preserved exactly. Overlapping statements and proofs are retained because their hypotheses and uses differ. Their total is not a count of independent discoveries.

The selected archived formal checkpoint covers nineteen scoped manuscript items. It records 323 explicit theorem declarations across twenty-five active Lean proof modules. These are different units. The strengthened Part II is outside that nineteen-item contract. This October monograph does not report a fresh Lean compilation. Exact hypotheses, declarations, inherited compiler records, and reproduction instructions appear in the Lean companion.

Five explicitly named conjectures remain open in the finite core: anomaly-free refinement, uniform Feshbach control, Lorentzian reconstruction, the Einstein infrared limit, and massless spin-two packets. Bounded direct limits and norm limits provide useful criteria. Each physical target still needs the additional structure stated in its conjecture.

\begin{keyidea}[title={Four distinctions to keep in view}]
A seed's algebraic orbit is different from a unitary state-preparation orbit. Compression is different from an effective resolvent. A normalized channel is different from an unnormalized error family. A finite spectrum is different from a continuum particle interpretation. Keeping these distinctions visible makes the later results easier to use.
\end{keyidea}

Read Part I first for the packet grammar and the finite statements. Turn to Part II for bounded operator transport, spectral and channel constructions, and refinement towers. Use Part III when the geometric bundle construction or a formal specialization needs more detail. The visual atlas follows the same topics through explicit examples. It illustrates the statements with finite data and does not replace their proofs.
\clearpage
\tableofcontents
\clearpage
\listoffigures
\clearpage
\pagenumbering{arabic}
\startreadingpart{Corrected finite quantum seed calculus}
\begin{readerroadmap}
The finite core starts with seed rank, then resolves symmetry and circuit transfer into packets. Its geometric statements use the actual range bundle, including the Lie-bracket term. Its channel and Feshbach statements keep normalization and invertibility visible. The three examples below introduce those distinctions with exact finite data.
\end{readerroadmap}
\input{monograph_examples.tex}
\graphicspath{{assets/core/}}
\input{parts/01_finite_core.tex}
\startreadingpart{Strengthened quantum packet calculus and continuum towers}
\begin{readerroadmap}
This part broadens the packet calculus while keeping its hypotheses explicit. Read the nonzero-space, common-space, self-adjointness, finite-address, and trace-class assumptions with each result. Its bounded limits do not establish the five open physical continuum targets.
\end{readerroadmap}
\graphicspath{{assets/strengthened/}}
\section{Introduction}\label{strong:section:001}\label{strong:sec:introduction}

Quantum mechanics repeatedly uses the same operation. A Hilbert space is split into sectors, an operator is resolved between those sectors, and the pieces are propagated or recombined. Symmetry sectors, eigenspaces, syndrome spaces, adiabatic subspaces, and retained model spaces are familiar examples. This article treats them as instances of one finite packet calculus.

The elementary coordinate object is a complete family of orthogonal projectors
\[
P_aP_b=\delta_{ab}P_a,
\qquad
\sum_{a\in I}P_a=I.
\]
Every vector has packet coordinates \(P_a\psi\). Every operator has transfer blocks
\[
A_{ba}=P_bAP_a.
\]
These identities are simple. Their value appears when the same coordinates are used across representation theory, circuits, channels, error correction, effective dynamics, and geometric limits.

Finite dimensional quantum information and operator theory supply the standard background \cite{NielsenChuang2010,Watrous2018,Bhatia1997}. Finite group projectors come from character theory \cite{Serre1977,FultonHarris1991}. Moving spectral subspaces and geometric phase originate in the work of Kato, Berry, Wilczek, and Zee \cite{Kato1950,Kato1995,Berry1984,WilczekZee1984}. Completely positive maps and quantum dynamical semigroups are treated in \cite{Stinespring1955,Choi1975,Kraus1983,Lindblad1976}. Exact error correction is governed by the Knill Laflamme condition \cite{KnillLaflamme1997,SchumacherNielsen1996,BenyKempfKribs2007}. Projection methods and hidden state memory are linked by the Feshbach, Nakajima, and Zwanzig constructions \cite{Feshbach1958,Nakajima1958,Zwanzig1960,BravyiDiVincenzoLoss2011}.

The first part begins with quantum seeds. A state is cyclic only relative to a declared observable algebra. For an algebra in multiplicity form, the local orbit dimension is controlled by the rank of a coefficient matrix. This recovers the relation between Schmidt rank and the dimension reached by local operations. The finite group section then constructs central isotypic projectors. Cyclic symmetry gives a sharper charge selection rule. An operator with charge \(q\) can move a state only from packet \(j\) to packet \(j+q\).

A circuit is resolved by inserting the identity between successive gates. This produces a sum over packet paths. The target probability contains diagonal path terms and interference between distinct paths. Symmetry can remove entire path families before any matrix multiplication is performed. Unitarity gives exact row and column relations among the transfer blocks.

The moving packet section treats time dependent projectors. The canonical generator
\[
K(t)=\sum_a\dot P_a(t)P_a(t)
\]
is skew adjoint and transports every packet. In the moving frame, the physical Hamiltonian acquires the geometric term \(-\ii W(t)^*\dot W(t)\). A leakage estimate shows that the relevant off packet coupling is not \(QHP\) alone. It is
\[
Q(t)\bigl(H(t)-\ii K(t)\bigr)P(t).
\]
This is the exact finite distinction between following a moving subspace and remaining in a fixed one.

Density matrices require two packet labels. The superoperators
\[
\Pcal_{ab}(X)=P_aXP_b
\]
form a complete orthogonal packet system on Hilbert Schmidt operator space. Dephasing is the diagonal projection. Every channel has input and output superblocks, with the same composition law as ordinary matrix blocks. A covariant channel preserves the isotypic packets of the conjugation action. In the cyclic case these are coherence orders. The Choi matrix carries the equivalent symmetry under \(U\otimes\overline U\).

Error correction is expressed through compression to a code. Let \(P\) be the code projector. The class of \(PAP\) modulo scalar multiples of \(P\) is the obstruction to making \(A\) invisible inside the code. The Knill Laflamme criterion says that this obstruction vanishes for every product \(E_a^*E_b\) from the error family. Diagonalizing the resulting positive coefficient matrix produces orthogonal syndrome packets and an explicit recovery on the error span.

The effective dynamics section returns to ordinary Hilbert space packets. Schur complementation gives the exact retained resolvent and a frequency dependent self energy. The same hidden block gives a causal memory kernel in time. For a self adjoint hidden Hamiltonian, the self energy has a fixed sign in the upper half plane, and the time kernel is positive definite. Uniform spectral separation gives direct bounds that can be carried through a refinement tower.

The geometric part has two finite components. A positive transfer matrix gives a reflection positive Hankel kernel and reconstructs a nonnegative Hamiltonian through its logarithm. A finite Laplacian gives a heat trace and a scale dependent spectral dimension. The article proves fixed time convergence from periodic lattices to the continuum torus. Poisson summation then gives the continuum small time dimension.

Constraint closure needs a different obstruction. If \(C_a\) are self adjoint constraints and \(P_{\phys}\) projects onto their common kernel, then
\[
P_{\phys}[C_a,C_b]P_{\phys}=0
\]
holds automatically. This compression cannot test closure. The correct finite obstruction is the class of \([C_a,C_b]\) in the quotient of \(\End(\Hcal)\) by the span of the constraints. Vanishing means that the constraint span is closed under commutators.

The last part gives two bounded continuum tower theorems. Exact compatible finite Hamiltonians define a bounded self adjoint operator on the Hilbert direct limit. A norm summable sequence in a common Hilbert space converges to a bounded self adjoint limit, and its unitary and Euclidean evolutions converge uniformly on bounded time intervals. Reflection positivity and fixed algebraic closure pass to limits under the stated convergence assumptions.

These results do not construct an interacting continuum quantum field theory. They do not prove anomaly freedom on an unbounded common domain, Osterwalder Schrader reconstruction for a continuum measure, an Einstein infrared action, or a positive massless spin two sector. Each of those conclusions needs hypotheses absent from the finite calculations.

\begin{figure}[t]
\centering
\includegraphics[width=\textwidth]{01_quantum_architecture.pdf}
\caption{The finite mathematical layers developed in the article and the separate limit criteria needed at the continuum boundary.}
\label{strong:fig:architecture}
\end{figure}

\begin{table}[t]
\centering
\small
\begin{tabular}{>{\raggedright\arraybackslash}p{.20\textwidth}>{\raggedright\arraybackslash}p{.31\textwidth}>{\raggedright\arraybackslash}p{.39\textwidth}}
\toprule
Object & Packet coordinate & Exact structural condition \\
\midrule
State & \(P_a\psi\) & completeness of the projector family \\
Operator & \(P_bAP_a\) & block multiplication \\
Circuit & intermediate packet path & insertion of complete resolutions \\
Density matrix & \(P_a\rho P_b\) & two sided operator packets \\
Channel & superoperator block & complete positivity and covariance when declared \\
Code & syndrome packet & scalar compression of error products \\
Effective model & Feshbach map & invertibility of the hidden resolvent \\
Euclidean model & positive transfer matrix & positivity and self adjointness \\
Constraint system & quotient obstruction & commutator closure modulo the constraint span \\
Continuum tower & direct or norm limit & compatibility plus explicit convergence \\
\bottomrule
\end{tabular}
\caption{The packet coordinates used in the article.}
\label{strong:tab:dictionary}
\end{table}

\section{Quantum seeds and observable algebras}\label{strong:section:002}\label{strong:sec:seeds}

\begin{definition}[Finite quantum seed system]\label{strong:definition:001}\label{strong:def:quantum_seed}
A finite quantum seed system is a triple
\[
(\Hcal,\Acal,S),
\]
where \(\Hcal\) is a finite dimensional Hilbert space, \(\Acal\subseteq\End(\Hcal)\) is a unital star subalgebra, and \(S\subseteq\Hcal\) is a finite set. Its generated hull is
\[
\Acal S
=
\operatorname{span}\{A\psi:A\in\Acal,\ \psi\in S\}.
\]
The seed set is cyclic when \(\Acal S=\Hcal\). The quantum seed rank is the least cardinality of a cyclic seed set.
\end{definition}

The algebra is part of the data. A vector may be cyclic for \(\End(\Hcal)\) and fail to be cyclic for a restricted observable or gate algebra.

Let
\[
\Acal
\cong
\bigoplus_{r=1}^{s}
\bigl(M_{d_r}(\CC)\otimes I_{m_r}\bigr)
\]
act on
\[
\Hcal
\cong
\bigoplus_{r=1}^{s}
\bigl(\CC^{d_r}\otimes\CC^{m_r}\bigr).
\]
Assume \(d_r,m_r\ge1\) for each displayed block. For an empty block family, the seed rank and the displayed finite supremum are defined as zero. A vector in one block is identified with a \(d_r\) by \(m_r\) coefficient matrix.

\begin{theorem}[Orbit dimension of one seed]\label{strong:theorem:001}\label{strong:thm:orbit_dimension}
For integers \(d\ge1\), \(m\ge0\), and the action of \(M_d(\CC)\otimes I_m\) on \(\CC^d\otimes\CC^m\), let \(C_\psi\) be the coefficient matrix of \(\psi\). Then
\[
\dim\operatorname{span}\{(A\otimes I_m)\psi:A\in M_d(\CC)\}
=
 d\rank C_\psi.
\]
The vector is cyclic exactly when \(\rank C_\psi=m\).
\end{theorem}

\begin{proof}\phantomsection\label{strong:proof:001}
Under the coefficient matrix identification, \((A\otimes I_m)\psi\) corresponds to \(AC_\psi\). Every row of \(AC_\psi\) lies in the row space of \(C_\psi\). Conversely, a matrix whose rows lie in that row space equals \(AC_\psi\) for a suitable \(A\). There are \(d\) independent row choices in a space of dimension \(\rank C_\psi\). The orbit dimension is therefore \(d\rank C_\psi\). It equals \(dm\) exactly when \(C_\psi\) has full column rank.
\phantomsection\label{strong:proof:001:end}\end{proof}

\begin{theorem}[Quantum seed rank]\label{strong:theorem:002}\label{strong:thm:quantum_seed_rank}
For the displayed multiplicity decomposition,
\[
\operatorname{srank}_{\Acal}(\Hcal)
=
\max_{1\le r\le s}
\left\lceil\frac{m_r}{d_r}\right\rceil.
\]
\end{theorem}

\begin{proof}\phantomsection\label{strong:proof:002}
A single seed contributes a row space of dimension at most \(d_r\) to the \(r\)th multiplicity factor. Hence \(k\) seeds cover at most \(kd_r\) multiplicity directions, which gives the lower bound.

Set
\[
k=\max_r\left\lceil\frac{m_r}{d_r}\right\rceil.
\]
For each \(r\), partition a basis of \(\CC^{m_r}\) into at most \(k\) groups of size at most \(d_r\). Choose the coefficient matrix of the \(j\)th seed so that its row space spans the \(j\)th group. The combined row spaces span every multiplicity factor. The preceding theorem then shows that the \(k\) seeds generate every summand.
\phantomsection\label{strong:proof:002:end}\end{proof}

\begin{corollary}[Local orbit and Schmidt rank]\label{strong:corollary:001}\label{strong:cor:schmidt_orbit}
Let \(\psi\in\Hcal_A\otimes\Hcal_B\), let \(d_A=\dim\Hcal_A\ge1\), and let \(r\) be the Schmidt rank of \(\psi\). Then
\[
\dim\operatorname{span}\{(A\otimes I_B)\psi:A\in\End(\Hcal_A)\}
=
 d_A r.
\]
Local operators on \(A\) generate the whole bipartite Hilbert space exactly when the reduced state on \(B\) has full support.
\end{corollary}

\begin{proof}\phantomsection\label{strong:proof:003}
The Schmidt rank is the rank of the coefficient matrix in any product basis. Apply \cref{strong:thm:orbit_dimension}.
\phantomsection\label{strong:proof:003:end}\end{proof}

\begin{figure}[t]
\centering
\includegraphics[width=.68\textwidth]{02_seed_rank.pdf}
\caption{The dimension reached by all local operators is linear in Schmidt rank.}
\label{strong:fig:seed_rank}
\end{figure}

\begin{proposition}[Computational address orbit]\label{strong:proposition:001}\label{strong:prop:computational_orbit}
For \(n\) qubits, let
\[
\psi_0=\ket{0}^{\otimes n}
\]
and let \(X_j\) be the Pauli \(X\) operator on the \(j\)th qubit. Then
\[
\left\{
X_1^{a_1}\cdots X_n^{a_n}\psi_0:
(a_1,\ldots,a_n)\in\{0,1\}^n
\right\}
\]
is the computational basis.
\end{proposition}

\begin{proof}\phantomsection\label{strong:proof:004}
The addressed vector is \(\ket{a_1\cdots a_n}\). Distinct binary addresses produce the \(2^n\) distinct basis vectors.
\phantomsection\label{strong:proof:004:end}\end{proof}

\begin{remark}\label{strong:remark:001}\label{strong:rem:seed_scope}
The seed rank records generation by a declared algebra. It does not measure circuit depth, locality cost, control complexity, or physical reachability under a restricted time dependent Hamiltonian.
\end{remark}

\section{Symmetry packets and selection rules}\label{strong:section:003}\label{strong:sec:symmetry}

Let \(G\) be a finite group and let
\[
U:G\longrightarrow\Ucal(\Hcal)
\]
be a unitary representation. Write \(\widehat G\) for the set of irreducible unitary representations. For \(\lambda\in\widehat G\), let \(d_\lambda\) and \(\chi_\lambda\) denote its dimension and character.

\begin{theorem}[Finite group isotypic projectors]\label{strong:theorem:003}\label{strong:thm:isotypic_projectors}
The operators
\[
P_\lambda
=
\frac{d_\lambda}{\abs{G}}
\sum_{g\in G}
\overline{\chi_\lambda(g)}U(g)
\]
are orthogonal projectors. They satisfy
\[
P_\lambda P_\mu=\delta_{\lambda\mu}P_\lambda,
\qquad
\sum_{\lambda\in\widehat G}P_\lambda=I_{\Hcal}.
\]
The range of \(P_\lambda\) is the \(\lambda\) isotypic component of \(\Hcal\).
\end{theorem}

\begin{proof}\phantomsection\label{strong:proof:005}
Complete reducibility gives a unitary decomposition
\[
\Hcal
\cong
\bigoplus_{\mu\in\widehat G}
V_\mu\otimes M_\mu,
\qquad
U(g)
\cong
\bigoplus_\mu
\pi_\mu(g)\otimes I_{M_\mu}.
\]
On the \(\mu\) summand, the displayed average is central in the irreducible representation \(\pi_\mu\), so it is a scalar by Schur's lemma. Taking its trace and applying character orthogonality shows that the scalar is one when \(\mu=\lambda\) and zero otherwise. The asserted identities follow on the direct sum.
\phantomsection\label{strong:proof:005:end}\end{proof}

The central projectors separate irreducible types but do not choose coordinates inside the representation factor. Matrix units provide that finer resolution.

\begin{theorem}[Representation matrix units]\label{strong:theorem:004}\label{strong:thm:representation_matrix_units}
Fix an irreducible unitary representation \(\pi_\lambda:G\to U(d_\lambda)\) and an orthonormal basis of its representation space. Define
\[
E_{ij}^{\lambda}
=
\frac{d_\lambda}{\abs{G}}
\sum_{g\in G}
\bigl[\pi_\lambda(g^{-1})\bigr]_{ji}U(g).
\]
Then
\[
E_{ij}^{\lambda}E_{k\ell}^{\mu}
=
\delta_{\lambda\mu}\delta_{jk}E_{i\ell}^{\lambda},
\qquad
\bigl(E_{ij}^{\lambda}\bigr)^*=E_{ji}^{\lambda},
\]
and
\[
P_\lambda=\sum_{i=1}^{d_\lambda}E_{ii}^{\lambda}.
\]
\end{theorem}

\begin{proof}\phantomsection\label{strong:proof:006}
Apply the matrix coefficient form of Schur orthogonality to the two group averages. It gives the matrix unit multiplication law. Unitarity gives the adjoint relation after replacing \(g\) by \(g^{-1}\). Summing the diagonal coefficients gives the character average in \cref{strong:thm:isotypic_projectors}.
\phantomsection\label{strong:proof:006:end}\end{proof}

\subsection{Cyclic charge packets}\label{strong:subsec:cyclic_charge}

Let \(N\ge1\) be an integer, \(S\) unitary with \(S^N=I\), and \(\omega=\e^{2\pi\ii/N}\).

\begin{corollary}[Cyclic charge projectors]\label{strong:corollary:002}\label{strong:cor:cyclic_projectors}
The operators
\[
P_j
=
\frac1N
\sum_{r=0}^{N-1}
\omega^{-jr}S^r,
\qquad
j\in\ZZ/N\ZZ,
\]
form a complete orthogonal packet system and satisfy
\[
SP_j=\omega^jP_j.
\]
\end{corollary}

\begin{proof}\phantomsection\label{strong:proof:007}
Apply \cref{strong:thm:isotypic_projectors} to the cyclic group. Its irreducible representations are the one dimensional characters \(r\mapsto\omega^{jr}\).
\phantomsection\label{strong:proof:007:end}\end{proof}

\begin{definition}[Operator charge]\label{strong:definition:002}\label{strong:def:operator_charge}
An operator \(A\in\End(\Hcal)\) has cyclic charge \(q\in\ZZ/N\ZZ\) when
\[
SAS^*=\omega^qA.
\]
\end{definition}

\begin{theorem}[Charge selection rule]\label{strong:theorem:005}\label{strong:thm:charge_selection}
If \(A\) has charge \(q\), then
\[
P_kAP_j=0
\]
unless \(k=j+q\) in \(\ZZ/N\ZZ\). In fact,
\[
AP_j=P_{j+q}AP_j.
\]
\end{theorem}

\begin{proof}\phantomsection\label{strong:proof:008}
For \(x\in P_j\Hcal\),
\[
S(Ax)
=(SAS^*)Sx
=\omega^{q+j}Ax.
\]
Thus \(Ax\in P_{j+q}\Hcal\). The block statement follows.
\phantomsection\label{strong:proof:008:end}\end{proof}

\begin{corollary}[Charge addition]\label{strong:corollary:003}\label{strong:cor:charge_addition}
If \(A\) and \(B\) have charges \(q\) and \(r\), then \(AB\) has charge \(q+r\). The adjoint \(A^*\) has charge \(-q\).
\end{corollary}

\begin{proof}\phantomsection\label{strong:proof:009}
Conjugation by \(S\) is a star automorphism. Apply it to \(AB\) and \(A^*\).
\phantomsection\label{strong:proof:009:end}\end{proof}

\begin{figure}[t]
\centering
\includegraphics[width=.82\textwidth]{03_symmetry_selection.pdf}
\caption{A charge one operator moves each cyclic packet to the next packet and has no other transfer blocks.}
\label{strong:fig:symmetry_selection}
\end{figure}

\subsection{Tensor refinement}\label{strong:subsec:tensor_refinement}

Suppose the \(a\)th tensor factor carries a cyclic packet family
\[
\bigl(P_j^{(a)}\bigr)_{j\in\ZZ/N\ZZ}.
\]
The fine tensor packets are
\[
P_{\bm j}^{\rm fine}
=
P_{j_1}^{(1)}\otimes\cdots\otimes P_{j_m}^{(m)}.
\]

\begin{theorem}[Tensor charge refinement]\label{strong:theorem:006}\label{strong:thm:tensor_refinement}
For \(r\in\ZZ/N\ZZ\), define
\[
P_r^{\rm total}
=
\sum_{j_1+\cdots+j_m=r}
P_{j_1}^{(1)}\otimes\cdots\otimes P_{j_m}^{(m)}.
\]
Then the \(P_r^{\rm total}\) form a complete orthogonal packet system. Every total charge packet is the exact sum of its fine tensor descendants.
\end{theorem}

\begin{proof}\phantomsection\label{strong:proof:010}
The fine tensor projectors are mutually orthogonal and sum to the identity. The total residue classes partition the fine address set. Summing the fine projectors over each class gives the result.
\phantomsection\label{strong:proof:010:end}\end{proof}

\begin{corollary}[Tensor selection rule]\label{strong:corollary:004}\label{strong:cor:tensor_selection}
If \(A_a\) has local charge \(q_a\), then
\[
A_1\otimes\cdots\otimes A_m
\]
has total charge \(q_1+\cdots+q_m\).
\end{corollary}

\begin{proof}\phantomsection\label{strong:proof:011}
Conjugate the tensor product by the total cyclic symmetry and multiply the local phase factors.
\phantomsection\label{strong:proof:011:end}\end{proof}

\begin{remark}\label{strong:remark:002}\label{strong:rem:nonabelian_fusion}
For a nonabelian group, tensor products are refined by irreducible fusion multiplicities rather than residue addition. The central isotypic projectors remain exact, but a complete tensor refinement also requires a choice of Clebsch Gordan intertwiners.
\end{remark}

\section{Quantum circuits as packet paths}\label{strong:section:004}\label{strong:sec:circuits}

Let
\[
U=U_mU_{m-1}\cdots U_1
\]
be a finite circuit with \(m\ge1\), and let \((P_a)_{a\in I}\) be a complete orthogonal packet family. For the empty circuit, \(U=I\) and its transfer block is \(P_bP_a=\delta_{ba}P_a\), whose endpoint sum reconstructs \(I\).

\begin{theorem}[Exact circuit path expansion]\label{strong:theorem:007}\label{strong:thm:circuit_path}
For source packet \(a_0\) and target packet \(a_m\),
\[
P_{a_m}UP_{a_0}
=
\sum_{a_1,\ldots,a_{m-1}}
P_{a_m}U_mP_{a_{m-1}}
\cdots
P_{a_1}U_1P_{a_0}.
\]
Summing over \(a_0\) and \(a_m\) reconstructs \(U\).
\end{theorem}

\begin{proof}\phantomsection\label{strong:proof:012}
Insert \(I=\sum_aP_a\) between every adjacent pair of gates and distribute the finite sums. Completeness at the two endpoints gives the full operator.
\phantomsection\label{strong:proof:012:end}\end{proof}

For a packet path
\[
\gamma=(a_0,a_1,\ldots,a_m),
\]
write
\[
U_\gamma
=
P_{a_m}U_mP_{a_{m-1}}
\cdots
P_{a_1}U_1P_{a_0}.
\]

\begin{theorem}[Packet unitarity identities]\label{strong:theorem:008}\label{strong:thm:packet_unitarity}
If \(U\) is unitary and \(U_{ba}=P_bUP_a\), then
\[
\sum_bU_{ba}^*U_{bc}=\delta_{ac}P_a
\]
and
\[
\sum_aU_{ba}U_{da}^*=\delta_{bd}P_b.
\]
\end{theorem}

\begin{proof}\phantomsection\label{strong:proof:013}
Use completeness and unitarity:
\[
\sum_bU_{ba}^*U_{bc}
=
P_aU^*\left(\sum_bP_b\right)UP_c
=P_aP_c.
\]
The second identity follows in the same way.
\phantomsection\label{strong:proof:013:end}\end{proof}

\begin{corollary}[Probability conservation from one packet]\label{strong:corollary:005}\label{strong:cor:packet_probability}
If \(\psi\in P_a\Hcal\), then
\[
\sum_b\norm{P_bU\psi}^2=\norm{\psi}^2.
\]
\end{corollary}

\begin{proof}\phantomsection\label{strong:proof:014}
Take the inner product of the first identity in \cref{strong:thm:packet_unitarity} with \(\psi\) and set \(c=a\).
\phantomsection\label{strong:proof:014:end}\end{proof}

\begin{theorem}[Path interference formula]\label{strong:theorem:009}\label{strong:thm:path_interference}
Fix a source vector \(\psi\in P_{a_0}\Hcal\) and a target packet \(b\). Let \(\Gamma_b\) be the set of packet paths from \(a_0\) to \(b\), and set
\[
\eta_\gamma=U_\gamma\psi.
\]
Then
\[
\norm{P_bU\psi}^2
=
\sum_{\gamma\in\Gamma_b}\norm{\eta_\gamma}^2
+
2\operatorname{Re}
\sum_{\substack{\gamma,\gamma'\in\Gamma_b\\ \gamma<\gamma'}}
\ip{\eta_\gamma}{\eta_{\gamma'}}.
\]
Consequently,
\[
\norm{P_bU\psi}
\leq
\sum_{\gamma\in\Gamma_b}\norm{\eta_\gamma}.
\]
If the path outputs are pairwise orthogonal, the target probability is the sum of the path probabilities.
\end{theorem}

\begin{proof}\phantomsection\label{strong:proof:015}
The path expansion gives
\[
P_bU\psi=\sum_{\gamma\in\Gamma_b}\eta_\gamma.
\]
Expand the squared norm and group conjugate cross terms. The norm bound is the triangle inequality. Orthogonality removes every cross term.
\phantomsection\label{strong:proof:015:end}\end{proof}

\begin{proposition}[Charge pruning of circuit paths]\label{strong:proposition:002}\label{strong:prop:charge_pruning}
Suppose every gate has a finite charge decomposition
\[
U_r=\sum_qU_r^{(q)},
\qquad
SU_r^{(q)}S^*=\omega^qU_r^{(q)}.
\]
A term with gate charges \(q_1,\ldots,q_m\) can contribute from source charge \(a_0\) only to target charge
\[
a_m=a_0+q_1+\cdots+q_m.
\]
Every packet path violating one of the intermediate charge updates vanishes.
\end{proposition}

\begin{proof}\phantomsection\label{strong:proof:016}
Apply \cref{strong:thm:charge_selection} after each gate component. The packet address is forced at every step.
\phantomsection\label{strong:proof:016:end}\end{proof}

\begin{figure}[t]
\centering
\includegraphics[width=.88\textwidth]{04_circuit_paths.pdf}
\caption{A circuit block is the coherent sum of its allowed packet paths. Selection rules remove forbidden paths before amplitudes are evaluated.}
\label{strong:fig:circuit_paths}
\end{figure}

\begin{example}[Bell circuit]\label{strong:example:001}\label{strong:ex:bell_circuit}
Starting from \(\ket{00}\), apply a Hadamard gate to the first qubit and then a controlled not gate. The final state is
\[
\frac{\ket{00}+\ket{11}}{\sqrt2}.
\]
It lies in the even parity packet. The intermediate state has one even and one odd computational contribution, so the final even amplitude is reconstructed by more than one packet path.
\end{example}

\section{Symmetry resolved dynamics and leakage}\label{strong:section:005}\label{strong:sec:dynamics}

Let \(H=H^*\) and let \((P_a)_{a\in I}\) be a complete orthogonal packet system.

\begin{theorem}[Exact invariant packets]\label{strong:theorem:010}\label{strong:thm:invariant_packets}
The following conditions are equivalent:
\begin{enumerate}
\item \([H,P_a]=0\) for every \(a\)
\item \(P_bHP_a=0\) whenever \(a\neq b\)
\item every packet \(P_a\Hcal\) is invariant under \(\e^{-\ii tH}\) for every real \(t\)
\end{enumerate}
When they hold,
\[
H=\sum_aP_aHP_a
\]
and
\[
\e^{-\ii tH}=\sum_aP_a\e^{-\ii tH}P_a.
\]
\end{theorem}

\begin{proof}\phantomsection\label{strong:proof:017}
Commutation is equivalent to vanishing of every off diagonal block. The same block decomposition passes to polynomials in \(H\), then to the exponential by norm convergence. Invariance of every packet implies invariance under both \(\e^{-\ii tH}\) and its adjoint. Differentiation at \(t=0\) gives commutation with \(H\).
\phantomsection\label{strong:proof:017:end}\end{proof}

Fix one packet projector \(P\) and set \(Q=I-P\).

\begin{lemma}[Off diagonal norm]\label{strong:lemma:001}\label{strong:lem:off_diagonal_norm}
For self adjoint \(H\),
\[
\norm{[H,P]}=\norm{QHP}=\norm{PHQ}.
\]
\end{lemma}

\begin{proof}\phantomsection\label{strong:proof:018}
Relative to \(P\Hcal\oplus Q\Hcal\),
\[
[H,P]
=
\begin{pmatrix}
0&-PHQ\\
QHP&0
\end{pmatrix}.
\]
Its square has diagonal blocks \(-PHQHP\) and \(-QHPHQ\). The norm is therefore the largest singular value of \(QHP\).
\phantomsection\label{strong:proof:018:end}\end{proof}

\begin{theorem}[Fixed packet leakage]\label{strong:theorem:011}\label{strong:thm:fixed_leakage}
For every real \(t\),
\[
\norm{Q\e^{-\ii tH}P}
\leq
\abs{t}\norm{QHP}
=
\abs{t}\norm{[H,P]}.
\]
\end{theorem}

\begin{proof}\phantomsection\label{strong:proof:019}
Let
\[
H_{\rm d}=PHP+QHQ,
\qquad
H_{\rm o}=PHQ+QHP.
\]
The propagator generated by \(H_{\rm d}\) preserves \(P\Hcal\) and \(Q\Hcal\). Duhamel's formula relative to \(H_{\rm d}\) gives
\[
Q\e^{-\ii tH}P
=
-\ii
\int_0^t
Q\e^{-\ii(t-s)H}H_{\rm o}\e^{-\ii sH_{\rm d}}P
\dd s.
\]
Since \(H_{\rm o}P=QHP\) and both propagators are unitary, the norm is at most \(\abs{t}\norm{QHP}\). Use \cref{strong:lem:off_diagonal_norm} for the last equality.
\phantomsection\label{strong:proof:019:end}\end{proof}

\begin{figure}[t]
\centering
\includegraphics[width=.72\textwidth]{05_leakage_bound.pdf}
\caption{Exact transition amplitude in a two packet model and the linear commutator bound.}
\label{strong:fig:leakage}
\end{figure}

\begin{corollary}[Leakage probability]\label{strong:corollary:006}\label{strong:cor:leakage_probability}
If \(\psi\in P\Hcal\) is a unit vector, then
\[
\norm{Q\e^{-\ii tH}\psi}^2
\leq
 t^2\norm{QHP}^2.
\]
\end{corollary}

\begin{proof}\phantomsection\label{strong:proof:020}
Apply \cref{strong:thm:fixed_leakage} and square the result.
\phantomsection\label{strong:proof:020:end}\end{proof}

\section{Moving packets and geometric transport}\label{strong:section:006}\label{strong:sec:moving_packets}

Let \((P_a(t))_{a\in I}\) be a twice continuously differentiable family of complete orthogonal packet systems on a finite dimensional Hilbert space.

\begin{theorem}[Canonical packet connection]\label{strong:theorem:012}\label{strong:thm:kato_connection}
Define
\[
K(t)
=
\sum_a\dot P_a(t)P_a(t).
\]
Then
\[
K(t)^*=-K(t)
\]
and
\[
[K(t),P_a(t)]=\dot P_a(t)
\]
for every \(a\).
\end{theorem}

\begin{proof}\phantomsection\label{strong:proof:021}
Differentiating \(P_a^2=P_a\) gives
\[
\dot P_aP_a+P_a\dot P_a=\dot P_a.
\]
Summing over \(a\) and using \(\sum_a\dot P_a=0\) yields
\[
\sum_aP_a\dot P_a=-\sum_a\dot P_aP_a.
\]
This proves skew adjointness. For fixed \(b\),
\[
KP_b=\dot P_bP_b.
\]
Differentiating \(P_bP_a=0\) for \(a\neq b\) gives
\[
P_b\dot P_a=-\dot P_bP_a.
\]
Hence
\[
P_bK=-\dot P_b(I-P_b).
\]
Subtracting gives \([K,P_b]=\dot P_b\).
\phantomsection\label{strong:proof:021:end}\end{proof}

\begin{theorem}[Exact Kato transport]\label{strong:theorem:013}\label{strong:thm:kato_transport}
Let \(W(t)\) solve
\[
\dot W(t)=K(t)W(t),
\qquad
W(0)=I.
\]
Then \(W(t)\) is unitary and
\[
P_a(t)W(t)=W(t)P_a(0)
\]
for every packet \(a\).
\end{theorem}

\begin{proof}\phantomsection\label{strong:proof:022}
Skew adjointness gives
\[
\frac{\dd}{\dd t}W(t)^*W(t)=0.
\]
Thus \(W(t)\) is unitary. Next,
\[
\frac{\dd}{\dd t}
\left(W(t)^*P_a(t)W(t)\right)
=
W(t)^*
\left(
\dot P_a-[K,P_a]
\right)
W(t)=0.
\]
The expression equals \(P_a(0)\) at \(t=0\), which proves the intertwining identity.
\phantomsection\label{strong:proof:022:end}\end{proof}

\begin{theorem}[Moving frame Hamiltonian]\label{strong:theorem:014}\label{strong:thm:moving_frame}
Let \(U(t)\) solve
\[
\ii\dot U(t)=H(t)U(t),
\qquad
U(0)=I,
\]
where \(H(t)=H(t)^*\) is continuous on the compact time interval under consideration. Set
\[
V(t)=W(t)^*U(t).
\]
Then
\[
\ii\dot V(t)=H_{\rm mov}(t)V(t),
\]
where
\[
H_{\rm mov}(t)
=
W(t)^*H(t)W(t)-\ii W(t)^*\dot W(t)
=
W(t)^*\bigl(H(t)-\ii K(t)\bigr)W(t).
\]
The moving frame Hamiltonian is self adjoint.
\end{theorem}

\begin{proof}\phantomsection\label{strong:proof:023}
Differentiate \(V=W^*U\). Since \(\dot W=KW\) and \(K^*=-K\),
\[
\dot W^*=-W^*K.
\]
Substitution gives the equation. The first term in \(H_{\rm mov}\) is self adjoint. The operator \(-\ii W^*\dot W\) is self adjoint because \(W^*\dot W\) is skew adjoint.
\phantomsection\label{strong:proof:023:end}\end{proof}

\begin{theorem}[Moving packet leakage]\label{strong:theorem:015}\label{strong:thm:moving_leakage}
Fix a packet \(P(t)=P_a(t)\) and put \(Q(t)=I-P(t)\). For \(0\leq t\leq T\),
\[
\norm{Q(t)U(t)P(0)}
\leq
\int_0^t
\norm{Q(s)\bigl(H(s)-\ii K(s)\bigr)P(s)}
\dd s.
\]
In particular, the physical evolution follows the moving packet exactly on \([0,T]\) when the displayed off packet block vanishes for every time in that interval.
\end{theorem}

\begin{proof}\phantomsection\label{strong:proof:024}
In the Kato frame, \(P(0)\) is fixed and
\[
Q(t)U(t)P(0)
=
W(t)Q(0)V(t)P(0).
\]
Apply the time dependent version of the Duhamel argument in \cref{strong:thm:fixed_leakage} to \(H_{\rm mov}(t)\). The off diagonal block is
\[
Q(0)H_{\rm mov}(s)P(0)
=
W(s)^*Q(s)\bigl(H(s)-\ii K(s)\bigr)P(s)W(s).
\]
Unitary conjugation preserves the norm.
\phantomsection\label{strong:proof:024:end}\end{proof}

The geometric term is therefore part of the exact coupling test. A Hamiltonian may be packet diagonal at every fixed time and still fail to transport the moving packet because the packet itself rotates.

\subsection{Discrete packet transport}\label{strong:subsec:discrete_transport}

Let \((P_a)\) and \((Q_a)\) be two complete packet systems with the same index set. Define
\[
S(\Qcal,\Pcal)=\sum_aQ_aP_a.
\]

\begin{theorem}[Direct packet rotation]\label{strong:theorem:016}\label{strong:thm:direct_rotation}
Assume
\[
\norm{S(\Qcal,\Pcal)-I}<1.
\]
Then \(S(\Qcal,\Pcal)\) is invertible. Its unitary polar factor
\[
\Ucal(\Qcal,\Pcal)
=
S(S^*S)^{-1/2}
\]
satisfies
\[
Q_a\Ucal(\Qcal,\Pcal)
=
\Ucal(\Qcal,\Pcal)P_a
\]
for every \(a\).
\end{theorem}

\begin{proof}\phantomsection\label{strong:proof:025}
The norm condition gives invertibility by the Neumann series. The identity
\[
Q_aS=SP_a=Q_aP_a
\]
shows that \(S^*S\) commutes with every \(P_a\). Its inverse square root also commutes with every \(P_a\). Hence
\[
Q_a\Ucal
=Q_aS(S^*S)^{-1/2}
=SP_a(S^*S)^{-1/2}
=\Ucal P_a.
\]
The polar factor is unitary because \(S\) is invertible.
\phantomsection\label{strong:proof:025:end}\end{proof}

\begin{theorem}[Convergence of short packet rotations]\label{strong:theorem:017}\label{strong:thm:short_rotation_convergence}
Let \((P_a(t))\) be twice continuously differentiable on \([0,T]\). For a partition
\[
0=t_0<t_1<\cdots<t_n=T,
\]
let
\[
W_\pi
=
\Ucal\bigl(\Pcal(t_n),\Pcal(t_{n-1})\bigr)
\cdots
\Ucal\bigl(\Pcal(t_1),\Pcal(t_0)\bigr).
\]
For sufficiently fine partitions the factors are defined, and
\[
\norm{W_\pi-W(T)}
\leq
C_T\abs{\pi}
\]
for a constant depending on the packet path but not on the partition.
\end{theorem}

\begin{proof}\phantomsection\label{strong:proof:026}
Taylor expansion gives
\[
S\bigl(\Pcal(t+h),\Pcal(t)\bigr)
=
I+hK(t)+O(h^2)
\]
uniformly on \([0,T]\). Since \(K(t)^*=-K(t)\),
\[
S^*S=I+O(h^2),
\]
so the polar factor satisfies
\[
\Ucal\bigl(\Pcal(t+h),\Pcal(t)\bigr)
=
I+hK(t)+O(h^2).
\]
This is a first order one step approximation to \(\dot W=KW\). A telescoping product estimate over the partition gives a total error bounded by a constant times \(\sum_j(t_{j+1}-t_j)^2\), which is at most \(T\abs{\pi}\).
\phantomsection\label{strong:proof:026:end}\end{proof}

\begin{figure}[t]
\centering
\includegraphics[width=.62\textwidth]{06_kato_transport.pdf}
\caption{A moving rank one packet on the Bloch sphere. The Kato evolution transports the range without choosing a phase convention for its eigenvector.}
\label{strong:fig:kato_transport}
\end{figure}

\subsection{Local frames and curvature}\label{strong:subsec:curvature}

Let \(P(x)\) be a smooth rank \(r\) orthogonal projector over a parameter manifold. Choose a local isometry
\[
V(x):\CC^r\longrightarrow\Hcal,
\qquad
V(x)V(x)^*=P(x).
\]

\begin{theorem}[Nonabelian packet connection]\label{strong:theorem:018}\label{strong:thm:nonabelian_connection}
The matrix valued one form
\[
\Acal=V^*\dd V
\]
is skew adjoint. Under a frame change \(V'=VG\), where \(G(x)\in U(r)\),
\[
\Acal'
=
G^*\Acal G+G^*\dd G.
\]
Its curvature
\[
\Fcal=\dd\Acal+\Acal\wedge\Acal
\]
satisfies
\[
\Fcal=V^*(\dd P\wedge\dd P)V.
\]
\end{theorem}

\begin{proof}\phantomsection\label{strong:proof:027}
Differentiate \(V^*V=I\) to obtain \(\Acal^*=-\Acal\). The gauge transformation law follows by differentiating \(VG\). Since \(P=VV^*\) and \((I-P)\dd V=\dd P\,V\), direct expansion gives
\[
\dd\Acal+\Acal\wedge\Acal
=
\dd V^*(I-P)\wedge\dd V
=
V^*(\dd P\wedge\dd P)V.
\]
\phantomsection\label{strong:proof:027:end}\end{proof}

\begin{corollary}[Bloch sphere curvature]\label{strong:corollary:007}\label{strong:cor:bloch_curvature}
Let
\[
P(\theta,\phi)
=
\frac12\left(I+\bm n(\theta,\phi)\cdot\bm\sigma\right),
\]
where
\[
\bm n=(\sin\theta\cos\phi,\sin\theta\sin\phi,\cos\theta).
\]
With the convention
\[
\Fcal_{\theta\phi}
=-\ii\Tr\bigl(P[\partial_\theta P,\partial_\phi P]\bigr),
\]
one has
\[
\Fcal_{\theta\phi}=\frac12\sin\theta
\]
and
\[
\frac1{2\pi}
\int_0^{2\pi}\int_0^\pi
\Fcal_{\theta\phi}\dd\theta\dd\phi
=1.
\]
\end{corollary}

\begin{proof}\phantomsection\label{strong:proof:028}
The Pauli identity
\[
(\bm a\cdot\bm\sigma)(\bm b\cdot\bm\sigma)
=(\bm a\cdot\bm b)I
+\ii(\bm a\times\bm b)\cdot\bm\sigma
\]
gives
\[
-\ii\Tr\bigl(P[\partial_\theta P,\partial_\phi P]\bigr)
=\frac12
\bm n\cdot
(\partial_\theta\bm n\times\partial_\phi\bm n)
=\frac12\sin\theta.
\]
Integration gives one.
\phantomsection\label{strong:proof:028:end}\end{proof}

\section{Operator packets, measurements, and quantum channels}\label{strong:section:007}\label{strong:sec:channels}

Density matrices and observables carry both an input packet and an output packet. Equip \(\End(\Hcal)\) with the Hilbert Schmidt inner product
\[
\ip{X}{Y}_{\HS}=\Tr(X^*Y).
\]

\begin{definition}[Two sided operator packet]\label{strong:definition:003}\label{strong:def:operator_packet}
For a complete orthogonal family \((P_a)\), define
\[
\Pcal_{ab}(X)=P_aXP_b.
\]
The range
\[
\Bcal_{ab}=P_a\End(\Hcal)P_b
\]
is the operator packet with output address \(a\) and input address \(b\).
\end{definition}

\begin{theorem}[Complete operator packet system]\label{strong:theorem:019}\label{strong:thm:operator_packets}
The superoperators \((\Pcal_{ab})_{a,b\in I}\) are mutually orthogonal projections on Hilbert Schmidt operator space. They satisfy
\[
X=\sum_{a,b}P_aXP_b
\]
and
\[
\norm{X}_{\HS}^2
=
\sum_{a,b}\norm{P_aXP_b}_{\HS}^2.
\]
\end{theorem}

\begin{proof}\phantomsection\label{strong:proof:029}
Idempotence and completeness follow from the projector identities. For Hilbert Schmidt orthogonality,
\[
\begin{aligned}
\ip{P_aXP_b}{P_cYP_d}_{\HS}
&=
\Tr(P_bX^*P_aP_cYP_d)\\
&=
\delta_{ac}\delta_{bd}
\Tr(P_bX^*P_aYP_b).
\end{aligned}
\]
The Pythagorean identity follows.
\phantomsection\label{strong:proof:029:end}\end{proof}

\begin{figure}[t]
\centering
\includegraphics[width=.58\textwidth]{07_operator_packets.pdf}
\caption{A density matrix or observable carries two packet addresses. Diagonal packets contain populations and internal observables. Off diagonal packets contain coherences.}
\label{strong:fig:operator_packets}
\end{figure}

\begin{definition}[Packet dephasing]\label{strong:definition:004}\label{strong:def:dephasing}
The packet dephasing map is
\[
\Dcal_{\Pcal}(X)=\sum_aP_aXP_a.
\]
\end{definition}

\begin{theorem}[Dephasing as conditional expectation]\label{strong:theorem:020}\label{strong:thm:dephasing_expectation}
The map \(\Dcal_{\Pcal}\) is completely positive, trace preserving, unital, and idempotent. It is the Hilbert Schmidt orthogonal projection onto the block diagonal algebra
\[
\Acal_{\Pcal}
=
\{X:[X,P_a]=0\text{ for every }a\}.
\]
In particular,
\[
\norm{\Dcal_{\Pcal}(X)}_{\HS}
\leq
\norm{X}_{\HS}.
\]
\end{theorem}

\begin{proof}\phantomsection\label{strong:proof:030}
The projectors \(P_a\) form a Kraus family because \(\sum_aP_a^*P_a=I\). Thus the map is completely positive and trace preserving. It is unital for the same reason. Orthogonality and idempotence follow from \cref{strong:thm:operator_packets}. Its range consists exactly of the diagonal operator packets, which is the commutant of the projector family.
\phantomsection\label{strong:proof:030:end}\end{proof}

\begin{corollary}[Block diagonal measurement invariance]\label{strong:corollary:008}\label{strong:cor:measurement_invariance}
Let \((E_k)\) be a positive operator valued measurement, and suppose
\[
E_k=\Dcal_{\Pcal}(E_k)
\]
for every \(k\). Then every density matrix \(\rho\) satisfies
\[
\Tr(E_k\rho)
=
\Tr\bigl(E_k\Dcal_{\Pcal}(\rho)\bigr).
\]
\end{corollary}

\begin{proof}\phantomsection\label{strong:proof:031}
Since \(\Dcal_{\Pcal}\) is self adjoint for the Hilbert Schmidt inner product,
\[
\Tr(E_k\rho)
=
\Tr\bigl(\Dcal_{\Pcal}(E_k)\rho\bigr)
=
\Tr\bigl(E_k\Dcal_{\Pcal}(\rho)\bigr).
\]
\phantomsection\label{strong:proof:031:end}\end{proof}

\subsection{Channel superblocks}\label{strong:subsec:channel_blocks}

Let \(\Phi:\End(\Hcal_X)\to\End(\Hcal_Y)\) be linear. Choose packet systems \((P_a)\) on \(\Hcal_X\) and \((Q_c)\) on \(\Hcal_Y\). Define
\[
\Phi_{cd,ab}
=
\Qcal_{cd}\Phi\Pcal_{ab}.
\]

\begin{theorem}[Exact channel block calculus]\label{strong:theorem:021}\label{strong:thm:channel_blocks}
Every linear map has the exact reconstruction
\[
\Phi=\sum_{a,b,c,d}\Phi_{cd,ab}.
\]
If \(\Psi:\End(\Hcal_Y)\to\End(\Hcal_Z)\), then
\[
(\Psi\Phi)_{ef,ab}
=
\sum_{c,d}
\Psi_{ef,cd}\Phi_{cd,ab}.
\]
\end{theorem}

\begin{proof}\phantomsection\label{strong:proof:032}
Insert the complete operator packet resolutions before and after \(\Phi\). For composition, insert the middle resolution between \(\Psi\) and \(\Phi\) and collect the output and input packets.
\phantomsection\label{strong:proof:032:end}\end{proof}

\begin{standardfact}[Finite Kraus and Choi criteria]\label{strong:standardfact:001}\label{strong:fact:kraus_choi}
A linear map \(\Phi:M_n(\CC)\to M_m(\CC)\) is completely positive exactly when it has a Kraus representation
\[
\Phi(X)=\sum_rK_rXK_r^*.
\]
It is trace preserving exactly when
\[
\sum_rK_r^*K_r=I.
\]
Equivalently, its Choi matrix
\[
J(\Phi)
=
\sum_{i,j}\Phi(\ket{i}\bra{j})\otimes\ket{i}\bra{j}
\]
is positive, and trace preservation is
\[
\Tr_{\rm out}J(\Phi)=I_{\rm in}.
\]
\end{standardfact}

These finite criteria are classical \cite{Stinespring1955,Choi1975,Kraus1983,Watrous2018}.

\subsection{Covariance and coherence order}\label{strong:subsec:channel_covariance}

Let \(S\) be unitary and \(N\ge1\) an integer with \(S^N=I\), and define the conjugation action
\[
\Ad_S(X)=SXS^*.
\]

\begin{theorem}[Coherence order projectors]\label{strong:theorem:022}\label{strong:thm:coherence_projectors}
For \(q\in\ZZ/N\ZZ\), define
\[
\Lcal_q(X)
=
\frac1N
\sum_{r=0}^{N-1}
\omega^{-qr}S^rXS^{-r}.
\]
Then the \(\Lcal_q\) form a complete orthogonal packet system on Hilbert Schmidt operator space. Their ranges are the operators satisfying
\[
SXS^*=\omega^qX.
\]
If \(P_j\) are the state charge projectors, then
\[
\Ran\Lcal_q
=
\bigoplus_jP_{j+q}\End(\Hcal)P_j.
\]
\end{theorem}

\begin{proof}\phantomsection\label{strong:proof:033}
Apply the cyclic character projector formula to the unitary representation \(\Ad_S\) on operator space. The last identity follows from \cref{strong:thm:charge_selection} applied to each operator packet.
\phantomsection\label{strong:proof:033:end}\end{proof}

\begin{theorem}[Covariant channels preserve coherence order]\label{strong:theorem:023}\label{strong:thm:covariant_channel}
Suppose
\[
\Phi(SXS^*)=S\Phi(X)S^*
\]
for every \(X\). Then
\[
\Phi\Lcal_q=\Lcal_q\Phi
\]
for every \(q\). Thus \(\Phi\) cannot mix distinct coherence orders.
\end{theorem}

\begin{proof}\phantomsection\label{strong:proof:034}
The covariance identity says that \(\Phi\) commutes with every power of \(\Ad_S\). It therefore commutes with each finite character average \(\Lcal_q\).
\phantomsection\label{strong:proof:034:end}\end{proof}

\begin{theorem}[Choi symmetry criterion]\label{strong:theorem:024}\label{strong:thm:choi_covariance}
Let \(U:G\to\Ucal(\Hcal)\) be a finite group representation. A channel \(\Phi\) is covariant,
\[
\Phi\bigl(U(g)XU(g)^*\bigr)
=
U(g)\Phi(X)U(g)^*,
\]
if and only if
\[
\bigl(U(g)\otimes\overline{U(g)}\bigr)
J(\Phi)
\bigl(U(g)^*\otimes U(g)^{\mathsf T}\bigr)
=
J(\Phi)
\]
for every \(g\in G\).
\end{theorem}

\begin{proof}\phantomsection\label{strong:proof:035}
Let
\[
\Omega=\sum_i\ket{i}\otimes\ket{i}.
\]
Then
\[
J(\Phi)=(\Phi\otimes\id)(\ket{\Omega}\bra{\Omega}).
\]
The vector \(\Omega\) is invariant under \(U\otimes\overline U\). Covariance therefore implies the displayed Choi invariance. Conversely, the displayed transformation is the Choi matrix of
\[
\Ad_{U(g)}\circ\Phi\circ\Ad_{U(g)^*}.
\]
The Choi correspondence is injective, so invariance implies that this map equals \(\Phi\), which is covariance.
\phantomsection\label{strong:proof:035:end}\end{proof}

\begin{corollary}[Choi packet decomposition]\label{strong:corollary:009}\label{strong:cor:choi_packets}
The Choi matrix of a covariant channel is block diagonal with respect to the isotypic decomposition of \(U\otimes\overline U\).
\end{corollary}

\begin{proof}\phantomsection\label{strong:proof:036}
Apply \cref{strong:thm:isotypic_projectors} to the representation \(U\otimes\overline U\) and use \cref{strong:thm:choi_covariance}.
\phantomsection\label{strong:proof:036:end}\end{proof}

\begin{example}[Amplitude damping]\label{strong:example:002}\label{strong:ex:amplitude_damping}
For \(0\leq\gamma\leq1\), let
\[
K_0=
\begin{pmatrix}
1&0\\
0&\sqrt{1-\gamma}
\end{pmatrix},
\qquad
K_1=
\begin{pmatrix}
0&\sqrt\gamma\\
0&0
\end{pmatrix}.
\]
The channel
\[
\Phi_\gamma(\rho)=K_0\rho K_0^*+K_1\rho K_1^*
\]
is trace preserving. It sends the off diagonal packet \(\ket0\bra1\) to \(\sqrt{1-\gamma}\ket0\bra1\), while the excited population transfers partly into the ground population.
\end{example}

\begin{figure}[t]
\centering
\includegraphics[width=.82\textwidth]{08_channel_coherence.pdf}
\caption{Amplitude damping transfers population and contracts the coherence packet.}
\label{strong:fig:channel_coherence}
\end{figure}


\subsection{Quantum dynamical semigroups}\label{strong:subsec:lindblad}

\begin{standardfact}[Finite Lindblad generator]\label{strong:standardfact:002}\label{strong:fact:lindblad}
A norm continuous semigroup of completely positive trace preserving maps on a finite matrix algebra, in the Schr\"odinger picture, has a generator of the form
\[
\Gcal(X)
=
-\ii[H,X]
+
\sum_r
\left(
L_rXL_r^*
-
\frac12L_r^*L_rX
-
\frac12XL_r^*L_r
\right),
\]
where \(H=H^*\). Conversely, every operator of this form generates a completely positive trace preserving semigroup \cite{Lindblad1976}.
\end{standardfact}

\begin{theorem}[Charge preservation by a Lindblad generator]\label{strong:theorem:025}\label{strong:thm:lindblad_charge}
Let \(S\) be unitary and \(N\ge1\) an integer with \(S^N=I\). Assume
\[
SHS^*=H
\]
and, for every jump operator, there is a charge \(q_r\) such that
\[
SL_rS^*=\omega^{q_r}L_r.
\]
Then the Lindblad generator \(\Gcal\) preserves every coherence order:
\[
\Gcal\Lcal_p=\Lcal_p\Gcal.
\]
The quantum dynamical semigroup \(\e^{t\Gcal}\) has the same property.
\end{theorem}

\begin{proof}\phantomsection\label{strong:proof:037}
Let \(X\) have coherence order \(p\). The Hamiltonian commutator has order \(p\) because \(H\) has charge zero. The jump term \(L_rXL_r^*\) has charge
\[
q_r+p-q_r=p.
\]
Also \(L_r^*L_r\) has charge zero, so both anticommutator terms have charge \(p\). Thus \(\Gcal\) preserves \(\Ran\Lcal_p\). The exponential series then shows that \(\e^{t\Gcal}\) preserves the same packets.
\phantomsection\label{strong:proof:037:end}\end{proof}

\begin{corollary}[Independent coherence evolution]\label{strong:corollary:010}\label{strong:cor:no_coherence_mixing}
Under the hypotheses of \cref{strong:thm:lindblad_charge}, each coherence component of an initial operator evolves within its own order. In particular, a density matrix in the zero order sector remains in that sector.
\end{corollary}

\begin{proof}\phantomsection\label{strong:proof:038}
Decompose the initial operator by \cref{strong:thm:coherence_projectors} and apply the invariant packet statement to each component.
\phantomsection\label{strong:proof:038:end}\end{proof}

The theorem does not say that every diagonal quantum master equation closes to a classical Markov chain on packet populations. Scalar population closure needs additional isotropy inside each packet.

\section{Quantum error correction as a compression obstruction}\label{strong:section:008}\label{strong:sec:qec}

Let \(P\) be the orthogonal projector onto a code space
\[
\Ccal=P\Hcal,
\qquad
k=\dim\Ccal\ge1.
\]
The compression of \(A\in\End(\Hcal)\) to the code is \(PAP\in\End(\Ccal)\). The zero code has only the zero supported operator and no normalized code state; that branch is treated in the mathematical completion in Part III.

\begin{definition}[Code compression obstruction]\label{strong:definition:005}\label{strong:def:code_obstruction}
Let
\[
q_P:\End(\Ccal)\longrightarrow\End(\Ccal)/(\CC I_{\Ccal})
\]
be the quotient map. Define
\[
\Ob_P(A)=q_P(PAP).
\]
Equivalently, one may use the traceless representative
\[
\ob_P(A)
=
PAP-
\frac{\Tr(PAP)}{k}P.
\]
\end{definition}

The class vanishes exactly when \(A\) acts as a scalar after compression to the code. Such an operator cannot distinguish normalized code vectors by expectation value.

Exact correction of an error span means that one completely positive trace preserving recovery $\Rcal$ satisfies $\Rcal(F\rho F^*)=q_F\rho$ for every $F$ in that span and every code-supported matrix $\rho$, with a nonnegative scalar $q_F$ depending on $F$ but not on the logical input. For every normalized channel with Kraus operators in the same span, the same recovery restores the code input with total scalar one. Individual unnormalized errors are not required to have unit recovery weight.

\begin{theorem}[Knill Laflamme as obstruction vanishing]\label{strong:theorem:026}\label{strong:thm:kl_obstruction}
Let \(E_1,\ldots,E_r\) be an error family. The code \(\Ccal\) corrects their linear span exactly if and only if
\[
\Ob_P(E_a^*E_b)=0
\]
for all \(a,b\). Equivalently, there is a matrix \(c=(c_{ab})\) such that
\[
PE_a^*E_bP=c_{ab}P
\]
for all \(a,b\).
\end{theorem}

\begin{proof}\phantomsection\label{strong:proof:039}
The equivalence between the quotient statement and the scalar compression identity is immediate from the definition. The latter identity is the finite Knill Laflamme criterion \cite{KnillLaflamme1997,SchumacherNielsen1996}.
\phantomsection\label{strong:proof:039:end}\end{proof}

The next result gives the packet construction behind the criterion.

\begin{theorem}[Orthogonal syndrome construction]\label{strong:theorem:027}\label{strong:thm:syndrome_construction}
Assume the Knill Laflamme identities. Then the coefficient matrix \(c\) is positive semidefinite. Choose a unitary matrix \(W\) such that
\[
WcW^*=\diag(\lambda_1,\ldots,\lambda_r),
\qquad
\lambda_j\geq0,
\]
and define
\[
F_j=\sum_a\overline{W_{ja}}E_a.
\]
For every \(j\) with \(\lambda_j>0\), the operator
\[
V_j=\lambda_j^{-1/2}F_jP
\]
is an isometry from the code into \(\Hcal\). Its range projector
\[
Q_j=V_jV_j^*
\]
is a syndrome packet, and the nonzero \(Q_j\) are pairwise orthogonal.
\end{theorem}

\begin{proof}\phantomsection\label{strong:proof:040}
For \(z\in\CC^r\), let \(F_z=\sum_az_aE_a\). Then
\[
PF_z^*F_zP=(z^*cz)P\geq0,
\]
so \(c\geq0\). The unitary change of error coordinates gives
\[
PF_j^*F_\ell P
=
\lambda_j\delta_{j\ell}P.
\]
Thus
\[
V_j^*V_\ell=\delta_{j\ell}P.
\]
Each \(V_j\) is an isometry on \(P\Hcal\), and distinct ranges are orthogonal. Their range projectors are therefore pairwise orthogonal.
\phantomsection\label{strong:proof:040:end}\end{proof}

\begin{theorem}[Exact recovery on the error span]\label{strong:theorem:028}\label{strong:thm:exact_recovery}
Under the hypotheses of \cref{strong:thm:syndrome_construction}, define on the syndrome support
\[
Q=\sum_{\lambda_j>0}Q_j
\]
the completely positive map
\[
\Rcal_0(X)
=
\sum_{\lambda_j>0}
V_j^*Q_jXQ_jV_j.
\]
It extends to a quantum channel \(\Rcal\) on \(\End(\Hcal)\). If
\[
\Phi(X)=\sum_aE_aXE_a^*
\]
is trace preserving on code states, then
\[
\Rcal\Phi(\rho)=\rho
\]
for every density matrix \(\rho=P\rho P\).
\end{theorem}

\begin{proof}\phantomsection\label{strong:proof:041}
The map \(\Rcal_0\) is completely positive. It is trace preserving on operators supported in \(Q\Hcal\) because
\[
\Tr\bigl(V_j^*Q_jXQ_jV_j\bigr)
=
\Tr(Q_jXQ_j).
\]
Fix a code state \(\rho_0\) and write \(Q^\perp=I-Q\). Then
\[
\Rcal(X)
=
\Rcal_0(QXQ)
+
\Tr(Q^\perp X)\rho_0
\]
is completely positive and trace preserving. It extends \(\Rcal_0\) on operators supported in \(Q\Hcal\).

Unitary mixing of a Kraus family does not change the channel, so
\[
\Phi(\rho)
=
\sum_jF_j\rho F_j^*
=
\sum_{\lambda_j>0}
\lambda_jV_j\rho V_j^*.
\]
Recovery gives
\[
\Rcal\Phi(\rho)
=
\left(\sum_j\lambda_j\right)\rho.
\]
Trace preservation on a code state implies \(\sum_j\lambda_j=1\).
\phantomsection\label{strong:proof:041:end}\end{proof}

\begin{corollary}[Syndrome measurement preserves logical information]\label{strong:corollary:011}\label{strong:cor:syndrome_preserves_logic}
For a channel trace preserving on code states as in the preceding theorem, the completed projective measurement consisting of \((Q_j)_{\lambda_j>0}\) and \(I-\sum_{\lambda_j>0}Q_j\) identifies the orthogonal error syndrome without resolving the encoded state inside the code. The complement outcome has probability zero on the corrupted code state.
\end{corollary}

\begin{proof}\phantomsection\label{strong:proof:042}
On syndrome \(j\), the encoded state is carried by the isometry \(V_j\). The outcome probability is \(\lambda_j\), independent of the normalized code state, while the postmeasurement state is \(V_j\rho V_j^*\).
\phantomsection\label{strong:proof:042:end}\end{proof}

\begin{definition}[Approximate Knill Laflamme defect]\label{strong:definition:006}\label{strong:def:approximate_kl}
For a finite error family define, with the maximum of the empty family set to zero,
\[
\varepsilon_{\rm KL}
=
\max_{a,b}
\inf_{c\in\CC}
\norm{PE_a^*E_bP-cP}.
\]
\end{definition}

\begin{proposition}[Meaning of the approximate defect]\label{strong:proposition:003}\label{strong:prop:approximate_kl_meaning}
For each pair \((a,b)\), choose a minimizing scalar \(c_{ab}\). For unit code vectors \(\psi,\phi\),
\[
\abs{
\ip{E_a\psi}{E_b\phi}
-c_{ab}\ip{\psi}{\phi}
}
\leq
\varepsilon_{\rm KL}.
\]
Thus \(\varepsilon_{\rm KL}\) controls both state dependence of error overlaps and nonorthogonality of nominal syndrome images.
\end{proposition}

\begin{proof}\phantomsection\label{strong:proof:043}
The left side is
\[
\abs{
\ip{\psi}{(PE_a^*E_bP-c_{ab}P)\phi}
},
\]
which is bounded by the operator norm.
\phantomsection\label{strong:proof:043:end}\end{proof}

No recovery fidelity estimate is inferred from this norm alone. Approximate error correction requires a separate continuity theorem and a declared channel metric.

\begin{example}[Three qubit repetition code]\label{strong:example:003}\label{strong:ex:bit_flip_code}
Let
\[
\ket{0_L}=\ket{000},
\qquad
\ket{1_L}=\ket{111},
\]
and take the errors
\[
I,
\quad
X\otimes I\otimes I,
\quad
I\otimes X\otimes I,
\quad
I\otimes I\otimes X.
\]
Their images of the code are four pairwise orthogonal two dimensional syndrome packets whose sum is the full eight dimensional Hilbert space.
\end{example}

\begin{figure}[t]
\centering
\includegraphics[width=.68\textwidth]{09_qec_syndromes.pdf}
\caption{The syndrome packets of the three qubit repetition code. Each packet has dimension two and carries the full logical state.}
\label{strong:fig:qec_syndromes}
\end{figure}

\begin{standardfact}[Correctable observable algebras]\label{strong:standardfact:003}\label{strong:fact:operator_qec}
Let \(\Phi\) be a finite completely positive trace preserving channel with Kraus operators \((E_a)\). A finite dimensional unital star subalgebra \(\Mcal\subseteq P\End(\Hcal)P\) is correctable when a completely positive trace preserving recovery \(\Rcal\) satisfies \(P(\Phi^*\circ\Rcal^*)(X)P=X\) for every \(X\in\Mcal\), where stars on maps denote Heisenberg adjoints. This holds exactly when every compressed product \(PE_a^*E_bP\) belongs to the commutant \(\Mcal'\) on the code. The scalar code criterion is the special case \(\Mcal=P\End(\Hcal)P\) \cite{BenyKempfKribs2007}.
\end{standardfact}

\section{Feshbach effective dynamics and hidden packet memory}\label{strong:section:009}\label{strong:sec:feshbach}

Let \(P\) be an orthogonal projector, put \(Q=I-P\), and write a self adjoint Hamiltonian as
\[
H=
\begin{pmatrix}
A&B\\
B^*&D
\end{pmatrix}
\]
relative to \(P\Hcal\oplus Q\Hcal\).

\begin{definition}[Feshbach map and self energy]\label{strong:definition:007}\label{strong:def:feshbach_map}
For \(z\notin\Spec D\), define
\[
\Sigma(z)=B(z-D)^{-1}B^*
\]
and
\[
F_P(z)=z-A-\Sigma(z).
\]
\end{definition}

\begin{theorem}[Exact projected resolvent]\label{strong:theorem:029}\label{strong:thm:feshbach_resolvent}
If \(z\notin\Spec D\) and \(F_P(z)\) is invertible, then
\[
P(z-H)^{-1}P=F_P(z)^{-1}
\]
as an operator on \(P\Hcal\).
\end{theorem}

\begin{proof}\phantomsection\label{strong:proof:044}
The lower right block of \(z-H\) is \(z-D\). Block Gaussian elimination gives
\[
\begin{pmatrix}
I&B(z-D)^{-1}\\
0&I
\end{pmatrix}
(z-H)
\begin{pmatrix}
I&0\\
(z-D)^{-1}B^*&I
\end{pmatrix}
=
\begin{pmatrix}
F_P(z)&0\\
0&z-D
\end{pmatrix}.
\]
Invert the factors and read the upper left block.
\phantomsection\label{strong:proof:044:end}\end{proof}

\begin{theorem}[Spectral equivalence]\label{strong:theorem:030}\label{strong:thm:feshbach_spectrum}
Let \(z\notin\Spec D\). Then
\[
z\in\Spec H
\quad\Longleftrightarrow\quad
0\in\Spec F_P(z).
\]
The map
\[
x\longmapsto
\begin{pmatrix}
x\\
(z-D)^{-1}B^*x
\end{pmatrix}
\]
is an isomorphism from \(\Ker F_P(z)\) onto \(\Ker(z-H)\).
\end{theorem}

\begin{proof}\phantomsection\label{strong:proof:045}
The equation \((z-H)(x,y)^{\mathsf T}=0\) is
\[
(z-A)x-By=0,
\qquad
-B^*x+(z-D)y=0.
\]
The second equation gives \(y=(z-D)^{-1}B^*x\). Substitution into the first gives \(F_P(z)x=0\). The construction is reversible.
\phantomsection\label{strong:proof:045:end}\end{proof}

\begin{figure}[t]
\centering
\includegraphics[width=.78\textwidth]{10_feshbach_spectrum.pdf}
\caption{The zeros of the Feshbach determinant recover the eigenvalues of the full finite Hamiltonian away from the hidden spectrum.}
\label{strong:fig:feshbach_spectrum}
\end{figure}

\begin{theorem}[Upper half plane sign]\label{strong:theorem:031}\label{strong:thm:self_energy_sign}
If \(\operatorname{Im}z>0\), then
\[
\operatorname{Im}\Sigma(z)
=
-(\operatorname{Im}z)
B(z-D)^{-1}(\overline z-D)^{-1}B^*
\leq0.
\]
Thus \(-\Sigma\) is an operator valued Herglotz function.
\end{theorem}

\begin{proof}\phantomsection\label{strong:proof:046}
For \(R(z)=(z-D)^{-1}\), the resolvent identity gives
\[
R(z)-R(z)^*
=(\overline z-z)R(z)R(z)^*.
\]
Divide by \(2\ii\) and sandwich the result between \(B\) and \(B^*\).
\phantomsection\label{strong:proof:046:end}\end{proof}

\begin{proposition}[Uniform self energy bounds]\label{strong:proposition:004}\label{strong:prop:self_energy_bounds}
Let
\[
\delta=\dist(z,\Spec D)>0,
\qquad
g=\norm B.
\]
Then
\[
\norm{\Sigma(z)}\leq\frac{g^2}{\delta}
\]
and
\[
\norm{\Sigma'(z)}\leq\frac{g^2}{\delta^2}.
\]
For real \(E,E'\) in an interval at distance at least \(\delta\) from \(\Spec D\),
\[
\norm{\Sigma(E)-\Sigma(E')}
\leq
\frac{g^2}{\delta^2}\abs{E-E'}.
\]
\end{proposition}

\begin{proof}\phantomsection\label{strong:proof:047}
The self adjoint resolvent satisfies
\[
\norm{(z-D)^{-1}}=\frac1{\dist(z,\Spec D)}.
\]
Differentiate the resolvent to obtain
\[
\Sigma'(z)=-B(z-D)^{-2}B^*.
\]
The first two estimates follow. The last is the mean value estimate along the real interval.
\phantomsection\label{strong:proof:047:end}\end{proof}

\begin{proposition}[Low energy expansion]\label{strong:proposition:005}\label{strong:prop:low_energy_expansion}
Assume \(D\) is invertible and \(\abs z\norm{D^{-1}}<1\). Then
\[
\Sigma(z)
=-
\sum_{n=0}^{\infty}
z^nBD^{-n-1}B^*
\]
with norm convergence. In particular,
\[
F_P(z)
=
 z-A+BD^{-1}B^*+O\bigl(\abs z\norm B^2\norm{D^{-1}}^2\bigr).
\]
\end{proposition}

\begin{proof}\phantomsection\label{strong:proof:048}
Write
\[
(z-D)^{-1}
=-D^{-1}(I-zD^{-1})^{-1}
=-\sum_{n=0}^{\infty}z^nD^{-n-1}
\]
and sandwich by \(B\) and \(B^*\).
\phantomsection\label{strong:proof:048:end}\end{proof}

\begin{proposition}[Bare truncation error identity]\label{strong:proposition:006}\label{strong:prop:bare_truncation}
If both \(F_P(z)\) and \(z-A\) are invertible, then
\[
P(z-H)^{-1}P-(z-A)^{-1}
=
F_P(z)^{-1}\Sigma(z)(z-A)^{-1}.
\]
\end{proposition}

\begin{proof}\phantomsection\label{strong:proof:049}
Use the inverse difference identity
\[
X^{-1}-Y^{-1}=X^{-1}(Y-X)Y^{-1}
\]
with \(X=F_P(z)\) and \(Y=z-A\).
\phantomsection\label{strong:proof:049:end}\end{proof}

\subsection{Time memory}\label{strong:subsec:time_memory}

Write a Schrödinger state as \(\psi=x+y\), where \(x=P\psi\) and \(y=Q\psi\).

\begin{theorem}[Exact retained memory equation]\label{strong:theorem:032}\label{strong:thm:memory_equation}
The retained coordinate satisfies
\[
\begin{aligned}
\ii\dot x(t)
={}&Ax(t)+B\e^{-\ii tD}y(0)\\
&-\ii\int_0^tK(t-s)x(s)\dd s,
\end{aligned}
\]
where
\[
K(t)=B\e^{-\ii tD}B^*.
\]
For \(\operatorname{Re}s>0\),
\[
\int_0^\infty\e^{-st}K(t)\dd t
=
B(s+\ii D)^{-1}B^*.
\]
\end{theorem}

\begin{proof}\phantomsection\label{strong:proof:050}
The hidden equation is
\[
\ii\dot y=B^*x+Dy.
\]
Variation of constants gives
\[
y(t)=\e^{-\ii tD}y(0)
-\ii\int_0^t\e^{-\ii(t-s)D}B^*x(s)\dd s.
\]
Substitute into \(\ii\dot x=Ax+By\). The Laplace formula follows by diagonalizing \(D\) or integrating the norm convergent finite dimensional exponential.
\phantomsection\label{strong:proof:050:end}\end{proof}

\begin{theorem}[Positive definite quantum memory]\label{strong:theorem:033}\label{strong:thm:positive_memory}
The kernel \(K(t-s)\) is positive definite. For times \(t_1,\ldots,t_m\) and retained vectors \(v_1,\ldots,v_m\),
\[
\sum_{j,k=1}^{m}
\ip{v_j}{K(t_j-t_k)v_k}
\geq0.
\]
\end{theorem}

\begin{proof}\phantomsection\label{strong:proof:051}
Set
\[
w=\sum_{k=1}^{m}\e^{\ii t_kD}B^*v_k.
\]
Then
\[
\norm w^2
=
\sum_{j,k}
\ip{B^*v_j}{\e^{-\ii(t_j-t_k)D}B^*v_k}
=
\sum_{j,k}
\ip{v_j}{K(t_j-t_k)v_k}.
\]
\phantomsection\label{strong:proof:051:end}\end{proof}

\begin{corollary}[Finite oscillatory memory expansion]\label{strong:corollary:012}\label{strong:cor:oscillatory_memory}
If
\[
D=\sum_{r}E_rQ_r
\]
is the spectral decomposition of the hidden Hamiltonian, then
\[
K(t)=\sum_r\e^{-\ii tE_r}BQ_rB^*.
\]
Every coefficient \(BQ_rB^*\) is positive semidefinite.
\end{corollary}

\begin{proof}\phantomsection\label{strong:proof:052}
Insert the spectral decomposition into the exponential.
\phantomsection\label{strong:proof:052:end}\end{proof}

\begin{figure}[t]
\centering
\includegraphics[width=.76\textwidth]{11_memory_kernel.pdf}
\caption{A finite hidden spectrum produces an oscillatory memory kernel rather than an instantaneous retained Hamiltonian.}
\label{strong:fig:memory_kernel}
\end{figure}

\begin{theorem}[Exact local closure criterion]\label{strong:theorem:034}\label{strong:thm:local_closure}
Assume \(y(0)=0\). The retained equation is local and equal to
\[
\ii\dot x=Ax
\]
for every retained initial state exactly when \(B=0\). Equivalently, exact local closure holds exactly when the retained and hidden packets reduce \(H\).
\end{theorem}

\begin{proof}\phantomsection\label{strong:proof:053}
If \(B=0\), the retained and hidden equations decouple. Conversely, suppose that every solution with \(y(0)=0\) satisfies \(x(t)=\e^{-\ii tA}x(0)\). The full block equation gives
\[
\ddot x(0)
=
-\bigl(A^2+BB^*\bigr)x(0),
\]
while the declared local equation gives
\[
\ddot x(0)=-A^2x(0).
\]
Since this holds for every retained initial vector, \(BB^*=0\). Hence \(B=0\).
\phantomsection\label{strong:proof:053:end}\end{proof}

\subsection{Stability across models}\label{strong:subsec:feshbach_stability}

\begin{theorem}[Stability of the self energy]\label{strong:theorem:035}\label{strong:thm:self_energy_stability}
Let \(D=D^*\) and \(\widetilde D=\widetilde D^*\) act on one fixed hidden Hilbert space, and let \(B,\widetilde B\) map that hidden space to one fixed retained Hilbert space. Define
\[
\Sigma(z)=B(z-D)^{-1}B^*,
\qquad
\widetilde\Sigma(z)=\widetilde B(z-\widetilde D)^{-1}\widetilde B^*.
\]
Assume
\[
\norm B,\norm{\widetilde B}\leq g
\]
and
\[
\dist(z,\Spec D),
\dist(z,\Spec\widetilde D)
\geq\delta>0.
\]
Then
\[
\norm{\Sigma(z)-\widetilde\Sigma(z)}
\leq
\frac{2g}{\delta}\norm{B-\widetilde B}
+
\frac{g^2}{\delta^2}\norm{D-\widetilde D}.
\]
\end{theorem}

\begin{proof}\phantomsection\label{strong:proof:054}
Add and subtract \(\widetilde B(z-D)^{-1}B^*\) and \(\widetilde B(z-D)^{-1}\widetilde B^*\). The first two differences give the coupling term. For the last, use
\[
(z-D)^{-1}-(z-\widetilde D)^{-1}
=(z-D)^{-1}(D-\widetilde D)(z-\widetilde D)^{-1}.
\]
Apply the resolvent bounds.
\phantomsection\label{strong:proof:054:end}\end{proof}

\begin{corollary}[Uniform Feshbach control]\label{strong:corollary:013}\label{strong:cor:uniform_feshbach}
For a family represented on fixed retained and hidden Hilbert spaces, with self adjoint hidden Hamiltonians, a common hidden spectral separation \(\delta\), and a common coupling bound \(g\), the self energies are uniformly bounded and uniformly Lipschitz on the declared spectral region. If the couplings and hidden Hamiltonians converge in norm, then the self energies converge uniformly there.
\end{corollary}

\begin{proof}\phantomsection\label{strong:proof:055}
Apply \cref{strong:prop:self_energy_bounds,strong:thm:self_energy_stability} uniformly.
\phantomsection\label{strong:proof:055:end}\end{proof}

This result identifies the finite hypotheses needed for cutoff independent Feshbach control. It does not prove that a chosen geometric refinement tower satisfies them.

\section{Euclidean transfer and finite reflection positivity}\label{strong:section:010}\label{strong:sec:euclidean}

Let \(H=H^*\) on a nonzero finite dimensional Hilbert space and let
\[
E_0=\min\Spec H.
\]
For a lattice time step \(a>0\), define
\[
T=\e^{-a(H-E_0I)}.
\]

\begin{theorem}[Finite Euclidean transfer]\label{strong:theorem:036}\label{strong:thm:euclidean_transfer}
The operator \(T\) is positive definite, self adjoint, and contractive. Its spectrum lies in \((0,1]\), and
\[
T^n=\e^{-na(H-E_0I)}
\]
for every \(n\in\NN\).
\end{theorem}

\begin{proof}\phantomsection\label{strong:proof:056}
Diagonalize \(H\). Each eigenvalue of \(T\) is
\[
\e^{-a(E_j-E_0)}
\]
and therefore lies in \((0,1]\). The power law follows from the functional calculus.
\phantomsection\label{strong:proof:056:end}\end{proof}

\begin{theorem}[Hamiltonian reconstruction from transfer]\label{strong:theorem:037}\label{strong:thm:transfer_logarithm}
Let \(T\) be positive definite, self adjoint, and contractive. Then
\[
H_T=-a^{-1}\log T
\]
is a nonnegative self adjoint operator and
\[
T^n=\e^{-naH_T}.
\]
It is the unique nonnegative self adjoint operator with \(T=\e^{-aH_T}\).
\end{theorem}

\begin{proof}\phantomsection\label{strong:proof:057}
Apply the real logarithm to the positive spectrum of \(T\). The spectral calculus gives existence and the power identity. Injectivity of the scalar exponential on \([0,\infty)\) gives uniqueness.
\phantomsection\label{strong:proof:057:end}\end{proof}

\begin{theorem}[Finite reflection positivity]\label{strong:theorem:038}\label{strong:thm:finite_reflection_positivity}
For vectors \(v_1,\ldots,v_m\in\Hcal\) and nonnegative integers \(n_1,\ldots,n_m\), the matrix
\[
G_{jk}=\ip{v_j}{T^{n_j+n_k}v_k}
\]
is positive semidefinite. Equivalently,
\[
\sum_{j,k}\overline{c_j}c_k
\ip{v_j}{T^{n_j+n_k}v_k}
\geq0
\]
for every \(c\in\CC^m\).
\end{theorem}

\begin{proof}\phantomsection\label{strong:proof:058}
Since \(T\) is positive and self adjoint,
\[
\ip{v_j}{T^{n_j+n_k}v_k}
=
\ip{T^{n_j}v_j}{T^{n_k}v_k}.
\]
The quadratic form is the squared norm of
\[
\sum_kc_kT^{n_k}v_k.
\]
\phantomsection\label{strong:proof:058:end}\end{proof}

\begin{corollary}[Reflection positive correlators]\label{strong:corollary:014}\label{strong:cor:reflection_correlators}
Fix a vector \(\Omega\) and operators \(A_1,\ldots,A_m\). The matrix
\[
C_{jk}
=
\ip{A_j\Omega}{T^{n_j+n_k}A_k\Omega}
\]
is positive semidefinite.
\end{corollary}

\begin{proof}\phantomsection\label{strong:proof:059}
Apply \cref{strong:thm:finite_reflection_positivity} with \(v_j=A_j\Omega\).
\phantomsection\label{strong:proof:059:end}\end{proof}

\begin{figure}[t]
\centering
\includegraphics[width=.68\textwidth]{12_reflection_positivity.pdf}
\caption{Eigenvalues of a finite reflection positive Gram matrix. Small numerical negative values would signal either roundoff or failure of the transfer assumptions.}
\label{strong:fig:reflection_positivity}
\end{figure}

\begin{proposition}[Ground packet convergence]\label{strong:proposition:007}\label{strong:prop:ground_packet}
Let \(P_0\) be the spectral projector of \(H\) at \(E_0\), and let
\[
\Delta=\min\{E_j-E_0:E_j>E_0\}.
\]
If an excited eigenvalue exists, then
\[
\norm{T^n-P_0}=\e^{-na\Delta}.
\]
\end{proposition}

\begin{proof}\phantomsection\label{strong:proof:060}
In the spectral decomposition, \(T^n-P_0\) has eigenvalues zero on the ground space and \(\e^{-na(E_j-E_0)}\) on the excited eigenspaces. The largest is \(\e^{-na\Delta}\).
\phantomsection\label{strong:proof:060:end}\end{proof}

\begin{proposition}[Symmetry resolved transfer]\label{strong:proposition:008}\label{strong:prop:symmetry_transfer}
If \([H,P_a]=0\) for a complete packet family, then
\[
T=\sum_aP_aTP_a
\]
and each reflection positive Gram matrix decomposes as the sum of its packet Gram matrices.
\end{proposition}

\begin{proof}\phantomsection\label{strong:proof:061}
The functional calculus preserves commutation. Decompose every vector into packet components and use orthogonality.
\phantomsection\label{strong:proof:061:end}\end{proof}

Finite reflection positivity is a necessary structural property of transfer data. It is not by itself the continuum Osterwalder Schrader theorem, which also requires a suitable algebra of Euclidean fields, covariance, regularity, clustering, and a controlled continuum limit \cite{OsterwalderSchrader1973,OsterwalderSchrader1975,GlimmJaffe1987}.

\section{Finite spectral geometry and diffusion dimension}\label{strong:section:011}\label{strong:sec:spectral_geometry}

A finite spectral datum consists of a finite dimensional Hilbert space, a nonnegative self adjoint operator \(L\), and any declared packet family. One may take \(L=D^2\) for a finite Dirac type operator or a graph Laplacian. Related geometric viewpoints appear in \cite{Connes1994,Vassilevich2003}.

\begin{definition}[Heat trace and spectral dimension]\label{strong:definition:008}\label{strong:def:spectral_dimension}
Let the eigenvalues of \(L\) be \(\lambda_1,\ldots,\lambda_N\), with at least one zero eigenvalue. Define
\[
Z_L(t)=\Tr(\e^{-tL})
\]
and
\[
d_L(t)
=-2\frac{\dd\log Z_L(t)}{\dd\log t}
=
2t
\frac{\Tr(L\e^{-tL})}{\Tr(\e^{-tL})}.
\]
\end{definition}

\begin{proposition}[Finite scale limits]\label{strong:proposition:009}\label{strong:prop:finite_spectral_limits}
For every finite nonnegative \(L\) with nontrivial kernel,
\[
\lim_{t\downarrow0}d_L(t)=0
\]
and
\[
\lim_{t\to\infty}d_L(t)=0.
\]
For every \(t>0\),
\[
0\leq d_L(t)\leq2t\norm L.
\]
\end{proposition}

\begin{proof}\phantomsection\label{strong:proof:062}
The formula in \cref{strong:def:spectral_dimension} is twice \(t\) times a weighted average of the eigenvalues. This gives the bound and the limit at zero. At large time, the denominator converges to \(\dim\Ker L\), while every positive eigenvalue contribution in the numerator decays exponentially.
\phantomsection\label{strong:proof:062:end}\end{proof}

Thus a finite system cannot have a nonzero spectral dimension at the extreme ultraviolet or infrared ends. A geometric dimension can appear only in a refinement window.

\begin{proposition}[Heat packet decomposition]\label{strong:proposition:010}\label{strong:prop:heat_packets}
If \([L,P_a]=0\) for every packet, then
\[
Z_L(t)=\sum_aZ_{L,a}(t),
\qquad
Z_{L,a}(t)=\Tr(P_a\e^{-tL}).
\]
The global spectral dimension is the heat weight average
\[
d_L(t)
=
\sum_{a:Z_{L,a}(t)>0}
\frac{Z_{L,a}(t)}{Z_L(t)}
 d_{L,a}(t),
\]
where
\[
d_{L,a}(t)
=2t
\frac{\Tr(P_aL\e^{-tL})}{Z_{L,a}(t)}
\]
for nonempty packet heat traces.
\end{proposition}

\begin{proof}\phantomsection\label{strong:proof:063}
Commutation gives a block diagonal heat operator. Add the traces and the corresponding first spectral moments, then divide by the total heat trace.
\phantomsection\label{strong:proof:063:end}\end{proof}

\subsection{Periodic lattices}\label{strong:subsec:periodic_lattices}

Let \(N\geq2\) and let \(\Lambda_N=(\ZZ/N\ZZ)^d\) be the unit flat torus grid with spacing \(1/N\). The standard periodic difference Laplacian has eigenvalues
\[
\lambda_N(k)
=
4N^2
\sum_{j=1}^{d}
\sin^2\left(\frac{\pi k_j}{N}\right),
\]
where \(k\) runs through centered residue representatives. This is the standard periodic finite difference Laplacian \cite{Strikwerda2004}.

\begin{theorem}[Low mode dispersion in dimension \(d\)]\label{strong:theorem:039}\label{strong:thm:dispersion_bound}
For every centered mode \(k\),
\[
0
\leq
4\pi^2\abs{k}^2-\lambda_N(k)
\leq
\frac{4\pi^4}{3N^2}
\sum_{j=1}^{d}k_j^4.
\]
\end{theorem}

\begin{proof}\phantomsection\label{strong:proof:064}
For \(x\in[0,\pi/2]\),
\[
\sin x\leq x
\]
and
\[
\sin x\geq x-\frac{x^3}{6}.
\]
Thus
\[
0\leq x^2-\sin^2x\leq\frac{x^4}{3}.
\]
Apply this with \(x=\pi\abs{k_j}/N\), multiply by \(4N^2\), and sum over \(j\).
\phantomsection\label{strong:proof:064:end}\end{proof}

\begin{theorem}[Fixed time heat trace convergence]\label{strong:theorem:040}\label{strong:thm:heat_trace_convergence}
For every \(t>0\),
\[
\sum_{k\in\Lambda_N^*}
\e^{-t\lambda_N(k)}
\longrightarrow
\sum_{k\in\ZZ^d}
\e^{-4\pi^2t\abs{k}^2}
\]
as \(N\to\infty\). The first spectral moments also converge:
\[
\sum_{k\in\Lambda_N^*}
\lambda_N(k)\e^{-t\lambda_N(k)}
\longrightarrow
\sum_{k\in\ZZ^d}
4\pi^2\abs{k}^2\e^{-4\pi^2t\abs{k}^2}.
\]
Consequently, the discrete spectral dimensions converge to the continuum torus spectral dimension at every fixed positive time.
\end{theorem}

\begin{proof}\phantomsection\label{strong:proof:065}
For a fixed integer mode, \cref{strong:thm:dispersion_bound} gives
\[
\lambda_N(k)\longrightarrow4\pi^2\abs{k}^2.
\]
For centered modes, the inequality
\[
\sin x\geq\frac{2x}{\pi},
\qquad
0\leq x\leq\frac\pi2,
\]
gives
\[
\lambda_N(k)\geq16\abs{k}^2.
\]
Also \(\lambda_N(k)\leq4\pi^2\abs{k}^2\). Extend the discrete summands by zero outside the centered residue box. The heat trace terms are dominated by
\[
\e^{-16t\abs{k}^2},
\]
and the moment terms by
\[
4\pi^2\abs{k}^2\e^{-16t\abs{k}^2}.
\]
Both dominators are summable on \(\ZZ^d\). Dominated convergence proves the two limits. Their quotient gives spectral dimension convergence.
\phantomsection\label{strong:proof:065:end}\end{proof}

\begin{standardfact}[Theta transformation]\label{strong:standardfact:004}\label{strong:fact:theta_transform}
The continuum flat torus heat trace satisfies
\[
Z_d(t)
=
\sum_{k\in\ZZ^d}\e^{-4\pi^2t\abs{k}^2}
=
(4\pi t)^{-d/2}
\sum_{n\in\ZZ^d}\e^{-\abs{n}^2/(4t)}.
\]
This is the Poisson summation formula for a Gaussian.
\end{standardfact}

\begin{theorem}[Continuum torus small time dimension]\label{strong:theorem:041}\label{strong:thm:continuum_dimension}
Let
\[
d_{\rm torus}(t)
=-2\frac{\dd\log Z_d(t)}{\dd\log t}.
\]
Then
\[
d_{\rm torus}(t)
=
 d+O\left(t^{-1}\e^{-1/(4t)}\right)
\]
as \(t\downarrow0\).
\end{theorem}

\begin{proof}\phantomsection\label{strong:proof:066}
By \cref{strong:fact:theta_transform},
\[
Z_d(t)
=(4\pi t)^{-d/2}(1+R(t)),
\]
where
\[
R(t)=\sum_{n\neq0}\e^{-\abs{n}^2/(4t)}
=O\bigl(\e^{-1/(4t)}\bigr).
\]
Differentiation gives
\[
tR'(t)=O\left(t^{-1}\e^{-1/(4t)}\right).
\]
Insert this in the logarithmic derivative.
\phantomsection\label{strong:proof:066:end}\end{proof}

\begin{corollary}[Rigorous intermediate dimension window]\label{strong:corollary:015}\label{strong:cor:dimension_window}
For every \(\varepsilon>0\), there is a diffusion time \(t_\varepsilon>0\) and an integer \(N_\varepsilon\) such that
\[
\abs{d_{L_N}(t_\varepsilon)-d}<\varepsilon
\]
for every \(N\geq N_\varepsilon\).
\end{corollary}

\begin{proof}\phantomsection\label{strong:proof:067}
Choose \(t_\varepsilon\) so that \cref{strong:thm:continuum_dimension} is within \(\varepsilon/2\) of \(d\). Then apply \cref{strong:thm:heat_trace_convergence} to choose \(N_\varepsilon\) so that the discrete spectral dimension is within \(\varepsilon/2\) of the continuum value.
\phantomsection\label{strong:proof:067:end}\end{proof}

\begin{figure}[t]
\centering
\includegraphics[width=.78\textwidth]{13_spectral_dimension.pdf}
\caption{Finite periodic lattices show an intermediate dimension window. The ultraviolet and infrared limits return to zero because the state space is finite.}
\label{strong:fig:spectral_dimension}
\end{figure}

The order of limits matters. Taking \(t\downarrow0\) at fixed \(N\) gives zero. Taking \(N\to\infty\) at fixed positive \(t\), then taking \(t\downarrow0\), gives the continuum dimension.

\section{Constraint closure and anomaly packets}\label{strong:section:012}\label{strong:sec:constraints}

Let \(C_1,\ldots,C_m\) be self adjoint constraint operators on a finite dimensional Hilbert space. The distinction between constraint action and constraint closure is standard in constrained Hamiltonian theory \cite{Dirac1964,HenneauxTeitelboim1992}. Write
\[
\mathfrak c=\operatorname{span}_{\CC}\{C_1,\ldots,C_m\}
\subseteq\End(\Hcal).
\]

\begin{definition}[Universal closure obstruction]\label{strong:definition:009}\label{strong:def:closure_obstruction}
Let
\[
q_{\mathfrak c}:\End(\Hcal)\longrightarrow\End(\Hcal)/\mathfrak c
\]
be the quotient map. Define
\[
\Ob_{ab}^{\rm cl}
=
q_{\mathfrak c}([C_a,C_b]).
\]
\end{definition}

\begin{theorem}[Finite constraint closure criterion]\label{strong:theorem:042}\label{strong:thm:constraint_closure}
The constraint span \(\mathfrak c\) is closed under commutators exactly when
\[
\Ob_{ab}^{\rm cl}=0
\]
for all \(a,b\).
\end{theorem}

\begin{proof}\phantomsection\label{strong:proof:068}
The quotient class vanishes exactly when \([C_a,C_b]\in\mathfrak c\). This holds for every pair exactly when \(\mathfrak c\) is a Lie algebra under the commutator.
\phantomsection\label{strong:proof:068:end}\end{proof}

If structure coefficients \(f_{ab}^{\phantom{ab}c}\) are prescribed, define the fixed structure anomaly
\[
\mathfrak A_{ab}
=
[C_a,C_b]
-\ii\sum_cf_{ab}^{\phantom{ab}c}C_c.
\]
Exact realization of the prescribed algebra is the stronger condition \(\mathfrak A_{ab}=0\).

\begin{proposition}[Physical compression is not a closure test]\label{strong:proposition:011}\label{strong:prop:physical_compression_trivial}
Let
\[
\Hcal_{\phys}=\bigcap_a\Ker C_a
\]
and let \(P_{\phys}\) be its orthogonal projector. Then
\[
P_{\phys}[C_a,C_b]P_{\phys}=0
\]
for every pair \(a,b\), whether or not the constraint span closes.
\end{proposition}

\begin{proof}\phantomsection\label{strong:proof:069}
Self adjointness and \(C_aP_{\phys}=0\) give \(P_{\phys}C_a=0\). Both terms in the compressed commutator therefore vanish.
\phantomsection\label{strong:proof:069:end}\end{proof}

This is why the quotient in \cref{strong:def:closure_obstruction} is needed. Compression to physical states tests action on the common kernel, not closure of the operator algebra.

\begin{figure}[t]
\centering
\includegraphics[width=.88\textwidth]{15_constraint_obstruction.pdf}
\caption{The closure obstruction lives in a quotient by the constraint span. Physical compression is automatic for self adjoint constraints and cannot replace this quotient.}
\label{strong:fig:constraint_obstruction}
\end{figure}

\begin{theorem}[Invariance of the physical space]\label{strong:theorem:043}\label{strong:thm:physical_invariance}
Let \(H=H^*\). Suppose
\[
[H,C_a]
=
\sum_cb_a^{\phantom a c}C_c
\]
for every \(a\). Then \(H\) preserves \(\Hcal_{\phys}\), and so does \(\e^{-\ii tH}\) for every real \(t\).
\end{theorem}

\begin{proof}\phantomsection\label{strong:proof:070}
For \(\psi\in\Hcal_{\phys}\),
\[
C_aH\psi
=HC_a\psi-[H,C_a]\psi
=0.
\]
Thus \(H\psi\in\Hcal_{\phys}\). The exponential preserves every invariant finite dimensional subspace.
\phantomsection\label{strong:proof:070:end}\end{proof}

\subsection{Constraint refinement}\label{strong:subsec:constraint_refinement}

Let \((\Hcal_\ell,C_{\ell,a})\) and \((\Hcal_{\ell+1},C_{\ell+1,a})\) be two levels, and let
\[
R_\ell:\Hcal_\ell\longrightarrow\Hcal_{\ell+1}
\]
be an isometry.

\begin{theorem}[Exact constraint transport]\label{strong:theorem:044}\label{strong:thm:constraint_transport}
Assume
\[
C_{\ell+1,a}R_\ell=R_\ell C_{\ell,a}
\]
for every \(a\). Then
\[
[C_{\ell+1,a},C_{\ell+1,b}]R_\ell
=
R_\ell[C_{\ell,a},C_{\ell,b}].
\]
If the same structure coefficients are used at both levels, then
\[
\mathfrak A_{\ell+1,ab}R_\ell
=
R_\ell\mathfrak A_{\ell,ab}.
\]
The isometry maps coarse physical states into fine physical states.
\end{theorem}

\begin{proof}\phantomsection\label{strong:proof:071}
Apply the intertwining identity twice to each product in the commutator. The anomaly identity follows by subtracting the same linear combination of constraints. If \(C_{\ell,a}\psi=0\), then
\[
C_{\ell+1,a}R_\ell\psi=R_\ell C_{\ell,a}\psi=0.
\]
\phantomsection\label{strong:proof:071:end}\end{proof}

\begin{proposition}[Approximate constraint transport]\label{strong:proposition:012}\label{strong:prop:approx_constraint_transport}
Assume
\[
\norm{C_{\ell+1,a}R_\ell-R_\ell C_{\ell,a}}
\leq\varepsilon_a
\]
and
\[
\norm{C_{\ell,a}},\norm{C_{\ell+1,a}}\leq M.
\]
Then a coarse physical unit vector \(\psi\) satisfies
\[
\norm{C_{\ell+1,a}R_\ell\psi}\leq\varepsilon_a.
\]
Also,
\[
\begin{aligned}
&\norm{
[C_{\ell+1,a},C_{\ell+1,b}]R_\ell
-R_\ell[C_{\ell,a},C_{\ell,b}]
}\\
&\hspace{4em}\leq
2M(\varepsilon_a+\varepsilon_b).
\end{aligned}
\]
\end{proposition}

\begin{proof}\phantomsection\label{strong:proof:072}
The first claim follows by applying the intertwining defect to \(\psi\). Write
\[
E_a=C_{\ell+1,a}R_\ell-R_\ell C_{\ell,a}.
\]
Expanding the two commutators gives four error terms:
\[
E_aC_{\ell,b},
\quad
C_{\ell+1,a}E_b,
\quad
-E_bC_{\ell,a},
\quad
-C_{\ell+1,b}E_a.
\]
Their norms give the bound.
\phantomsection\label{strong:proof:072:end}\end{proof}

Approximate transport controls finite violations. It does not turn a nonzero anomaly into an exact continuum algebra without a limit theorem and a common operator domain.

\section{Quantum refinement towers}\label{strong:section:013}\label{strong:sec:towers}

\begin{definition}[Finite quantum refinement tower]\label{strong:definition:010}\label{strong:def:quantum_tower}
A finite quantum refinement tower consists of Hilbert spaces \(\Hcal_\ell\), isometries
\[
R_\ell:\Hcal_\ell\longrightarrow\Hcal_{\ell+1},
\]
self adjoint Hamiltonians \(H_\ell\), packet projectors \(P_{\ell,a}\), constraints \(C_{\ell,r}\), and positive transfer operators \(T_\ell\). Exact compatibility means that every declared object intertwines with \(R_\ell\). Approximate compatibility is recorded by the corresponding operator norm defects.
\end{definition}

\begin{figure}[t]
\centering
\includegraphics[width=.82\textwidth]{14_continuum_tower.pdf}
\caption{A refinement tower carries Hamiltonians, packets, constraints, and transfer data. Each family has its own compatibility defect.}
\label{strong:fig:continuum_tower}
\end{figure}

\subsection{Exact bounded direct limits}\label{strong:subsec:exact_direct_limit}

Let \(\Hcal_\infty\) be the Hilbert direct limit of the isometric system, and let
\[
\iota_\ell:\Hcal_\ell\longrightarrow\Hcal_\infty
\]
be the canonical isometries.

\begin{theorem}[Bounded Hamiltonian direct limit]\label{strong:theorem:045}\label{strong:thm:bounded_direct_limit}
Assume
\[
H_{\ell+1}R_\ell=R_\ell H_\ell
\]
and
\[
\sup_\ell\norm{H_\ell}<\infty.
\]
Then there is a unique bounded self adjoint operator \(H_\infty\) on \(\Hcal_\infty\) such that
\[
H_\infty\iota_\ell=\iota_\ell H_\ell
\]
for every \(\ell\). It satisfies
\[
\e^{-\ii tH_\infty}\iota_\ell
=
\iota_\ell\e^{-\ii tH_\ell}
\]
and
\[
\e^{-tH_\infty}\iota_\ell
=
\iota_\ell\e^{-tH_\ell}
\]
when the Hamiltonians are nonnegative.
\end{theorem}

\begin{proof}\phantomsection\label{strong:proof:073}
Define \(H_\infty\) on the algebraic union of the embedded levels by
\[
H_\infty\iota_\ell x=\iota_\ell H_\ell x.
\]
Compatibility makes this well defined. The uniform norm bound makes it bounded on the dense union, so it extends uniquely. It is symmetric because every \(H_\ell\) is self adjoint. A bounded symmetric operator defined on the whole Hilbert space is self adjoint. The functional calculus intertwining follows first for polynomials and then for exponentials by uniform convergence.
\phantomsection\label{strong:proof:073:end}\end{proof}

\begin{theorem}[Exact packet direct limit]\label{strong:theorem:046}\label{strong:thm:packet_direct_limit}
Assume the packet families have one fixed finite address set at every level and each packet address \(a\) is compatible in the sense
\[
P_{\ell+1,a}R_\ell=R_\ell P_{\ell,a}.
\]
Then there is an orthogonal projector \(P_{\infty,a}\) on \(\Hcal_\infty\) satisfying
\[
P_{\infty,a}\iota_\ell=\iota_\ell P_{\ell,a}.
\]
Finite completeness and orthogonality pass to the direct limit.
\end{theorem}

\begin{proof}\phantomsection\label{strong:proof:074}
Use the same construction as in \cref{strong:thm:bounded_direct_limit}. Every projector has norm at most one, so no additional boundedness assumption is needed. The projector identities hold on the dense algebraic union and hence on the completion.
\phantomsection\label{strong:proof:074:end}\end{proof}

\subsection{Norm summable approximation}\label{strong:subsec:norm_summable}

The next theorem uses a common Hilbert space. It is appropriate for bounded effective Hamiltonians, not for an unbounded field Hamiltonian.

\begin{theorem}[Norm summable Hamiltonian tower]\label{strong:theorem:047}\label{strong:thm:norm_summable_tower}
Let \((H_n)\) be bounded self adjoint operators on one Hilbert space and assume
\[
\sum_{n=0}^{\infty}\norm{H_{n+1}-H_n}<\infty.
\]
Then \(H_n\) converges in norm to a bounded self adjoint operator \(H\). For every \(T>0\),
\[
\sup_{\abs t\leq T}
\norm{\e^{-\ii tH_n}-\e^{-\ii tH}}
\leq
T\norm{H_n-H}.
\]
If every \(H_n\geq0\), then
\[
\sup_{0\leq t\leq T}
\norm{\e^{-tH_n}-\e^{-tH}}
\leq
T\norm{H_n-H}.
\]
\end{theorem}

\begin{proof}\phantomsection\label{strong:proof:075}
The summability condition makes \((H_n)\) Cauchy in operator norm. The limit is self adjoint. Duhamel's identity gives
\[
\e^{-\ii tH_n}-\e^{-\ii tH}
=-\ii\int_0^t
\e^{-\ii(t-s)H_n}(H_n-H)\e^{-\ii sH}\dd s.
\]
Both unitary factors have norm one. The Euclidean estimate is identical because the nonnegative semigroups are contractions.
\phantomsection\label{strong:proof:075:end}\end{proof}

\begin{corollary}[Limit of packet leakage]\label{strong:corollary:016}\label{strong:cor:limit_leakage}
Let \(P_n\) be orthogonal projectors and \(H_n\) bounded self adjoint operators on one Hilbert space, with \(P_n\to P\) and \(H_n\to H\) in norm. Then for every bounded time interval,
\[
(I-P_n)\e^{-\ii tH_n}P_n
\longrightarrow
(I-P)\e^{-\ii tH}P
\]
uniformly in time.
\end{corollary}

\begin{proof}\phantomsection\label{strong:proof:076}
Add and subtract terms so that one factor changes at a time. The projector terms converge uniformly because \(P_n\to P\) in norm. Duhamel's identity gives
\[
\sup_{\abs t\leq T}
\norm{\e^{-\ii tH_n}-\e^{-\ii tH}}
\leq
T\norm{H_n-H},
\]
which proves uniform convergence of the middle factor.
\phantomsection\label{strong:proof:076:end}\end{proof}

\begin{theorem}[Reflection positivity is closed under limits]\label{strong:theorem:048}\label{strong:thm:reflection_limit}
Suppose \(C_n(j,k)\) are reflection positive kernels on \(\NN\), meaning every finite matrix
\[
\bigl(C_n(j_r,j_s)\bigr)_{r,s}
\]
is positive semidefinite. If
\[
C_n(j,k)\longrightarrow C(j,k)
\]
for every pair, then \(C\) is reflection positive.
\end{theorem}

\begin{proof}\phantomsection\label{strong:proof:077}
For a fixed finite set of times and coefficients, the corresponding quadratic forms are nonnegative and converge to the quadratic form of \(C\). The limit is nonnegative.
\phantomsection\label{strong:proof:077:end}\end{proof}

\begin{theorem}[Constraint closure under norm convergence]\label{strong:theorem:049}\label{strong:thm:constraint_limit}
Let \(C_{n,a}\to C_a\) in norm for every \(a\), and fix structure coefficients \(f_{ab}^{\phantom{ab}c}\). If
\[
\norm{
[C_{n,a},C_{n,b}]
-\ii\sum_cf_{ab}^{\phantom{ab}c}C_{n,c}
}
\longrightarrow0,
\]
then
\[
[C_a,C_b]
=\ii\sum_cf_{ab}^{\phantom{ab}c}C_c.
\]
\end{theorem}

\begin{proof}\phantomsection\label{strong:proof:078}
Operator multiplication and the commutator are continuous in operator norm. Pass to the limit in the anomaly identity.
\phantomsection\label{strong:proof:078:end}\end{proof}

\begin{theorem}[Heat trace convergence from trace norm]\label{strong:theorem:050}\label{strong:thm:trace_norm_heat}
Let \(L_n,L\) be nonnegative self adjoint operators on one Hilbert space. For a fixed \(t>0\), assume their heat operators are trace class and that
\[
\e^{-tL_n}\longrightarrow\e^{-tL}
\]
in trace norm. Then
\[
\Tr(\e^{-tL_n})\longrightarrow\Tr(\e^{-tL}).
\]
If also the heat moment operators are trace class and
\[
L_n\e^{-tL_n}\longrightarrow L\e^{-tL}
\]
in trace norm and the limiting heat trace is nonzero, then the spectral dimensions converge at that time.
\end{theorem}

\begin{proof}\phantomsection\label{strong:proof:079}
The trace is continuous in trace norm. Apply this to the heat operators and to the first heat moments, then take the quotient in \cref{strong:def:spectral_dimension}.
\phantomsection\label{strong:proof:079:end}\end{proof}

\begin{remark}\label{strong:remark:003}\label{strong:rem:unbounded_limit}
The theorems above produce bounded continuum Hamiltonians. A physically useful continuum Hamiltonian is often unbounded. Its construction generally requires strong resolvent convergence, convergence of closed quadratic forms, or another domain sensitive method \cite{Kato1995,ReedSimon1972}. None of those conclusions follows from finite matrix convergence without additional hypotheses.
\end{remark}

\section{Continuum milestones and exact boundary}\label{strong:section:014}\label{strong:sec:boundary}

The finite theorems establish exact identities and bounded limit criteria. They do not collapse the steps between a finite transfer model and a continuum theory. This section records the separate mathematical tasks.

\begin{definition}[Anomaly free continuum constraint system]\label{strong:definition:011}\label{strong:def:anomaly_free_continuum}
A continuum constraint system consists of a Hilbert space, a common invariant dense domain \(\Dcal\), and operators \(C_a\) on \(\Dcal\). It is anomaly free for prescribed structure coefficients when
\[
[C_a,C_b]\psi
=
\ii\sum_cf_{ab}^{\phantom{ab}c}C_c\psi
\]
for every \(\psi\in\Dcal\), and when the associated gauge evolution is defined on the physical quotient or physical subspace.
\end{definition}

Norm closure of bounded finite constraints, as in \cref{strong:thm:constraint_limit}, is one special case. It does not treat unbounded constraints, changing domains, structure functions, or renormalized operator products.

\begin{definition}[Reflection positive Euclidean limit]\label{strong:definition:012}\label{strong:def:reflection_positive_limit}
A reflection positive Euclidean limit requires a limiting algebra of observables, Euclidean time translations, a reflection operation, and correlation functions whose reflected quadratic forms are nonnegative. A field theoretic reconstruction also requires the regularity and covariance assumptions needed to construct a Hilbert space, a vacuum, and a self adjoint time translation generator.
\end{definition}

\Cref{strong:thm:reflection_limit} shows that positivity of each fixed finite Gram matrix survives pointwise convergence. It does not prove existence, locality, clustering, or the remaining Osterwalder Schrader axioms.

\begin{definition}[Einstein infrared criterion]\label{strong:definition:013}\label{strong:def:einstein_criterion}
A quantum geometric tower meets an Einstein infrared criterion only if a controlled long wavelength effective action exists and has an expansion whose leading geometric terms take the form
\[
\Gamma_{\rm eff}[g]
=
\int
\sqrt{\abs g}
\left(
\Lambda
+
\frac{1}{16\pi G}R
+
\text{controlled higher order terms}
\right)
\dd^dx.
\]
The expansion must respect the limiting gauge identities and must include an estimate for the omitted terms on the declared infrared regime.
\end{definition}

A spectral dimension close to four does not imply this criterion. A heat trace records eigenvalue growth. It does not determine the tensor structure, gauge identities, sign, or interaction coefficients of an effective action.

\begin{definition}[Massless spin two criterion]\label{strong:definition:014}\label{strong:def:spin_two_criterion}
A positive massless spin two sector requires a physical two point function or resolvent with a massless pole in the transverse traceless sector, positive residue on physical states, absence of negative norm modes, and compatibility with the limiting gauge constraints. Interactions must preserve the constraint structure.
\end{definition}

A low energy eigenvalue, a spectral dimension plateau, or a finite rank two tensor packet does not establish a massless spin two particle.

\begin{figure}[t]
\centering
\includegraphics[width=.96\textwidth]{16_continuum_boundary.pdf}
\caption{The finite identities occupy the first step. Every later step needs a separate convergence or reconstruction theorem.}
\label{strong:fig:continuum_boundary}
\end{figure}

\begin{table}[t]
\centering
\small
\begin{tabular}{>{\raggedright\arraybackslash}p{.25\textwidth}>{\raggedright\arraybackslash}p{.31\textwidth}>{\raggedright\arraybackslash}p{.34\textwidth}}
\toprule
Milestone & Finite result available here & Missing continuum input \\
\midrule
Symmetry and constraints & exact packets and quotient obstruction & common unbounded domain and renormalized closure \\
Effective dynamics & exact Feshbach map and uniform gap bounds & cutoff independent spectral separation and locality \\
Euclidean transfer & positive finite transfer and finite reflection positivity & convergent Euclidean fields and reconstruction axioms \\
Spectral dimension & lattice heat trace convergence & relation to a physical metric and interacting observables \\
Lorentzian theory & bounded unitary direct limits & unbounded generator, causality, locality, and field algebra \\
Einstein regime & a definition of the required expansion & derivation and error control for the effective action \\
Spin two sector & a precise spectral criterion & positive physical pole and interaction consistency \\
\bottomrule
\end{tabular}
\caption{What the finite results establish and what remains outside the article.}
\label{strong:tab:continuum_milestones}
\end{table}

\section{Conclusion}\label{strong:section:015}\label{strong:sec:conclusion}

A complete packet family gives one coordinate language for finite quantum structure. Group characters produce symmetry packets. Charged operators obey exact selection rules. Circuits become coherent sums over packet paths. The canonical moving connection separates geometric transport from physical mixing. Density matrices and channels use two sided packets, while channel covariance becomes block preservation on operator space and symmetry of the Choi matrix.

Compression to a code gives a natural error correction obstruction. Its vanishing is the Knill Laflamme condition. Diagonalizing the resulting coefficient matrix produces orthogonal syndrome packets and an exact recovery on the error span. Schur Feshbach elimination gives an equally exact description of hidden Hamiltonian packets. Their frequency response is the self energy, and their time response is a positive definite memory kernel.

Positive Euclidean transfer gives finite reflection positivity. Periodic lattice heat traces converge to the continuum torus at every fixed positive diffusion time, and the continuum small time spectral dimension is recovered after the refinement limit. Constraint closure is measured by a quotient of operator space, not by compression to the physical kernel.

The bounded tower theorems show how exact compatibility and norm summability produce limiting quantum evolutions. They also show the boundary of the argument. Unbounded field Hamiltonians, continuum gauge constraints, reflection positive field measures, the Einstein infrared action, and a massless spin two sector require additional theorems. The finite packet calculus organizes those questions but does not settle them.



\startreadingpart{Mathematical proof completions and retained derivations}
\begin{readerroadmap}
The six supplements retain the substantive algebra and geometry of the formal completion. The final supplement constructs local range-bundle charts, their overlaps, and a fiber-valued connection. Three earlier derivations follow because they offer a distinct route through the same finite mathematics.
\end{readerroadmap}
\input{parts/03_supplement_01.tex}
\input{parts/03_supplement_02.tex}
\section{Actual orbits and sharp seed counts}\label{supp3:section:001}
Let $C\in M_{d\times m}(\C)$ be the coefficient matrix of one seed. Define its orbit by actual matrix multiplication:
\[
\mathcal O(C)=\{AC:A\in M_d(\C)\}.
\]
This set is already a linear subspace. The source proves its equality with the linear span of the same orbit, rather than assuming an abstract orbit-dimension formula.

A matrix $X$ belongs to $\mathcal O(C)$ exactly when each of its rows belongs to $\operatorname{Row}(C)$. One direction follows by expanding $AC$. For the converse, choose the coefficients producing each row separately and place them in the corresponding row of $A$. An explicit linear equivalence therefore gives
\[
\mathcal O(C)\simeq \operatorname{Row}(C)^d,
\qquad
\boxed{\dim_\C\mathcal O(C)=d\,\rank C.}
\]
If $d>0$, equality with the full $d\times m$ coefficient space is equivalent to $\rank C=m$. Consequently a cyclic seed requires $m\le d$. The nonempty-row assumption is essential for canceling the factor $d$.

\subsection{Several seeds: both necessity and construction}
For seeds $C_1,\ldots,C_k$, generation means the concrete statement
\[
\forall X\in M_{d\times m}(\C),\quad
\exists A_1,\ldots,A_k\in M_d(\C):\quad
X=\sum_{a=1}^k A_aC_a.
\]
Stack the $C_a$ vertically into a $kd\times m$ matrix $D$. The source proves the exact stacked-action identity and the equivalence
\[
\text{the seeds generate}\quad\Longleftrightarrow\quad\rank D=m.
\]
Thus $m\le kd$ is necessary. It is sufficient as well: embed the $m$ coordinate columns into $kd$ coordinate positions to obtain a full-column-rank matrix $D$, then split its rows back into $k$ seeds. This is an existence proof of actual generating matrices, not only an inequality on possible dimensions.

\subsection{Simultaneous direct-sum generation}
For the stipulated block representation
\[
\mathcal A=\bigoplus_i\bigl(M_{d_i}(\C)\otimes I_{m_i}\bigr),\qquad d_i>0,
\]
the proof establishes that $k$ seeds generate all blocks simultaneously exactly when $m_i\le kd_i$ for every $i$. It follows that
\[
\boxed{\srank_{\mathcal A}(\mathcal H)=\max_i\left\lceil\frac{m_i}{d_i}\right\rceil.}
\]
The final Lean proposition is an \code{IsLeast} statement over the set of feasible seed counts. It proves both attainment and minimality. For an empty block family the finite supremum is defined as $0$, matching the zero module.





\section{Schmidt rank and the reduced state}\label{supp3:section:002}
Use standard product-basis coordinates $\psi_{ij}$ and let $C_{ij}=\psi_{ij}$. An explicit linear equivalence identifies the product-index state vector with its coefficient matrix. The implementation checks the actual Kronecker action:
\[
\operatorname{coeff}\bigl((A\otimes I_B)\psi\bigr)=AC.
\]
It also proves that the transported matrix orbit equals the span of these Kronecker-transformed state vectors. The state-space dimension formula is therefore connected to the physical local action, not only to a separate matrix model.

\subsection{An independent definition of Schmidt rank}
The source defines Schmidt rank as the least number of product-vector summands. A decomposition of length $s$ is the concrete equality
\[
C_{ij}=\sum_{a=1}^{s}u_a(i)v_a(j).
\]
It factors $C=UV$ through an $s$-dimensional intermediate space, so $\rank C\le s$. Conversely, choose a basis of the row space of $C$. Expanding each row in this basis constructs a decomposition with exactly $\rank C$ product terms. Thus
\[
\boxed{r(\psi)=\rank C.}
\]
The definition uses a least natural number, and the equality with matrix rank is proved. The proof does not define Schmidt rank to be matrix rank. It does not construct an orthonormal Schmidt decomposition or prove a new singular-value decomposition theorem.

The matrix-orbit dimension and the coefficient equivalence now yield
\[
\boxed{\dim_\C\spanof\{(A\otimes I_B)\psi:A\in M_{d_A}(\C)\}
=d_A r(\psi).}
\]

\subsection{The actual partial trace and its support}
The reduced matrix is defined by tracing the addressed index of the pure-state density matrix:
\[
(\rho_B)_{j\ell}=\sum_i\psi_{ij}\overline{\psi_{i\ell}}.
\]
For the coefficient convention used here, the exact Gram identity is
\[
\rho_B=C^{\mathsf T}(C^{\mathsf T})^*,
\qquad \rank\rho_B=\rank C=r(\psi).
\]
The transpose is important. With this convention, writing $C^*C$ instead would conjugate the reduced matrix. The Lean source checks the entrywise identity and positivity of the resulting matrix.

Full support is expressed as surjectivity of the reduced matrix's linear map. Assuming $d_A>0$, the checked equivalence is
\[
\boxed{\mathcal O_A(\psi)=\mathcal H_A\otimes\mathcal H_B
\quad\Longleftrightarrow\quad \Ran\rho_B=\mathcal H_B.}
\]
Normalization of $\psi$ is unnecessary for these dimension and support statements. The zero state has rank $0$ under the minimum-product-length definition.





\section{Constructed cyclic character projectors}\label{supp3:section:003}
Let $N>0$ and let $S$ be a unitary complex matrix with $S^N=I$. The implementation uses the genuine standard additive character of $\mathbb Z/N\mathbb Z$ and proves its exponential-root form:
\[
\omega=\exp(2\pi i/N),\qquad
\chi(r)=\omega^{\operatorname{val}(r)}.
\]
Here $\operatorname{val}(r)$ is the canonical representative in $\{0,\ldots,N-1\}$. The representation is constructed as
\[
\rho(r)=S^{\operatorname{val}(r)}.
\]
Reduction of exponents modulo $N$ is proved from $S^N=I$, giving $\rho(r+s)=\rho(r)\rho(s)$ and $\rho(0)=I$. No representation or character-orthogonality assumption is inserted in the final matrix theorem.

\subsection{Exact coefficient formula}
The packet definition is a finite sum over residue classes:
\[
P_j=\frac1N\sum_{r\in\mathbb Z/N\mathbb Z}\chi(-rj)\rho(r).
\]
The source proves
\[
\chi(-rj)=\bigl(\omega^{\operatorname{val}(r)\operatorname{val}(j)}\bigr)^{-1},
\]
so this is precisely the manuscript's $N$-term Fourier projector formula, with canonical representatives.

\subsection{Algebraic identities}
Primitivity of the standard character yields the checked finite sum
\[
\sum_r\chi(tr)=\begin{cases}N,&t=0,\\0,&t\ne0.\end{cases}
\]
Reindexing finite sums then gives
\[
\boxed{P_jP_k=\delta_{jk}P_j,\qquad \sum_jP_j=I,\qquad SP_j=\omega^jP_j.}
\]
Multiplicativity and completeness are first proved for a general representation in a complex algebra, then instantiated with the actual powers of $S$. This organization isolates reusable finite-sum algebra without weakening the final hypotheses.

\subsection{Orthogonal, not merely idempotent}
Unitarity supplies $\rho(r)^*=\rho(-r)$, while the character obeys $\overline{\chi(r)}=\chi(-r)$. Reindexing $r\mapsto-r$ proves $P_j^*=P_j$. The final concrete matrix theorem includes Hermiticity together with multiplication, completeness, and eigenaction. Thus the result establishes orthogonal spectral packets, not only algebraic idempotents.

The proof includes the degenerate modulus $N=1$: finite order then forces $S=I$ and the sole projector is $I$. It does not need an artificial assumption $N>1$.





\section{The actual Euclidean transfer exponential}\label{supp3:section:004}
Let $H$ be a Hermitian matrix on a nonempty finite complex state space. The source defines
\[
T_e(t)=\exp\bigl(-t(H-eI)\bigr)
\]
using Lean's power-series exponential. It does not introduce an unrelated transfer family or assume that a spectral expansion equals this exponential.

\subsection{Connecting power series and spectral calculus}
The shifted generator is proved Hermitian. The implementation then identifies the actual matrix exponential with continuous functional calculus:
\[
\boxed{T_e(t)=f_t(H),\qquad f_t(x)=e^{-t(x-e)}.}
\]
The formal proof uses composition, scalar multiplication, subtraction, and the established agreement between the functional-calculus exponential and the power-series exponential. Spectral mapping is consequently available for the operator that was actually defined:
\[
\Spec_\R T_e(t)=\{e^{-t(x-e)}:x\in\Spec_\R H\}.
\]
Hermiticity also identifies the full complex spectrum with the image of these real eigenvalues. A separate checked theorem states the complex-spectrum interval containment explicitly.

\subsection{A genuine minimum eigenvalue}
The ground energy is the minimum of the finite, nonempty eigenvalue set. The proof establishes both
\[
E_0\in\Spec_\R H,\qquad E_0\le x\quad(x\in\Spec_\R H).
\]
The nonempty-state-space assumption is recorded in the theorem because the minimum of an empty spectrum is not defined by the manuscript's formula.

\subsection{Positivity, semigroup composition, and the interval}
Every eigenvalue of $T_e(t)$ is strictly positive, for arbitrary real $e,t$, since the real exponential is positive. For $e=E_0$ and $t\ge0$ it is also at most one. Therefore
\[
\boxed{T_{E_0}(t)>0,\qquad \Spec T_{E_0}(t)\subseteq(0,1].}
\]
The two scalar multiples of $H-eI$ commute, so the actual power-series exponential obeys
\[
\boxed{T_e(t+s)=T_e(t)T_e(s),\qquad T_e(0)=I.}
\]
The semigroup identity is proved for all real $t,s$. Nonnegative time is needed for the upper spectral bound, not for the composition identity or strict positivity.

These are finite-dimensional operator results. No continuum transfer limit, gravitational path integral, or physical spacetime model is established by this theorem. Numerical figures in the original paper were not recomputed in this release.






\input{parts/03_supplement_04.tex}
\input{parts/03_supplement_05.tex}
\input{parts/03_supplement_06.tex}
\input{parts/04_alternative_derivations.tex}
\startreadingpart{Worked visual atlas}
\begingroup
\input{atlas_definitions.tex}
\setstretch{1.015}\fontsize{10.5}{12.6}\selectfont
\setlength{\parskip}{4pt}
\setlength{\abovedisplayskip}{6pt}\setlength{\belowdisplayskip}{6pt}
\vpage{1}{Three objects, three questions}{HOW TO READ THIS APPENDIX}{Traditional edition: Section II.1, Section II.7, Section II.9.}
A quantum calculation becomes easier to follow when its objects remain visible. This appendix follows a state vector, a density matrix, and a retained part of a dynamical system. They may describe the same experiment, but they do not answer the same question.
\vfig{three_objects}
\head{Amplitudes, not just occupied addresses}
For the Bell vector $\psi=(\ket{00}+\ket{11})/\sqrt2$, the two nonzero entries are amplitudes. The density matrix $\rho=\psi\psi^*$ contains additional off-diagonal entries. Those entries matter to measurements that compare the two addresses. Deleting them is a physical channel, not simply rewriting coordinates.
\head{Elimination, not disappearance}
When a Hamiltonian couples a retained state to a hidden state, the retained amplitude can leave and later return. Its exact equation then remembers the hidden coupling. A smaller matrix obtained by deleting rows and columns does not generally reproduce that equation.
\head{A route through the paper}
Pages \pagerefV{3}--\pagerefV{10} develop seeds, symmetry, and coherent circuit paths. Pages \pagerefV{11}--\pagerefV{15} track leakage and moving subspaces. Pages \pagerefV{16}--\pagerefV{20} distinguish measurement, channels, and recovery. Pages \pagerefV{21}--\pagerefV{27} follow hidden dynamics and finite geometric diagnostics.

This guide reads alongside the traditional volume. Its first 28 worked pages preserve the original appendix, with source lines revised to the final edition. Pages \pagerefV{29}--\pagerefV{40} develop additional strengthened results. Exact example checks are not new Lean certificates.
\takeaway{A packet coordinate change preserves information. A measurement, dephasing operation, or truncation may not. Always identify which operation is actually being performed.}

\vpage{2}{The objects on the page}{AN ATLAS OF TYPES AND CONVENTIONS}{Traditional edition: Definition II.2.1, Definition II.7.1.}
The same letter can stand for different kinds of data in different calculations. Keep the following types fixed while reading a formula. In this appendix, vectors are columns, operators compose from right to left, and $A^*$ denotes the conjugate transpose.
\par\vspace{4pt}
\begin{tabularx}{\linewidth}{@{}p{.13\linewidth}YY@{}}
\toprule
Symbol & Object & What it records\\\midrule
$\HH,\psi$ & Hilbert space and state vector & Addresses and coherent amplitudes.\\
$\AA$ & Declared operator algebra & Which linear operations are available.\\
$P_a$ & Orthogonal packet projector & The whole selected subspace, not one basis vector.\\
$\rho$ & Positive matrix of trace one & Populations and coherences.\\
$E_k$ & Positive measurement effect & An outcome probability $\Tr(E_k\rho)$.\\
$K_a$ & Kraus operator & One branch of a channel; generally not a projector.\\
$H$ & Self-adjoint Hamiltonian & Unitary real-time evolution $e^{-\ii tH}$.\\
$C,\nabla$ & Packet transport generator; connection & Motion of subspaces and differentiation within their ranges.\\
$P,Q$ & Retained and hidden projectors & A decomposition with $Q=I-P$.\\
$\Sigma(z)$ & Self-energy & Spectral influence of the hidden sector.\\
$L,T(\tau)$ & Nonnegative operator; Euclidean transfer & Decay factors and finite heat-trace diagnostics.\\\bottomrule
\end{tabularx}
\head{Normalization and indexing}
The bit order is the displayed left-to-right tensor order: $X_1=X\otimes I\otimes I$. A normalized pure state gives $\rho=\psi\psi^*$; a general coefficient matrix used to count a seed orbit need not be normalized. Scaling a nonzero seed does not change its generated linear span.

A Hamiltonian and a Hadamard gate are both traditionally denoted $H$. Where both appear, the gate is $H_{\rm Had}$ and the Hamiltonian keeps the letter $H$. Numerical labels on a drawing are coordinates or exact example values, not observations of a physical device.
\takeaway{Check the type before interpreting the picture. An idempotent projector, a density matrix, and a channel matrix have different defining equations.}

\vpage{3}{One seed, eight addresses}{CYCLICITY DEPENDS ON THE GRAMMAR}{Traditional edition: Proposition II.2.5.}
Start with $\psi_0=\ket{000}$ and allow the three single-bit flip operators $X_1,X_2,X_3$. For a binary address $a=(a_1,a_2,a_3)$, put
\[
W_a=X_1^{a_1}X_2^{a_2}X_3^{a_3}.
\qquad W_a\psi_0=\ket{a_1a_2a_3}.
\]
The drawing shows the actual eight states reached by these words. Dashed edges change the first bit between the two displayed squares; the other edges change one of the remaining bits.
\vfig{addresses}
\head{Two different meanings of ``reach''}
The eight vectors are orthonormal, so their \emph{linear span} is all of $\CC^8$. This is exactly the cyclicity statement for the algebra generated by the flips. It is not a statement that a product of flips alone prepares every normalized vector.

For example, no such word sends $\ket{000}$ to $(\ket{000}+\ket{001})/\sqrt2$. A Hadamard gate, a suitable linear combination in the observable algebra, or a larger preparation grammar is required for that operation. Counting linearly independent addressed vectors and synthesizing a physical gate sequence are separate tasks.

The exact check constructs each $8\times8$ word matrix, applies it to the seed, and verifies that the matrix of all eight outputs is the identity. It also tests the displayed superposition against every word output. No search over an unspecified gate set is involved.
\takeaway{A cyclic seed generates a vector space under an algebra. That conclusion does not automatically provide a unitary preparation procedure for every vector in it.}

\vpage{4}{Entanglement changes the local span}{SCHMIDT RANK AS ACCESSIBLE MULTIPLICITY}{Traditional edition: Theorem II.2.2, Corollary II.2.4.}
Write a bipartite vector as a coefficient matrix $C$, with the first subsystem indexing rows. Acting only on that subsystem transforms $C$ to $AC$. Every resulting row remains in the row space of $C$, and every independent choice of such rows can be obtained in this way.
\[
\dim\Span\{(A\otimes I)\psi:A\in M_d(\CC)\}
=d\,\rank C.
\]
\vfig{local_orbits}
The product vector $\ket{00}$ has matrix $C=\diag(1,0)$. Its local orbit span contains arbitrary first-subsystem vectors but only the second-subsystem direction $\ket0$, so its dimension is two. The Bell vector has $C=I/\sqrt2$, rank two, and a four-dimensional local orbit span.
\head{What the formula does not say}
The action here allows arbitrary matrices and then takes a linear span. It is not the orbit under local unitary gates alone. Local unitary operations preserve Schmidt coefficients; they do not turn the Bell vector into an arbitrary product vector as a deterministic pure-state transformation.

For the exact calculation, apply the four matrix units $E_{ab}$ to each seed matrix, flatten the outputs, and compute their column rank. The observed dimensions are two and four, respectively. This directly exhibits the rows counted by the theorem rather than inferring them from a drawing of entangled particles.
\takeaway{Schmidt rank measures the remote multiplicity directions accessible to the local operator span. It is not permission for unrestricted deterministic local state preparation.}

\vpage{5}{Why one seed is not always enough}{AN EXACT MULTIPLICITY LOWER BOUND}{Traditional edition: Theorem II.2.3.}
Consider $M_2(\CC)\otimes I_3$ acting on $\CC^2\otimes\CC^3$. A seed is a $2\times3$ coefficient matrix. Its two rows span at most two directions in the three-dimensional multiplicity factor. Therefore no single seed is cyclic for this algebra.
\[
C_1=\mat{1&0&0\\0&1&0},\qquad
C_2=\mat{0&0&1\\0&0&0}.
\]
These two matrices provide a complete set of row directions. Their individual orbit dimensions are four and two, and together those spans fill the six-dimensional state space.
\vfig{seed_rank}
\head{From the example to the finite formula}
For a block $M_d(\CC)\otimes I_m$, $k$ seed matrices contribute at most $kd$ independent multiplicity directions. Hence $k\ge\lceil m/d\rceil$. Partitioning a multiplicity basis into groups of size at most $d$ attains this bound. For a direct sum of such blocks, the worst block determines the required count:
\[
\operatorname{srank}_{\AA}(\HH)=\max_i\lceil m_i/d_i\rceil.
\]
This is an invariant of the declared representation of the algebra, not of the Hilbert-space dimension alone. If the available algebra were enlarged to the full matrix algebra on $\CC^6$, every nonzero seed would instead be cyclic.

The example checks both the ranks of the separate coefficient matrices and the rank three of their stacked row matrix. The argument explains why the construction meets the lower bound; a successful example alone would not prove minimality.
\takeaway{Missing multiplicity directions cannot be created by reusing the same left action. Additional seeds or a larger operator algebra are genuinely different resources.}

\vpage{6}{Four phases resolve a cyclic shift}{CHARACTER PROJECTORS YOU CAN MULTIPLY}{Traditional edition: Corollary II.3.3.}
On four sites, let $S\ket k=\ket{k+1\bmod4}$ and let $\omega=\ii$. The character projectors are the exact matrices
\[
P_j=\frac14\sum_{r=0}^3\ii^{-jr}S^r,
\qquad v_j(k)=\frac12\ii^{-jk},\qquad P_j=v_jv_j^*.
\]
All entries belong to the Gaussian rationals. Thus no approximation to a root of unity is needed in this model.
\vfig{cyclic_four}
\head{Reading the columns}
Each column $v_j$ is a normalized wave around the four sites. Advancing the address multiplies it by $\ii^j$. The entries show phase differences rather than four different probability distributions: every column has uniform site probability $1/4$.

Multiplying the actual projectors gives $P_jP_k=0$ for $j\ne k$, $P_j^2=P_j$, and $\sum_jP_j=I$. In particular, the complete packet tuple $(P_j\psi)_j$ retains the original vector exactly. Recording only the four packet norms is a different operation and can discard phase information between packets.

The four-site shift will return on \pagerefV{27}. There it diagonalizes the Hamiltonian $L=2I-S-S^*$, and its four character phases explain perfect transfer between opposite sites. The same fixed projectors serve as coordinates, symmetry sectors, and a complete interference ledger.
\takeaway{A character projector retains a complex subspace component. Its norm is only one scalar summary of that component, not a substitute for it.}

\vpage{7}{A coarse packet can still contain coherence}{TENSOR REFINEMENT BY TOTAL PARITY}{Traditional edition: Theorem II.3.7, Corollary II.3.8.}
For three qubits, the eight fine projectors are $P_a=\ket a\bra a$. Grouping addresses by Hamming-weight parity gives
\[
P_{\rm even}=P_{000}+P_{011}+P_{101}+P_{110}
=\frac{I+Z\otimes Z\otimes Z}{2},
\qquad P_{\rm odd}=I-P_{\rm even}.
\]
This is a sum of subspaces, not an averaging of their coordinates. The rank-four even packet retains all complex amplitudes on its four children.
\vfig{parity}
\head{Refinement is not measurement}
The fine packet description and the coarse packet description each reconstruct the full vector when every component is retained. Replacing a state by the mixture $\sum_aP_a\rho P_a$ is instead a dephasing channel. That operation removes cross-child coherence and need not agree with dephasing only between even and odd sectors.

For instance, a superposition of $\ket{000}$ and $\ket{011}$ lies entirely in the even packet. Its internal relative phase survives a parity measurement but is lost by a complete computational-basis measurement if outcomes are forgotten. The two measurements use projectors from the same refinement hierarchy but expose different information.

The exact check verifies all eight addresses, their parity grouping, the rank of both coarse projectors, and equality with the tensor expression above. The general residue version replaces parity by total label modulo $N$ and groups the tensor descendants by their summed residue.
\takeaway{A coarse packet is a whole subspace. Internal coherence survives until an operation actually discards it; changing the address description alone does not do so.}

\vpage{8}{Track the Bell state through the circuit}{INTERMEDIATE PACKETS ARE NOT FINAL PACKETS}{Traditional edition: Theorem II.4.1, Example II.4.6.}
The input is $\ket{00}$. First apply a Hadamard to the first qubit, then a controlled flip:
\[
\ket{00}\ \xmapsto{H_{\rm Had}\otimes I}\
\frac{\ket{00}+\ket{10}}{\sqrt2}
\ \xmapsto{\rm CNOT}\
\frac{\ket{00}+\ket{11}}{\sqrt2}.
\]
The displayed columns keep amplitudes, including their normalization, at every step. The middle vector occupies both parity packets. The final vector occupies only the even packet.
\vfig{bell_circuit}
\head{How the exact path expansion works}
Insert $P_{\rm even}+P_{\rm odd}=I$ between the two gates. Each summand records a particular intermediate packet, while the boundary projectors select its input and output packets. Summing all terms recovers the original $4\times4$ circuit matrix:
\[
\sum_{a,b,c}P_b\,\mathrm{CNOT}\,P_a
(H_{\rm Had}\otimes I)P_c
=\mathrm{CNOT}(H_{\rm Had}\otimes I).
\]
This expansion is an identity even though the intermediate parity is not measured. Physically inserting a measurement would alter the problem by removing interference between its alternatives. A coordinate resolution in a derivation should not be mistaken for an experimental intervention.

In the finite check, all eight parity-indexed matrix terms are constructed and summed. The matrix equality is checked before specializing to the input vector. It therefore verifies the decomposition for every input of this example, not merely agreement on the Bell seed.
\takeaway{A final invariant packet does not imply that each preceding gate preserved it. Follow the amplitudes through the actual intermediate resolutions.}

\vpage{9}{Why path probabilities do not add}{A COMPLETE DESTRUCTIVE-INTERFERENCE LEDGER}{Traditional edition: Theorem II.4.2, Theorem II.4.4.}
Two consecutive Hadamard gates multiply to the identity. Starting from $\ket0$, there are two resolved paths to output $\ket1$. Their signed products are
\[
\bra1H_{\rm Had}\ket0\bra0H_{\rm Had}\ket0=\frac12,
\qquad
\bra1H_{\rm Had}\ket1\bra1H_{\rm Had}\ket0=-\frac12.
\]
Their sum is zero. Squaring each separately gives a different answer.
\vfig{interference}
\head{Where the missing term goes}
For path amplitudes $a,b$, the outcome probability is $|a+b|^2=|a|^2+|b|^2+2\Re(a\bar b)$. Here the last term is $-1/2$, canceling the first two terms. The path picture retains this information only if its edges carry complex amplitudes rather than unsigned reachability or probabilities.

The packetwise unitarity identity is likewise an operator identity, including the off-diagonal cancellations:
\[
\sum_b U_{ba}^*U_{bc}=\delta_{ac}P_a,
\qquad U_{ba}=P_bUP_a.
\]
It cannot be replaced by the assertion that every path individually has unit norm. The exact checks verify all four choices of $a,c$ for the two rank-one packets, as well as the two opposite-output path products.

A measurement at the intermediate stage would remove the cross term. In that modified experiment the two alternatives really would contribute total probability $1/2$ to output one. Thus the false calculation accidentally describes a different experiment.
\takeaway{Insert projectors algebraically to resolve paths, then add amplitudes. Add probabilities only when the physical setup has actually made the alternatives incoherent.}

\vpage{10}{A symmetry reveals blocks, not approximations}{ONE HAMILTONIAN IN TWO ORDERS}{Traditional edition: Theorem II.5.1.}
Take the exact rational Hamiltonian
\[
H=X\otimes X+\frac38(Z\otimes I+I\otimes Z).
\]
It commutes with total parity $Z\otimes Z$. Reordering the computational basis into even addresses $(00,11)$ and odd addresses $(01,10)$ displays the invariant subspaces explicitly.
\vfig{hamiltonian}
\[
H_{\rm even}=\mat{3/4&1\\1&-3/4},\qquad
H_{\rm odd}=\mat{0&1\\1&0}.
\]
The even block squares to $(25/16)I$, so its eigenvalues are $\pm5/4$. The odd block squares to $I$, so its eigenvalues are $\pm1$. Exponentiating either block gives the corresponding part of the exact evolution; the off-packet transfer stays zero for every real time.
\head{What has changed?}
Only the order of basis vectors has changed. No matrix entries, amplitudes, or energy levels were discarded. This is why the decomposition is exact. Keeping just the even block would be a separate restriction of the state space, justified for initial data known to be even and for dynamics that preserve it.

The next page breaks a symmetry in a smaller model. There the cross-packet entries no longer vanish, and their effect must be controlled rather than hidden by reordering. The exact tests here multiply the commutator, verify the displayed permutation, and compute all four eigenvalues symbolically.
\takeaway{Block diagonalization uses an invariant decomposition already present in the operator. It is not evidence that an arbitrary chosen packet is invariant.}

\vpage{11}{Quantify the amplitude that escapes}{COMMUTATORS AND PROBABILITY BOUNDS}{Traditional edition: Theorem II.5.3, Corollary II.5.4.}
For a two-level system, choose $H=\varepsilon X$, $P=\ket0\bra0$, and $Q=I-P$. Direct exponentiation gives
\[
e^{-\ii tH}=\cos(\varepsilon t)I-\ii\sin(\varepsilon t)X,
\quad
\norm{Qe^{-\ii tH}P}=|\sin(\varepsilon t)|,
\quad\norm{[H,P]}=|\varepsilon|.
\]
The general commutator estimate becomes the elementary bound $|\sin u|\le |u|$. Unitarity also bounds the amplitude by one.
\vfig{leakage}
\head{From a norm to a probability}
For this initial state, the escape probability is $\sin^2(\varepsilon t)$, not $|\sin(\varepsilon t)|$. Consequently
\[
\Pr(\text{hidden packet at }t)\le\min\{1,\varepsilon^2t^2\}.
\]
At $\varepsilon t=\pi/6$, the amplitude is $1/2$ and the probability is $1/4$. Near zero, amplitudes are first order while escape probabilities are second order. Mixing those two quantities produces an incorrect interpretation of a leakage plot.

The figure samples closed-form functions for display. The supporting check verifies the exact exponential identity, its unitarity, the squared leakage matrix, and the special-time value. The bound itself is supplied by the paper's theorem, not established by sampling the curve.

At larger times amplitude can return from the hidden packet. Thus the commutator bound is a worst-case guarantee, not a prediction that leakage rises monotonically forever.
\takeaway{The norm estimate controls amplitude. Square it for a normalized-state probability bound, and retain the distinction between a bound and the actual oscillatory dynamics.}

\vpage{12}{A moving frame creates a connection term}{NOT ALL MOTION IS PHYSICAL EVOLUTION}{Traditional edition: Theorem II.6.3.}
Set the laboratory Hamiltonian to zero and keep $\psi(t)=\ket0$ fixed. Describe it in rotating coordinates,
\[
U(t)=\mat{\cos t&-\sin t\\\sin t&\cos t},
\qquad y(t)=U(t)^*\psi(t)=\mat{\cos t\\-\sin t}.
\]
The vector has not moved in the laboratory, but its coordinate entries have changed.
\vfig{moving_frame}
\head{The missing term can be computed exactly}
Let $J=\mat{0&-1\\1&0}$. Since $U^*\dot U=J$, the coordinate Hamiltonian is
\[
H_{\rm frame}=-\ii J=\mat{0&\ii\\-\ii&0},
\qquad \ii\dot y=H_{\rm frame}y.
\]
This matrix is Hermitian. Omitting it would predict a constant $y$, contradicting the displayed coordinates. In this example the norm of the omitted time derivative is exactly one, so the failure is not a small numerical residual.

For a nonzero laboratory Hamiltonian, the same product rule adds the transformed physical term $U^*HU$. Thus $U^*HU-\ii U^*\dot U$ separates the physical generator from the coordinate correction. Both pieces belong to the same equation but have different origins.

The drawing shows actual rotated axes at three angles, not three different laboratory trajectories. The check differentiates the exact matrices and verifies the coordinate Schr\"odinger equation symbolically. A finite-difference test is unnecessary for this model.
\takeaway{When coordinates move, even a stationary physical vector acquires changing components. The connection term accounts for that change; it is not an optional force added to the laboratory system.}

\vpage{13}{Differentiate the subspace itself}{THE CANONICAL PACKET CONNECTION}{Traditional edition: Theorem II.6.1, Theorem II.6.2.}
Continue with $P(t)=U(t)\diag(1,0)U(t)^*$ and $Q(t)=I-P(t)$. Differentiating $P^2=P$ gives
\[
\dot P P+P\dot P=\dot P,
\qquad P\dot P P=0.
\]
The second identity says that the derivative has no component entirely inside the current packet. It moves the subspace relative to its complement.
\vfig{connection}
For this complete pair of packets, define
\[
C=\dot P P+\dot Q Q.
\]
Exact calculation yields the constant matrix $C=J$, with $C^*=-C$ and $[C,P]=\dot P$. At $t=0$, the matrices in the plate let the reader verify these identities by multiplying four small arrays.
\head{From two packets to a complete family}
For differentiable orthogonal projectors $P_a$ summing to the identity, the paper constructs
\[
C=\sum_a\dot P_aP_a=\frac12\sum_a[\dot P_a,P_a].
\]
It transports every packet derivative simultaneously. No basis within a packet is needed for this formula. A basis choice would add internal coordinate data without changing which subspaces are selected.

The completeness of the family matters in deriving skew-Hermiticity and the simultaneous transport identities. Using only one summand of an incomplete list is not automatically the same connection. The example retains both $P$ and $Q$ throughout, even when the picture emphasizes only the first line.
\takeaway{The projector derivative describes motion between subspaces. The canonical connection reconstructs that motion from the complete family without choosing internal packet bases.}

\vpage{14}{A chart failure is not a singular fiber}{THE RANGE BUNDLE IS CONSTRUCTED LOCALLY}{Traditional edition: Theorem II.6.7. Range-atlas details: saved v0.6.0 supplement, pp. 2--4.}
The fiber of a projector field is the actual fixed subspace $E_x=\{v:P(x)v=v\}$. To compare it with a nearby anchor fiber $E_a$, the completion proof constructs
\[
T_a(x)=P(x)P(a)+(I-P(x))(I-P(a)).
\]
It equals the identity at the anchor and intertwines the two projectors. Where it is invertible, its restriction provides local fiber coordinates.
\vfig{range_atlas}
\head{An exact example of the domain restriction}
For the rotating line and anchor $P(0)$,
\[
T_0(t)=\cos t\,U(t),\qquad \det T_0(t)=\cos^2t.
\]
The chart fails at $t=\pi/2$, but $P(\pi/2)=\diag(0,1)$ is a perfectly regular rank-one projector. Choosing a vertical anchor supplies a new chart there. What failed was the comparison with one reference fiber, not the fiber itself.

In general, the open sets where $T_a(x)$ is invertible cover the base. The inverse chart maps and smooth overlap maps are constructed from the actual matrices. Their cocycle identities express consistent coordinate changes, rather than assuming a global basis exists.

The v0.6.0 supplement proves the intrinsic connection and its curvature using actual manifold derivatives. Its scope includes bounded idempotents, not only orthogonal ones. The geometric illustration here is finite-dimensional and orthogonal for transparency; that convenience is not promoted to an extra premise of the recorded theorem.
\takeaway{A smooth bundle is built from compatible local charts. The singularity of one chart determinant can be harmless, provided another anchor covers the same smooth fiber.}

\vpage{15}{Curvature measures a local area response}{AN EXACT BLOCH-PROJECTOR CALCULATION}{Traditional edition: Theorem II.6.7, Corollary II.6.8.}
On the unit sphere, set $n=(\sin\theta\cos\phi,\sin\theta\sin\phi,\cos\theta)$ and
\[
P(\theta,\phi)=\frac12(I+n\cdot\sigma).
\]
This is a smooth rank-one orthogonal projector. The intrinsic curvature is the endomorphism-valued two-form $P(\dd P\wedge\dd P)P$. With the paper's scalar convention, its component is
\[
F_{\theta\phi}=-\ii\Tr\!\left(P[\partial_\theta P,\partial_\phi P]\right)
=\frac12\sin\theta.
\]
\vfig{bloch_curvature}
\head{A local frame gives the same answer}
Away from the south pole, use $u_N=(\cos(\theta/2),e^{\ii\phi}\sin(\theta/2))^T$. Then $P=u_Nu_N^*$ and the scalar one-form is
\[
A_N=-\ii u_N^*\dd u_N=\frac{1-\cos\theta}{2}\,\dd\phi,
\qquad \dd A_N=\frac12\sin\theta\,\dd\theta\wedge\dd\phi.
\]
Thus the curvature flux through the cap $0\le\theta\le\theta_0$ is $\pi(1-\cos\theta_0)$. At $\theta_0=\pi/3$, it is exactly $\pi/2$. The plate's latitude is an illustration of this cap; the formula, not its apparent area on the page, supplies the value.

The vanishing of the coordinate density at a pole does not mean the geometric curvature disappears there. The spherical area form also carries the factor $\sin\theta$. A phase assigned to parallel transport additionally requires an orientation and a connection-sign convention; flux and phase should not be interchanged without those choices.
\takeaway{The projector formula and a local frame calculation agree. Curvature is intrinsic even when no single smooth frame covers the whole sphere.}

\vpage{16}{Which coherence can a measurement see?}{THE ANSWER DEPENDS ON ITS EFFECTS}{Traditional edition: Theorem II.7.2, Theorem II.7.4, Corollary II.7.5.}
For the Bell density matrix $\rho=\ket{\Phi^+}\bra{\Phi^+}$, parity dephasing does nothing: both nonzero amplitudes already lie in the even subspace. Computational-basis dephasing instead removes the entries between addresses $00$ and $11$.
\vfig{dephasing}
Every computational-basis effect $E_a=\ket a\bra a$ has the same expectation on both matrices, namely $1/2$ at $00$ and $11$ and zero elsewhere. These measurements read only the diagonal. They cannot distinguish the pure Bell state from this incoherent mixture.
\head{A measurement that resolves the difference}
Use the valid two-outcome measurement with effects $E=\rho$ and $I-E$. Then
\[
\Tr(E\rho)=1,\qquad
\Tr\!\left(E D_{\rm computational}(\rho)\right)=\frac12.
\]
The effect itself is not diagonal in the fine computational packets. It compares their coherent combination, so the paper's block-diagonal invariance hypothesis does not apply.

This example contains three operations that should not be conflated: representing a state by all of its packet matrix blocks, applying a dephasing channel, and reading one specified measurement. Only the second discards the off-diagonal blocks. The third may or may not notice that loss, depending on its effects.

The exact checks construct all four fine projectors, both parity projectors, both dephased states, and the outcome traces. They reject the stronger statement that \emph{every} measurement is invariant under dephasing.
\takeaway{A measurement's inability to see a coherence does not make that coherence absent. It states a restriction on the observable, not an equality of the underlying states.}

\vpage{17}{A channel acts on four matrix entries}{POPULATIONS AND COHERENCES TRANSFORM DIFFERENTLY}{Traditional edition: Theorem II.7.6, Example II.7.12.}
Choose amplitude damping with $\gamma=9/25$:
\[
K_0=\mat{1&0\\0&4/5},\qquad
K_1=\mat{0&3/5\\0&0},\qquad K_0^*K_0+K_1^*K_1=I.
\]
The channel $\Phi(\rho)=\sum_aK_a\rho K_a^*$ is trace preserving and completely positive. The special parameter keeps every displayed coefficient rational.
\vfig{damping}
The superoperator uses the entry order $(\rho_{00},\rho_{01},\rho_{10},\rho_{11})$. Population leaves the excited state with factor $9/25$, while coherences contract by $4/5$. These are not the same rate or the same coordinate.

For $\rho_+=\frac12\mat{1&1\\1&1}$, the exact output is
\[
\Phi(\rho_+)=\mat{17/25&2/5\\2/5&8/25}.
\]
Its trace is one and determinant is $36/625>0$, so it is a mixed positive state. Keeping only population-transfer arrows would not specify its off-diagonal entries.
\head{The complete-positivity certificate}
The Choi matrix is an explicit sum of two outer products of vectorized Kraus operators. Its eigenvalues are $41/25,9/25,0,0$, all nonnegative. This gives an independent finite certificate in addition to Kraus normalization. The Choi matrix has trace two under the unnormalized maximally entangled convention used by the check.
\takeaway{A packet channel acts on pairs of packet addresses, because density matrices have two indices. Population transfer alone does not reconstruct a quantum channel.}

\vpage{18}{An ancilla tests more than positivity}{WHY COMPLETE POSITIVITY IS THE RELEVANT CONDITION}{Traditional edition: Classical input II.7.7.}
Ordinary transpose $\mathsf T$ sends a positive single-system matrix to another positive matrix and preserves its trace. Those properties are insufficient for a quantum channel on a subsystem of an entangled state. The test must include $\id\otimes\mathsf T$ on the larger system.
\vfig{complete_positivity}
Apply this extended map to the Bell density matrix. The result is
\[
(\id\otimes\mathsf T)(\rho_{\rm Bell})
=\frac12\mat{1&0&0&0\\0&0&1&0\\0&1&0&0\\0&0&0&1}.
\]
Its middle block has one negative eigenvalue. More concretely, the normalized vector $w=(\ket{01}-\ket{10})/\sqrt2$ satisfies
\[
w^*(\id\otimes\mathsf T)(\rho_{\rm Bell})w=-\frac12.
\]
This single exact quadratic-form value certifies failure of positivity on the composite system. No numerical eigenvalue tolerance is needed.
\head{Why the Kraus form avoids this problem}
For every auxiliary dimension and positive $X$, each term $(I\otimes K_a)X(I\otimes K_a)^*$ is positive. Their sum is positive as well. This argument explains the word \emph{completely}; checking one input density matrix in the original dimension is not an equivalent test.

The transpose counterexample is an explanatory finite calculation, not a new addition to the formal theorem count. It makes visible the hypothesis that the manuscript's Kraus result already supplies. The exact check constructs the partial transpose with the stated tensor indexing and evaluates its negative witness directly.
\takeaway{A map can look harmless on an isolated system and fail on an entangled extension. Complete positivity retains compatibility with every auxiliary system.}

\vpage{19}{Find the error without reading the logical state}{FOUR EXACT SYNDROME PACKETS}{Traditional edition: Theorem II.8.2, Theorem II.8.4, Example II.8.8.}
Encode a qubit by $\ket{0_L}=\ket{000}$ and $\ket{1_L}=\ket{111}$. Let $P$ project onto this code and take errors $E_a\in\{I,X_1,X_2,X_3\}$. Their images are four orthogonal rank-two packets $Q_a=E_aPE_a^*$ that sum to the full identity.
\vfig{syndromes}
Use the logical state $\psi=(3\ket{0_L}+4\ket{1_L})/5$ and channel probabilities $w=(1,2,3,4)/10$. The noisy state is $\sum_aw_aE_a\psi\psi^*E_a^*$. A syndrome measurement returns label $a$ with probability $w_a$, independent of the logical amplitudes.
\head{The recovery is an actual channel}
The operators $R_a=E_a^*Q_a$ satisfy $\sum_aR_a^*R_a=I$. Applying the recovery $\mathcal R(X)=\sum_aR_aXR_a^*$ to the noisy state gives exactly $\psi\psi^*$. The off-diagonal logical coherence $12/25$ is recovered as well as the populations $9/25$ and $16/25$.

The Knill--Laflamme compressions are $PE_a^*E_bP=\delta_{ab}P$. Their scalar form says that the error labels do not resolve the encoded vector. The exact checks evaluate all sixteen compressions, all sixteen syndrome products, the four probabilities, the recovery normalization, and the full recovered matrix.

This code corrects the specified single-bit flips. It is not claimed to correct arbitrary errors on three qubits, or even arbitrary single-qubit errors. The next example exhibits an error outside its protected family.
\takeaway{A syndrome identifies an error subspace while leaving a two-dimensional logical subspace intact. Measuring the two logical basis states separately would destroy the information the recovery must preserve.}

\vpage{20}{A harmless-looking syndrome can hide a logical error}{THE SCALAR COMPRESSION TEST MATTERS}{Traditional edition: Theorem II.8.2, Proposition II.8.7.}
Apply $Z_1=Z\otimes I\otimes I$ to the same repetition code. It sends $\ket{000}$ to itself and $\ket{111}$ to its negative. Both states remain inside the original code, so the previous syndrome partition reports no bit flip.
\[
\frac{3\ket{000}+4\ket{111}}5
\ \longmapsto\
\frac{3\ket{000}-4\ket{111}}5.
\]
The populations are unchanged, but the relative phase is not.
\vfig{logical_error}
Let $V$ be the code isometry with columns $\ket{000},\ket{111}$. For the proposed error family $\{I,Z_1\}$, its off-diagonal error compression is
\[
V^*I^*Z_1V=\mat{1&0\\0&-1}.
\]
This matrix is not $cI$ for any scalar $c$. Thus the Knill--Laflamme condition fails for that family. Its trace is zero, but its action on the code is certainly nonzero. A trace test alone would miss the obstruction.
\head{Known correction versus family recovery}
If it is known with certainty that $Z_1$ occurred, applying $Z_1$ again reverses it. The failed criterion concerns a \emph{single recovery valid for the whole declared unknown-error family}, not invertibility of each error considered separately. This distinction prevents a misleading conclusion that no phase error can ever be undone.

The finite calculation checks the actual compressed matrix and rejects its equality to zero. The scalar obstruction is then transparent from the two unequal diagonal entries.
\takeaway{Correction must preserve all logical superpositions for the declared error family. Staying in the code and preserving basis probabilities are not sufficient.}

\vpage{21}{Hidden states return as spectral dependence}{TWO INVERTIBILITY CONDITIONS, NOT ONE}{Traditional edition: Theorem II.9.2, Theorem II.9.3.}
Take $H=\mat{0&1\\1&2}$ and retain its first coordinate. The eliminated block is $H_{QQ}=2$. Whenever $z-2$ and the full $zI-H$ are invertible,
\[
H_{\rm eff}(z)=\frac1{z-2},\qquad
\bigl[(zI-H)^{-1}\bigr]_{00}
=\frac1{z-1/(z-2)}=\frac{z-2}{z^2-2z-1}.
\]
This is the corrected hypothesis set recorded in the formal release.
\vfig{feshbach}
At $z=3$, the exact retained resolvent is $1/2$. Bare truncation uses only $H_{PP}=0$ and gives $1/3$, an error of $1/6$. The omitted coupling has changed the response even though the retained diagonal entry did not change.
\head{A pole can belong to the elimination chart}
At $z=2$, the eliminated inverse does not exist, so the displayed Schur formula cannot be applied there as written. Nevertheless, the full resolvent exists and its retained entry is zero. The rational expression has a continuation to this value. By contrast, $1\pm\sqrt2$ are genuine full-resolvent poles.

The full-invertibility hypothesis is independently necessary. At $z=0$, $H=\diag(0,-1)$ has an invertible eliminated block in $zI-H$, but the full matrix is singular. The original prose omitted this condition; the preserved amendments restore it. This appendix uses the corrected form throughout.
\takeaway{An exact elimination formula has a spectral domain. Check both inverses before evaluating it, and distinguish poles of an intermediate inverse from poles of the full physical resolvent.}

\vpage{22}{The smaller system remembers the discarded one}{AN EXACT QUANTUM MEMORY EQUATION}{Traditional edition: Theorem II.9.8, Theorem II.9.11.}
Choose $H=X$, initial state $(1,0)^T$, and retain the first coordinate. The coupled equations and their solution are
\[
\ii\dot x=y,\qquad \ii\dot y=x,\qquad
x(t)=\cos t,\quad y(t)=-\ii\sin t.
\]
The hidden probability returns periodically. Deleting the hidden row and column would instead leave the zero Hamiltonian and predict $x(t)=1$.
\vfig{memory}
\head{Eliminate exactly, not approximately}
Integrating the second equation with $y(0)=0$ gives $y(t)=-\ii\int_0^tx(s)\,ds$. Substitution produces
\[
\dot x(t)=-\int_0^tx(s)\,ds.
\]
The memory kernel is the constant function one. Differentiating once more recovers $\ddot x=-x$ with $x(0)=1$ and $\dot x(0)=0$, explaining the cosine solution. A nonzero hidden initial state would add its own forcing term; it cannot be silently ignored.

The full propagator is a unitary group, but its retained scalar entry is not a semigroup. At $t=s=\pi/4$,
\[
x(t+s)-x(t)x(s)=0-\frac12=-\frac12.
\]
There is no contradiction: compression of a product need not equal the product of the compressed factors when the hidden subspace carries intermediate amplitude.

The check verifies the differential and integral equations symbolically and the displayed nonzero semigroup defect exactly. The plotted curves are merely a visualization of these closed-form identities.
\takeaway{Exact elimination trades a hidden coordinate for memory and initial-data forcing. Replacing that memory by a constant generator is an additional approximation, not part of the exact theorem.}

\vpage{23}{Positive transfer is not a stochastic matrix}{SPECTRAL POSITIVITY IS BASIS INDEPENDENT}{Traditional edition: Theorem II.10.1, Theorem II.10.3.}
For $H=I-X$, the two eigenvalues are zero and two. Its Euclidean transfer is
\[
T(\tau)=e^{-\tau H}
=\frac12\mat{1+q&1-q\\1-q&1+q},\qquad q=e^{-2\tau}.
\]
For $\tau\ge0$, the eigenvalues $1,q$ lie in $(0,1]$, and multiplying two such matrices replaces $q,r$ by $qr$. This is the semigroup law in spectral coordinates.
\vfig{euclidean}
At $q=1/2$, the displayed $T$ happens also to have nonnegative entries and column sums one. Conjugate by the unitary $Z$. The spectrum and positive-definite quadratic form are preserved, but both off-diagonal entries become $-1/4$. Thus positivity as an operator is not entrywise positivity and does not by itself give a Markov interpretation.
\head{A useful finite positive kernel}
For this nonnegative $H$, vectors $u_a$ and times $t_a\ge0$ satisfy
\[
\sum_{a,b}\langle u_a,e^{-(t_a+t_b)H}u_b\rangle
=\left\|\sum_a e^{-t_aH}u_a\right\|^2\ge0.
\]
The factorization makes positivity transparent. It remains a finite operator identity. A continuum reconstruction theorem requires a compatible field structure, limits, positivity hypotheses, and the other assumptions of the chosen reconstruction framework. None follows just from two positive eigenvalues.

The exact check verifies the matrix semigroup law as a polynomial identity in $q,r$ and verifies the unchanged spectrum after conjugation. It rejects entrywise nonnegativity for the conjugated matrix.
\takeaway{Euclidean positivity refers to quadratic forms and spectral values. A stochastic interpretation or a continuum physical reconstruction needs additional structure.}

\vpage{24}{A dimension curve is a scale diagnostic}{A FOUR-SITE HEAT TRACE IN CLOSED FORM}{Traditional edition: Definition II.11.1, Proposition II.11.2.}
For the four-cycle Laplacian $L=2I-S-S^*$, the eigenvalues are $0,2,2,4$. The normalized heat trace and its spectral-dimension diagnostic are therefore
\[
K(\tau)=\frac{1+2e^{-2\tau}+e^{-4\tau}}4
=\frac{(1+e^{-2\tau})^2}4,
\qquad
 d_s(\tau)=-2\tau\frac{K'(\tau)}{K(\tau)}
=\frac{8\tau}{e^{2\tau}+1}.
\]
No numerical differentiation is needed for this finite example.
\vfig{spectral_dimension}
\head{Both ends are finite-size regimes}
As $\tau\to0$, every eigenmode contributes approximately one and the dimension diagnostic tends to zero. As $\tau\to\infty$, the zero mode dominates, so it again tends to zero. A finite spectrum cannot sustain a nonzero plateau at all scales.

For the Cartesian product of two identical four-cycles, the Laplacian eigenvalues add and the heat traces multiply: $K_2=K^2$. Consequently $d_{s,2}=2d_s$ exactly. The second curve is not fitted to the first; the product identity determines it.

At $\tau=\log2/2$, the first curve equals $4\log2/3$. This noninteger value is a scale-dependent spectral quantity, not a claim that the graph has a fractional topological dimension. Four sites are too few to promise a broad geometric plateau. Larger-grid experiments remain finite diagnostics; convergence requires the analytic hypotheses developed on \pagerefV{37}.
\takeaway{A convincing continuum dimension needs a stable scale window that persists under controlled refinement. One appealing finite curve is not that theorem.}

\vpage{25}{Refine the grid while holding the mode fixed}{AN EXACT ERROR ENCLOSURE, NOT A FITTED SLOPE}{Traditional edition: Theorem II.11.4.}
For the circle symbol, the continuum value is $k^2$ and the grid value is
\[
\lambda_{k,h}=\frac4{h^2}\sin^2\!\left(\frac{kh}{2}\right),\qquad h\ne0.
\]
For a fixed low mode, the paper bounds the error by $0\le k^2-\lambda_{k,h}\le k^4h^2/12$. The spacing condition matters because the displayed expression divides by $h^2$.
\vfig{dispersion}
\head{What is certified in the displayed table?}
Take $k=1$ and $h=1/4,1/8,1/16$. For $x=h/2\in[0,1]$, alternating Taylor sums give rational lower and upper bounds for $\sin x$. Squaring these positive endpoints and multiplying by $4/h^2$ yields exact rational enclosures for each error. The displayed decimals only locate those much narrower rational intervals.

Each interval lies strictly inside $(0,h^2/12)$. Dividing successive interval endpoints proves that both error ratios are between $3.99$ and $4$. The check verifies these inequalities by rational arithmetic, not by trusting floating-point sine values or fitting a line on a logarithmic plot.
\head{Why the low-mode qualification matters}
If instead $k=\pi/h$ moves with the grid, the ratio $\lambda_{k,h}/k^2=4/\pi^2$ stays constant, not one. The bound's factor $k^4$ warns that it is not uniform in arbitrary moving high modes. The displayed rational test uses the analytic symbol; it does not assert an integer periodic grid for every chosen spacing.
\takeaway{Second-order convergence is a quantified statement with a specified mode regime. Keep that regime fixed when reading an error curve or comparing refinements.}

\vpage{26}{A zero compressed anomaly can be automatic}{TEST THE CONSTRAINTS BEFORE THE PROJECTION}{Traditional edition: Theorem II.12.2, Proposition II.12.3.}
Take the self-adjoint matrices
\[
C_1=\diag(0,1,-1),\qquad
C_2=\mat{0&0&0\\0&0&1\\0&1&0},\qquad
P=\diag(1,0,0).
\]
Both constraints annihilate $P\CC^3$, and their common kernel is exactly that one-dimensional space. Nevertheless, their full commutator is nonzero.
\vfig{constraints}
\head{Why the compressed test cannot decide closure}
Since $C_1P=C_2P=0$, each product $C_1C_2P$ and $C_2C_1P$ vanishes. Therefore $P[C_1,C_2]P=0$ without any requirement that the unprojected commutator close in the chosen constraint span. In this example the skew off-diagonal commutator is not a linear combination of the two displayed symmetric constraints.

More generally, an anomaly expression made from their commutators and linear combinations of constraints also kills an exact common-kernel projector on the right. This is an important boundary of the diagnostic, recorded by the formal release. A nonzero compressed residual can reveal inconsistent proposed data, but an automatic zero does not construct an anomaly-free constraint algebra.

A substantive closure analysis must inspect the operator relations, their domains, the actual physical subspace, and compatibility under refinement. For infinite systems, these questions also involve domain and limit control that finite matrix arithmetic cannot replace.

The exact example computes both constraint products, the full commutator, and its compression. It deliberately rejects the inference from compressed zero to unprojected zero.
\takeaway{Before celebrating a vanishing residual, ask whether the way its projector was defined already forced it to vanish. A tautological diagnostic cannot supply the missing existence or closure theorem.}

\vpage{27}{One complete calculation: opposite-site transfer}{SEEDS, CHARACTERS, INTERFERENCE, AND DEPHASING}{Traditional edition: Corollary II.3.3, Theorem II.4.4, Theorem II.7.4.}
Return to the four-site shift $S$ and Laplacian $L=2I-S-S^*$. The same character projectors have energies $(0,2,4,2)$. At time $t=\pi/2$, their evolution phases are $(1,-1,1,-1)$, exactly the character eigenvalues of $S^2$. Hence
\[
e^{-\ii\pi L/2}=S^2,\qquad
\ket0\longmapsto\ket2.
\]
\vfig{perfect_transfer}
At halfway time $\pi/4$, the actual amplitude vector is $(-\ii,1,\ii,1)^T/2$, giving uniform site probabilities. The subsequent phases recombine those four amplitudes into a single opposite-site output. Uniform intermediate probabilities do not imply an incoherent intermediate state.
\head{A channel changes the conclusion predictably}
Before evolving, mix the seed density matrix with its complete momentum dephasing:
\[
\rho_\eta=(1-\eta)\ket0\bra0+\eta\sum_jP_j\ket0\bra0P_j
=(1-\eta)\ket0\bra0+\eta I/4,
\quad0\le\eta\le1.
\]
At transfer time, the opposite-site probability is exactly $1-3\eta/4$. Complete momentum dephasing reduces it to $1/4$. This calculation isolates how discarded coherence changes the interference result.

Finally, the shift grammar generates four independent vectors from $\ket0$, while the Hamiltonian powers generate only three: its repeated energy links two momentum directions. The same seed can therefore be cyclic for the shift algebra and noncyclic for the Hamiltonian algebra, even though perfect transfer occurs at one time.
\takeaway{Perfect transfer is an exact finite interference statement. It neither implies cyclicity for every operator grammar nor survives an unaccounted dephasing operation.}

\vpage{28}{What these pictures establish}{AN EVIDENCE GUIDE FOR THE READER}{Traditional edition: Section II.14. Saved formal scope: v0.6.0 ledger and completion supplement.}
The appendix uses small objects to make the paper's constructions inspectable. Every matrix and finite example is supplied in editable source, with a separate check record. The plates visualize those calculations; their apparent smoothness or symmetry is not itself proof.
\par\vspace{4pt}
\begin{tabularx}{\linewidth}{@{}p{.24\linewidth}YY@{}}
\toprule
Evidence & What it supports & What it does not supply\\\midrule
Exact matrices & The displayed finite identities, witnesses, and counterexamples. & A classification of arbitrary systems.\\
Symbolic derivatives & The stated closed-form rotating and Bloch examples. & A proof inferred from numerical differentiation.\\
Rational enclosures & The three displayed dispersion intervals and ratio inequalities. & A continuum theorem from three sample points.\\
Saved Lean release & Its 19 registered finite items under their recorded hypotheses. & A fresh kernel run for this appendix.\\
Numerical replay & The original finite software examples in the tested environment. & Physical experiments or unrestricted convergence.\\\bottomrule
\end{tabularx}
\head{Which source is authoritative?}
The saved formal checkpoint is v0.6.0: 323 recorded theorem declarations in 25 current modules, including supporting lemmas. These are inherited metadata, not a fresh kernel execution. The final traditional edition supplies the current reading references. The original proof sources and the separate v0.5.1 editorial sibling retain their own identities.

The five research boundaries remain separate: an anomaly-free coherent refinement, uniform effective-operator control, Lorentzian reconstruction, an Einstein infrared limit, and a positive physical massless spin-two sector. A finite projector bundle, positive transfer matrix, or spectral diagnostic does not settle them.
\head{A practical way to use the appendix}
Choose one page, reproduce its matrix calculation, and then read the cited final-edition theorem with its hypotheses. The original finite and symbolic suite is replayed for this edition, with its exact assertions and rejected controls recorded separately. Fresh standard-library checks and their bounds are described on \pagerefV{40}. In particular, retain the two resolvent assumptions on \pagerefV{21}, the channel-family distinction on \pagerefV{20}, and the common-kernel caveat on \pagerefV{26}.
\takeaway{Exact finite structure is valuable because its boundary is explicit. The appendix adds understanding and reproducibility, not unproved continuum claims or a new theorem count.}

\vpage{29}{Nonabelian packets carry multiplicity}{A SIX-DIMENSIONAL REGULAR REPRESENTATION}{Traditional edition: Theorem II.3.1, Theorem II.3.2.}
Let $\HH=\CC[S_3]$ have the six orthonormal basis vectors $\ket g$, with left action $U(h)\ket g=\ket{hg}$. The irreducible characters on the identity, transposition, and three-cycle classes are
\[
\begin{array}{c|rrr|r}
 & e & (12)&(123)&d_\lambda\\\hline
\text{trivial}&1&1&1&1\\
\text{sign}&1&-1&1&1\\
\text{standard}&2&0&-1&2
\end{array}
\qquad
P_\lambda=\frac{d_\lambda}{6}\sum_{g\in S_3}\chi_\lambda(g^{-1})U(g).
\]
All matrices in this example are rational. Multiplying them gives self-adjoint idempotents of ranks $1,1,4$, with pairwise zero products and sum $I_6$.
\xfig{s3}
\head{The standard packet is four-dimensional}
The standard representation occurs twice in the regular representation. Its isotypic component therefore has dimension $2\cdot2=4$: one factor carries the two-dimensional irreducible action and the other records multiplicity. A character projector selects the whole isotypic component. It does not select a particular copy or a particular internal vector.

Representation matrix units refine the irreducible factor, while leaving the multiplicity factor untouched. Their multiplication law is stronger information than a list of packet ranks. A basis chosen inside the rank-four packet is a coordinate choice; it must not be mistaken for a canonical decomposition into two individual copies.
\head{What is executed}
The fresh exact check constructs all six permutation matrices, tests all $36$ multiplication relations, computes the character averages, and checks ranks independently by Gaussian elimination and traces. It rejects the claim that the standard isotypic projector has rank two. This is one specified representation of $S_3$, not a computational classification of all finite group representations.
\takeaway{For a nonabelian symmetry, distinguish irreducible dimension from multiplicity. A central character projector retains both.}

\vpage{30}{Rotate whole packets by a polar factor}{AN EXACT LOCAL CHART AND ITS FAILURE}{Traditional edition: Theorem II.6.5, Theorem II.6.6.}
On $\HH=\CC^2$, put $P_0=\diag(1,0)$, $P_1=I-P_0$, and
\[
R=\mat{3/5&-4/5\\4/5&3/5},\qquad Q_a=RP_aR^*,
\qquad S=\sum_{a=0}^1Q_aP_a=\frac35R.
\]
The exact Gram matrices give
\[
(S-I)^*(S-I)=\frac{16}{25}I,\quad S^*S=\frac9{25}I,
\quad\norm{S-I}=\frac45<1.
\]
Consequently the unitary polar factor is $S(S^*S)^{-1/2}=R$, and $Q_aR=RP_a$ for both complete packets.
\xfig{polar}
\head{The norm condition selects a chart}
If the target rank-one packet is orthogonal to the initial one, the two labels are exchanged and $S=0$. Its polar inverse square root is undefined. This is a failure of this direct local construction, although a unitary carrying one packet system to the other still exists. The strict norm assumption is sufficient, not a global existence criterion for all packet rotations.

Short steps can be composed before reaching a distant endpoint. For example, the rational step $R_1=\mat{4/5&-3/5\\3/5&4/5}$ obeys
\[
R_1^2=\mat{7/25&-24/25\\24/25&7/25}.
\]
The multiplication is exact and each step has a valid local overlap. The convergence theorem has a separate analytic contract: a twice continuously differentiable complete projector path on a compact time interval, sufficiently fine partitions, and a path-dependent constant in an $O(|\pi|)$ operator-norm estimate.
\head{Evidence boundary}
The fresh run checks the Gram identities, polar factor, intertwining, and rational two-step product. It also checks that the orthogonal-endpoint overlap has rank zero. It does not certify the analytic convergence theorem by sampling a path.
\takeaway{A local polar chart can fail while the subspaces remain regular. Use the stated chart hypothesis and the separate convergence theorem when composing small rotations.}

\vpage{31}{Coherence order survives a covariant channel}{CHANNEL, CHOI, AND GENERATOR CHECKS}{Traditional edition: Theorem II.7.8, Theorem II.7.10, Theorem II.7.14.}
Take $S=Z=\diag(1,-1)$ on $\CC^2$. Conjugation $X\mapsto ZXZ$ has order zero on diagonal matrices and order one modulo two on off-diagonal matrices. Its projectors on Hilbert--Schmidt operator space are $\Lcal_0(X)=(X+ZXZ)/2$ and $\Lcal_1(X)=(X-ZXZ)/2$.

Choose amplitude damping with $K_0=\diag(1,4/5)$ and $K_1=(3/5)\ket0\bra1$. The channel is trace preserving, population flows from state one to state zero, and each coherence contracts by $4/5$.
\xfig{choi}
With the unnormalized Choi convention $J=\sum_{a,b}\Phi(\ket a\bra b)\otimes\ket a\bra b$ and basis $(00,01,10,11)$, the exact matrix is
\[
J=\mat{1&0&0&4/5\\0&9/25&0&0\\0&0&0&0\\4/5&0&0&16/25},
\qquad [J,Z\otimes Z]=0.
\]
The $00,11$ block is rank one and positive; the other nonzero eigenvalue is $9/25$. The fresh check also verifies every principal minor is nonnegative.
\head{Continuous-time symmetry has its own hypotheses}
For the Lindblad generator with Hamiltonian $H=\diag(0,2)$ and jump $L=(3/5)\ket0\bra1$, $H$ has charge zero and $L$ has charge one. Each matrix unit stays within its coherence order under the generator. Replacing $H$ by $X$ breaks that conclusion: the generator then sends a diagonal matrix to off-diagonal entries.

All four input matrix units are checked exactly for channel covariance, generator order preservation, and zero generator trace. The finite exponential-series argument in the paper extends generator invariance to its semigroup. No master-equation solver is executed here, and symmetry alone supplies no scalar population closure inside higher-rank packets.
\takeaway{Charge is a property of operator space. Covariance preserves its coherent sectors, while complete positivity is checked separately through Kraus or Choi data.}

\vpage{32}{Null Gram directions are removed before recovery}{A TOTAL CHANNEL FROM SINGULAR ERROR DATA}{Traditional edition: Theorem II.8.3, Theorem II.8.4.}
Let $\HH=\CC^3$, with nonzero code projector $P=\diag(1,1,0)$. On the code use the error family $E_1=(3/5)P$, $E_2=(4/5)P$, and $E_3=0$. Its Knill--Laflamme Gram matrix is
\[
c=\frac1{25}\mat{9&12&0\\12&16&0\\0&0&0},
\qquad PE_a^*E_bP=c_{ab}P.
\]
The errors are redundant rather than three distinct syndromes. A real unitary change of error coordinates gives
\[
W=\mat{3/5&4/5&0\\-4/5&3/5&0\\0&0&1},
\quad WcW^*=\diag(1,0,0),
\quad(F_1,F_2,F_3)=(P,0,0).
\]
\xfig{gram}
\head{Normalize only positive eigenvalues}
The positive branch has $V_1=P$ as an isometry \emph{on the code}, and syndrome support $Q=P$. No factor $\lambda^{-1/2}$ is formed on either zero branch: those branches vanish on the code. Treating them as extra recoverable packets would divide by zero.

The syndrome support is smaller than the ambient space. Define a total recovery with Kraus operators
\[
R_1=P,\qquad R_2=\ket0\bra2,
\qquad R_1^*R_1+R_2^*R_2=I_3.
\]
It acts as the identity on code-supported operators and sends the unused complement population into a chosen code state. The declared error map preserves the trace of code states, so recovery returns every code matrix unit, including both off-diagonal units, exactly. It need not preserve states on the unused complement.
\head{Exceptional cases and scope}
This construction assumes a nonzero code so that a complement replacement state exists. The zero-code case is vacuous and belongs to the separate formal boundary in the saved ledger. The exact run verifies all nine Gram compressions, the diagonalization, both null branches, total channel normalization, and all four code matrix units. It adds no Lean coverage and does not implement a numerical spectral decoder for arbitrary Gram data.
\takeaway{Positive Gram support defines genuine syndrome isometries. Null branches vanish; a separate complement rule makes the recovery a total channel.}

\vpage{33}{An approximate defect has a precise norm meaning}{A TWO-POINT OPTIMAL SCALAR PROBLEM}{Traditional edition: Definition II.8.6, Proposition II.8.7.}
On a two-dimensional code, consider the overlap compression
\[
A=PE_a^*E_bP=\diag(11/10,9/10).
\]
For instance, within the code take $E_a=I$ and $E_b=A$. This declares an error family, not by itself a normalized noise channel. The optimal scalar compression has
\[
\inf_{c\in\CC}\norm{A-cI}=\frac1{10},\qquad c=1\ \text{attains it}.
\]
Indeed, for every complex $c$, the triangle inequality gives
\[
\frac15=\left|\frac{11}{10}-\frac9{10}\right|
\leq\left|\frac{11}{10}-c\right|+\left|\frac9{10}-c\right|.
\]
Thus the larger of the two distances is at least $1/10$, while both distances equal $1/10$ at $c=1$.
\xfig{defect}
\head{What one pair contributes to the family defect}
For unit code vectors $\psi,\phi$, this pair therefore satisfies
\[
\left|\langle E_a\psi,E_b\phi\rangle-\langle\psi,\phi\rangle\right|
\leq\frac1{10}.
\]
The bound is attained on either coordinate vector. For the \emph{whole} declared family $\{I,A\}$, the $A^*A$ compression also enters the maximum: its diagonal entries are $121/100,81/100$, giving pair defect $1/5$. Hence the full family value is $\varepsilon_{\rm KL}=1/5$, not $1/10$.

This distinction prevents a tempting shortcut from one favorable overlap to an approximate correction claim. The quantity measures how close all overlap operators are to scalar action on the code. A recovery fidelity bound requires a declared normalized channel, a metric, and a separate continuity result.
\head{What is checked exactly}
For $\varepsilon\in\{0,1/10,1/4,1\}$ the script checks $(A-I)^*(A-I)=\varepsilon^2I$ for $A=\diag(1+\varepsilon,1-\varepsilon)$ and several rational scalar controls. The displayed triangle argument supplies the all-complex-scalar lower bound; a finite list of candidates cannot do that. The family maximum above is also recorded separately.
\takeaway{An overlap norm has a useful operational interpretation. Evaluate every error pair, then keep that statement separate from an unproved recovery fidelity estimate.}

\vpage{34}{Control the hidden spectral response}{SIGN, STABILITY, AND A CONVERGENT EXPANSION}{Traditional edition: Theorem II.9.4, Proposition II.9.5, Theorem II.9.12.}
Retain one coordinate in $H=\mat{0&1\\1&2}$, so $A=0$, $B=1$, and $D=2$. In its elimination domain, $\Sigma(z)=1/(z-2)$. At $z=3+\ii$,
\[
\Sigma(3+\ii)=\frac{1-\ii}{2},\qquad
\Im\Sigma(z)=-(\Im z)\left|(z-2)^{-1}\right|^2\leq0.
\]
The upper half-plane sign comes from self-adjoint hidden data. It is not an entrywise sign condition on an arbitrary effective matrix.
\xfig{stability}
\head{A stability estimate with a common spectral gap}
Compare $(B,D)=(1,2)$ with $(\widetilde B,\widetilde D)=(9/8,17/8)$ at $z=3$. Both hidden spectral distances are at least $\delta=3/4$, and both couplings are bounded by $g=9/8$. Then
\[
|\Sigma-\widetilde\Sigma|=\left|1-\frac{81}{56}\right|=\frac{25}{56},
\qquad
\frac{2g}{\delta}|B-\widetilde B|
+\frac{g^2}{\delta^2}|D-\widetilde D|=\frac{21}{32}.
\]
The theorem bounds the actual difference by this common-gap expression. The hypotheses are uniform spectral separation and coupling control; they must be proved for a refinement family before using the bound uniformly.
\head{A low-energy series and its actual remainder}
At $z=1/2$, $|z|\,|D^{-1}|=1/4<1$, so
\[
\Sigma(z)=-\sum_{n=0}^{\infty}\frac{z^n}{2^{n+1}}.
\]
After the first four terms, the exact remainder has magnitude $1/384$. The geometric-series argument controls the infinite tail. The fresh check verifies that value, the complex sign identity, and the projected-resolvent and bare-truncation identities at four declared real spectral parameters.

At $z=2$ the hidden inverse fails; at a full eigenvalue the retained Feshbach inverse fails. Neither point belongs to the resolvent identity's domain. Analytic continuation of a scalar expression does not erase its theorem hypotheses.
\takeaway{A uniform effective response requires a uniform gap and coupling bound. Exact small examples show how those hypotheses enter each estimate.}

\vpage{35}{Memory positivity is a Gram identity}{TWO HIDDEN FREQUENCIES AND A LOCAL-CLOSURE TEST}{Traditional edition: Theorem II.9.9, Theorem II.9.11.}
Use retained space $\CC$ and hidden space $\CC^2$, with $A=0$, $B=(3/5,4/5)$, and $D=\diag(1,2)$. The full Hamiltonian is $\mat{0&3/5&4/5\\3/5&1&0\\4/5&0&2}$. Its memory kernel is
\[
K(t)=\frac9{25}e^{-\ii t}+\frac{16}{25}e^{-2\ii t},
\qquad K(0)=1.
\]
Each oscillatory coefficient is positive although $K(t)$ need not be real or pointwise nonnegative. At $t=\pi/2$, it equals $-16/25-9\ii/25$.
\xfig{memory_positive}
\head{Positive definite means a matrix test at arbitrary times}
For retained scalars $v_j$ and real times $t_j$, the structural identity is
\[
\sum_{j,k}\overline v_jK(t_j-t_k)v_k
=\left\|\sum_k e^{\ii t_kD}B^*v_k\right\|^2\geq0.
\]
The exact worked check specializes to $t=(0,\pi/2,\pi)$ and $v=(1,2,-1)$. Its phases lie in $\QQ(\ii)$, so both sides are computed without evaluating trigonometric functions. This one Gram test illustrates the factorization; the written identity explains positivity for all times and all vectors.
\head{Why this does not close to a bare local equation}
With hidden initial state zero, the full system has $\ddot x(0)=-(A^2+BB^*)x(0)=-x(0)$. The bare retained equation $\ii\dot x=Ax$ would instead give zero acceleration. Therefore exact local closure for every retained initial state fails. More generally, equality of these accelerations forces $BB^*=0$ and hence $B=0$.

A nonzero initial hidden vector also supplies the forcing term $Be^{-\ii tD}y(0)$. It is additional data, not part of the convolution kernel. The checks cover the displayed kernel values, exact Gram factorization, and acceleration defect; no time discretization is used as proof of an integral equation.
\takeaway{Memory is positive definite through a Gram factorization, not through pointwise signs. Exact bare local closure requires the coupling to vanish.}

\vpage{36}{A positive transfer determines its Hamiltonian}{STRICT POSITIVITY AND THE GROUND GAP}{Traditional edition: Theorem II.10.2, Proposition II.10.5.}
On $\HH=\CC^3$, take the strictly positive contraction
\[
T=\diag(1,1/2,1/4),\qquad H=-\log T=\diag(0,\log2,2\log2).
\]
At unit Euclidean time, $T=e^{-H}$. The self-adjoint logarithm is unique because the logarithm acts on positive spectral values. At integer time $n\geq0$,
\[
T^n=\diag(1,2^{-n},4^{-n}),\qquad
P_0=\diag(1,0,0),\qquad\norm{T^n-P_0}=2^{-n}.
\]
The gap $\Delta=\log2$ determines the rate $e^{-n\Delta}$ exactly in this example.
\xfig{transfer_ground}
\head{Strict positivity is an essential finite hypothesis}
For $T'=\diag(1,0)$, the operator is positive semidefinite, but there is no finite self-adjoint $H'$ on all of $\CC^2$ with $T'=e^{-H'}$: every finite eigenvalue of $H'$ exponentiates to a strictly positive value. One may restrict to the positive support, but that changes the space on which the logarithm is defined.

For a family with nontrivial excited spectra, a uniform exponential rate needs a common positive lower bound on the gap. Ground multiplicity does not cause a problem: $P_0$ selects the entire minimum-energy eigenspace. For a Hamiltonian whose minimum is $E_0\ne0$, the normalized transfer is $e^{-t(H-E_0I)}$ before taking the limit.
\head{Positivity is retained at several Euclidean times}
For arbitrary vectors $u_j$ and nonnegative integer times $n_j$, the reflected kernel obeys
\[
\sum_{j,k}\langle u_j,T^{n_j+n_k}u_k\rangle
=\left\|\sum_jT^{n_j}u_j\right\|^2.
\]
This is a finite Gram factorization. The fresh run checks the exact integer semigroup and diagonal ground error for $n=0,\ldots,12$, and rejects a full-space logarithm for the singular control. Logarithms and the all-time functional-calculus statement are supplied by the written spectral argument, not by these bounded tests.
\takeaway{A positive definite transfer can be logarithmically reconstructed. Ground convergence needs a declared normalization and gap, while continuum reconstruction still needs its own structure.}

\vpage{37}{The dimension window has an order of limits}{FINITE HEAT DATA AND THE QUANTIFIERS}{Traditional edition: Theorem II.11.5, Corollary II.11.8.}
For the unit circle grid with $N$ sites, the scaled eigenvalues are $\lambda_N(k)=4N^2\sin^2(\pi k/N)$. Put $Z_N(t)=\sum_ke^{-t\lambda_N(k)}$ and $d_N(t)=2t\sum_k\lambda_N(k)e^{-t\lambda_N(k)}/Z_N(t)$.

At the single time $t=\log2/16$, two tiny grids admit exact rational heat weights:
\[
\begin{array}{c|c|c|c|c}
N&\text{spectrum}&Z_N&\sum\lambda e^{-t\lambda}&d_N/\log2\\\hline
2&(0,16)&3/2&8&2/3\\
4&(0,32,32,64)&25/16&20&8/5
\end{array}
\]
These are unnormalized traces. Normalizing by $N$ changes $Z_N$ but leaves the logarithmic derivative at that fixed grid unchanged.
\xfig{heat_limits}
\head{The theorem uses two separate choices}
For each $\varepsilon>0$, first choose a fixed positive time $t_\varepsilon$ so that the continuum torus dimension differs from one by less than $\varepsilon/2$. Then choose $N_\varepsilon$ so that every $N\geq N_\varepsilon$ differs from that continuum value by less than $\varepsilon/2$ at the \emph{same time}. This proves one controlled intermediate time for all sufficiently fine grids.

It does not assert a uniform interval width from that argument, a moving time $t_N$ with no extra estimates, or an all-scale dimension plateau. At fixed finite $N$, $d_N(t)\to0$ as $t\downarrow0$ and as $t\to\infty$. By contrast, first taking $N\to\infty$ at positive $t$ and then taking $t\downarrow0$ yields the continuum dimension.
\head{Where the infinite tail is controlled}
The fixed-time proof bounds centered modes below by $16k^2$ and above by $4\pi^2k^2$. The summable Gaussian dominators control both the heat trace and its first moment. A finite grid table supplies neither those tails nor the quantifiers. The fresh check verifies only the displayed two spectra, rational weights, traces, and dimension coefficients. It rejects confusing an unscaled four-cycle spectrum with the unit-circle scaled spectrum.
\takeaway{The strengthened window result is a convergence statement with a fixed-time choice. The exact table is a worked calculation, not its proof.}

\vpage{38}{Transport constraints without hiding closure}{THE IMAGE OF AN ISOMETRY IS NOT THE WHOLE LEVEL}{Traditional edition: Theorem II.12.2, Theorem II.12.5, Proposition II.12.6.}
On $\CC^3=\CC\ket0\oplus\CC^2$, use $C_a=0\oplus\sigma_a$ for $a=x,y,z$. The common kernel is $\CC\ket0$, and the full algebra closes:
\[
[C_x,C_y]=2\ii C_z,\qquad[C_y,C_z]=2\ii C_x,\qquad[C_z,C_x]=2\ii C_y.
\]
These are uncompressed identities. They supply an actual closure test in the quotient by the constraint span, unlike the automatic compressed zero on the physical kernel.
\xfig{constraint_image}
\head{Exact transport holds on the coarse image}
Embed by the coordinate isometry $R:\CC^3\to\CC^4$ and set $\widetilde C_a=C_a\oplus1$. Then $\widetilde C_aR=RC_a$, so
\[
[\widetilde C_a,\widetilde C_b]R=R[C_a,C_b].
\]
The coarse physical vector maps into the fine common kernel. However, the same prescribed Pauli structure does \emph{not} close on the whole fine level: on the added coordinate the commutator is zero, whereas $2\ii\widetilde C_z$ is $2\ii$. Exact transport constrains the image of $R$ and leaves its orthogonal complement as separate data.
\head{An approximate residual is also explicit}
Perturb only $\widetilde C_x$ by $(1/10)\ket0\bra0$. For the unit coarse physical vector, the fine constraint residual is exactly $1/10$. The intertwining defect has the same norm. The general approximate-transport result controls such residuals and commutator errors using both defect sizes and a common constraint norm bound; it does not replace a nonzero anomaly by an exact algebra.

The fresh checks verify all three coarse commutators, the isometry, each constraint intertwining, all nine commutator transports, and the residual. The fine-coordinate counterexample rejects extending an image statement to the whole space. Unbounded constraints would additionally require a common invariant domain.
\takeaway{Closure is an operator-algebra statement. Refinement intertwining carries that statement on the embedded image; control the new degrees of freedom separately.}

\vpage{39}{Two routes to a bounded Hamiltonian limit}{EXACT EMBEDDINGS AND A SUMMABLE TAIL}{Traditional edition: Theorem II.13.2, Theorem II.13.4, Corollary II.13.5.}
For the exact route, let $\HH_n=\CC^n$ with coordinate inclusions and
\[
H_n=\diag(0,1,1/2,\ldots,1/(n-1)),\qquad H_1=(0).
\]
The old energies are retained at every level, $H_{n+1}R_n=R_nH_n$, and $\norm{H_n}\leq1$. On the direct-limit Hilbert space $\ell^2(\NN_0)$, the diagonal sequence defines a bounded self-adjoint $H_\infty$. Even-address projectors intertwine as well and give a fixed two-packet complete system in the limit.
\xfig{tower_tail}
\head{Compatibility and boundedness are different assumptions}
The equally compatible tower $\diag(0,1,2,\ldots,n-1)$ has growing norms. It does not meet the bounded direct-limit theorem, although its natural diagonal operator has a separately definable unbounded domain. That domain-sensitive construction is not obtained by dropping the uniform norm hypothesis.

For the norm route, work instead on the \emph{single fixed space} $\CC^2$ and set
\[
H_n=\diag(0,1+2^{-n}),\qquad H=\diag(0,1).
\]
Here $\norm{H_{n+1}-H_n}=2^{-(n+1)}$ and $\norm{H_n-H}=2^{-n}$. The summable tail supplies an all-$n$ error bound. For $|t|\leq T$, the paper's Duhamel estimate gives
\[
\norm{e^{-\ii tH_n}-e^{-\ii tH}}\leq T2^{-n}.
\]
The Euclidean counterpart holds on $0\leq t\leq T$ because both Hamiltonians are nonnegative. If packet projectors also converge in norm, leakage operators converge uniformly on these time intervals.
\head{Execution scope}
Exact inclusions and projector compatibility are checked through level eight. The two-dimensional tail identities are checked for $n=0,\ldots,12$. Their infinite conclusion uses the displayed geometric-series argument and the theorem's analytic estimate. A finite prefix cannot certify an unspecified continuation of a tower.
\takeaway{Exact embedded compatibility uses a uniform norm bound. Approximate convergence uses a common space and summable defects. Keep those two contracts explicit.}

\vpage{40}{Choose the topology that preserves the quantity}{REFLECTION GRAMS, TRACES, AND FINAL SCOPE}{Traditional edition: Theorem II.13.6, Theorem II.13.8, Section II.14.}
Let $q_n=1/2+2^{-(n+2)}$ and define the scalar reflected kernel $C_n(j,k)=1+q_n^{j+k}$ on nonnegative integers. Each finite Gram matrix is a sum of two positive rank-one matrices, and pointwise convergence gives $C(j,k)=1+(1/2)^{j+k}$. For times $(0,1,2)$ and coefficients $(1,-2,1)$,
\[
\sum_{j,k}v_jC_n(j,k)v_k=(1-q_n)^4\geq0.
\]
This exact check illustrates the finite quadratic-form argument that passes positivity to the limit.
\xfig{trace_topology}
\head{Operator norm does not preserve trace in growing rank}
On $\ell^2$, let $P_n$ select the first $n$ coordinates and set $A_n=P_n/n$. Then $\norm{A_n}=1/n\to0$ while $\Tr A_n=1$ for every $n$. These are positive trace-class operators, used as a topology control; the limiting zero is not a finite-Hamiltonian heat operator.

For actual two-state transfers $T_n=\diag(1,q_n)=e^{-H_n}$, with $H_n=\diag(0,-\log q_n)$ at unit time, the trace-norm error is $2^{-(n+2)}$ and the heat traces converge to $3/2$. The trace-norm criterion for spectral-dimension convergence additionally assumes convergence of the first heat moments in trace norm and a nonzero limiting heat trace. Heat-trace convergence alone does not determine that quotient.
\head{What is fresh, inherited, and still open}
The extension has fresh bounded exact checks in $\QQ(\ii)$ using Python's standard-library \texttt{Fraction}. The original appendix's 180-check record, including its symbolic derivatives, remains available as inherited evidence. The original SymPy suite is replayed for this edition. Its exact assertions and rejected controls remain separate from the extended examples. All original model and vector sources remain supplied.

The saved v0.6.0 ledger lists 19 finite items; no Lean kernel run or expanded formal coverage is asserted. Its five original continuum/physical targets remain open. The strengthened edition supplies analytic proofs and bounded-limit criteria; neither these examples nor their check counts establish the unresolved physical existence claims.
\takeaway{Use pointwise convergence for finite positivity tests, operator norm for bounded dynamics and constraints, and trace norm for heat traces. Every conclusion keeps its own hypotheses and evidence scope.}

\endgroup

\clearpage\FloatBarrier
\section*{References}
\addcontentsline{toc}{section}{References}
\renewcommand{\refname}{}
\bibliographystyle{amsplain}
\bibliography{references}
\end{document}
